%  LaTeX support: latex@mdpi.com 
%  For support, please attach all files needed for compiling as well as the log file, and specify your operating system, LaTeX version, and LaTeX editor.

%=================================================================
\documentclass[algorithms,article,accept,pdftex,moreauthors]{Definitions/mdpi} 

% MDPI internal commands
\firstpage{1} 
\makeatletter 
\setcounter{page}{\@firstpage} 
\makeatother
\pubvolume{1}
\issuenum{1}
\articlenumber{0}
\pubyear{2023}
\copyrightyear{2023}
\externaleditor{Academic Editor: Firstname Lastname} %MDPI: Please add. %<-- Authors' reply: we don't know name of the Academic Editor. 

\datereceived{1 June 2023} 
\daterevised  {11 June 2023}
\dateaccepted{16 June 2023} 
\datepublished{} 
\hreflink{https://doi.org/} % If needed use \linebreak
%\doinum{}

%=================================================================
% Add packages and commands here. The following packages are loaded in our class file: fontenc, inputenc, calc, indentfirst, fancyhdr, graphicx, epstopdf, lastpage, ifthen, lineno, float, amsmath, setspace, enumitem, mathpazo, booktabs, titlesec, etoolbox, tabto, xcolor, soul, multirow, microtype, tikz, totcount, changepage, attrib, upgreek, cleveref, amsthm, hyphenat, natbib, hyperref, footmisc, url, geometry, newfloat, caption

\graphicspath{{figures/}} % path to folder with figures
\usepackage[labelformat=simple]{subcaption}
\renewcommand\thesubfigure{\alph{subfigure}}
\DeclareCaptionLabelFormat{subcaptionlabel}{\normalfont(\textbf{#2}\normalfont)}
\captionsetup[subfigure]{labelformat=subcaptionlabel}

\usepackage{listings}
\usepackage{xcolor}
\usepackage{gensymb} % for \textdegree
\usepackage{tabularx}
\usepackage[para,online,flushleft]{threeparttable}
\usepackage{longtable,booktabs,threeparttablex}
\usepackage{array}
\usepackage{makecell, multirow}
\usepackage{booktabs}
\usepackage{caption}
\usepackage{adjustbox}
\usepackage{longtable}

%=================================================================
% Full title of the paper (Capitalized)
\Title{Optimized Workflow Framework in Construction Projects to Control the Environmental Properties of Soil}

% MDPI internal command: Title for citation in the left column
\TitleCitation{Optimized Workflow Framework in Construction Projects to Control the Environmental Properties of Soil}

% Author Orchid ID: enter ID or remove command
\newcommand{\orcidauthorA}{0000-0002-0577-9936} %  Per Add \orcidA{} behind the author's name
\newcommand{\orcidauthorB}{0000-0002-5759-1089} % Polina Add \orcidB{} behind the author's name

% Authors, for the paper (add full first names)
\Author{Per %MDPI: Please carefully check the accuracy of names and affiliations. %<-- Authors' reply: yes, we confirm that the names and affiliations are correct.
 Lindh $^{1,2,}$*\orcidA{} and Polina Lemenkova $^{3}$\orcidB{}}

%\longauthorlist{yes}

% MDPI internal command: Authors, for metadata in PDF
\AuthorNames{Per Lindh and Polina Lemenkova}

% MDPI internal command: Authors, for citation in the left column
\AuthorCitation{Lindh, P.; Lemenkova, P.}
% If this is a Chicago style journal: Lastname, Firstname, Firstname Lastname, and Firstname Lastname.

% Affiliations / Addresses (Add [1] after \address if there is only one affiliation.)
\address{%
$^{1}$ \quad Department of Investments Technology and Environment, Swedish Transport Administration, Neptunigatan~52, P.O. Box %MDPI: Please check if this should be P.O. Box. %<-- Authors' reply: yes, P.O. Box..
 366, SE-201-23 Malmö, %MDPI: English writing is recommended, please try to provide.  %<-- Authors' reply: Malmö is the official name of the city, we cannot change it. Also in English Wikipedia: https://en.wikipedia.org/wiki/Malm%C3%B6
 Sweden \\
$^{2}$ \quad Division of Building Materials, Department of Building and Environmental Technology, 
Faculty of Engineering (LTH), %MDPI: Please keep the English version if possible. % <-- Authors' reply: corrected.
 Lund University, P.O. Box 118, SE-221-00 %MDPI: Please keep the postal format consistent from the same country. % <-- Authors' reply: corrected.
 Lund, Sweden\\
$^{3}$ \quad Laboratory of Image Synthesis and Analysis (LISA), École Polytechnique de Bruxelles (Brussels Faculty of Engineering), Université Libre de Bruxelles (ULB). Building L, Campus de Solbosch, ULB--LISA CP165/57, Avenue Franklin Roosevelt 50, Brussels 1000, Belgium; polina.lemenkova@ulb.be %<-- Authors' reply: École Polytechnique de Bruxelles is the official name; we need it for publication record of the faculty.
}

% Contact information of the corresponding author
\corres{Correspondence: per.lindh@byggtek.lth.se; %MDPI: We replaced it with the correspondence's email. Please confirm. 2. Please try to keep only one email. % <-- Authors' reply: corrected to one email.
 Tel.: +(46)-0771-921-921 %MDPI: We moved both tel. number here since Tel. number is not allowed in affiliation. Please keep only the correspondence Tel. number. % <-- Authors' reply: corrected.

}

% Current address and/or shared authorship
%\firstnote{Current address: Affiliation 3.} 

% Abstract (Do not insert blank lines, i.e., \\) 
\abstract{To optimize the workflow of civil engineering construction in a harbour, this paper developed a framework of the contaminant leaching assessment carried out on the stabilized/solidified dredged soil material. The specimens included {the sampled }sediments collected from the in situ %MDPI: Greek/Latin expressions should not be in italics. Please confirm. %<-- Authors' reply: yes, we confirm.
 fieldwork in Arendal and {Kongshavn}. The background levels of the concentration of pollutants were evaluated to {assess the} cumulative surface leaching of substances from samples {over two~months}. The {contamination of} soil was assessed using {a} structured workflow scheme on the following {toxic substances}, heavy metals---As, Pb, Cd, Cr, Hg, Ni, and Zn; organic compounds---PAH-16 and PCB; and organotin compounds---TBT. {The} numerical computation and data analysis {were} applied to {the results of }geochemical testing creating computerised solutions to soil quality evaluation in civil engineering. Data modelling {enabled} the estimation of leaching of {the }contaminants in one year. The estimated leaching of As is 0.9153 mg/m\textsuperscript{2}%MDPI: We removed italics for units. Please confirm. %<-- Authors' reply: yes, we confirm.
, for Ni---2.8178 mg/m\textsuperscript{2}, for total PAH-16 as 0.0507~mg/m\textsuperscript{2}, and for TBT---0.00061 mg/m\textsuperscript{2} per year. The performance of the sediments was examined with regard to permeability through a series of the controlled experiments. The environmental engineering tests were implemented in the Swedish Geotechnical Institute (SGI) in a triplicate mode over 64 days. The results were compared for several sites and showed that the amount of As is slightly higher in {Kongshavn} than for Arendal, while the content of Cd, Cr, and Ni is lower. For TBT, the levels are significantly lower than for those at Arendal. The algorithm of permeability tests evaluated the safety of foundation soil for construction of embankments and structures. The {optimized} assessment methods were applied for monitoring coastal areas through {the }evaluated permeability of soil and estimated leaching rates of heavy metals, PHB, PACs, and TBT in selected test sites in harbours of southern Norway. 
} 

% Keywords
\keyword{optimization; quality control; material science; chemical contamination; workflow; soil stabilisation; binders; solidification; leaching; decontamination} 

% The fields PACS, MSC, and JEL may be left empty or commented out if not applicable
\PACS{81.40.Cd; 81.40.Ef; 62.20.Qp; 83.50.Xa; 45.70.Mg; 92.40.Lg; 81.40.Lm; 62.20.M- %MDPI: Please check if this is correct. %<-- Authors' reply: yes, correct: https://ufn.ru/en/pacs/62.20.M-/
}
\MSC{76Axx; 74Exx; 74Fxx}
\JEL{Q00; Q01; Q24; Q55; Q56}

%%%%%%%%%%%%%%%%%%%%%%%%%%%%%%%%%%%%%%%%%%
\begin{document}

\section{Introduction}

Soil treatment is a challenging civil engineering task with many applications, such as soil stabilization~\cite{ref1,ma16113996,Yietal2014,app13084880}, soil {remediation}~\cite{Benschoten,Chu,Bhandari}, and increasing soil permeability~\cite{Rangeard,Lemenkova202222,Schultz}. {These} tasks aim at improving geotechnical and environmental properties of soil prior {to} construction works~\cite{Mohammedetal,Zayyat,SHEOB2023,Tiwari}. {The stabilization of soil foundations }with{ binders} improves {its }strength properties~\cite{DHARINI2023,ELHARIRI202347,Eissa,ZHANG2023131887,HogentoglerWillis}{, while }soil washing removes {the }pollutants and contaminants from soil through separating {it} by particle size or leaching~\cite{Fristad,Finney,Hanson}. 

{Improving} the environmental properties of {polluted }soil {aims at} to {removing} the contaminants{, which can be performed }using{ diverse} washing solutions~\cite{He,Mizutani2016}. Leaching is then performed by cleaning the soil from organic and inorganic {compounds} that are bind to soil particles. {At the same time, the} granulometric type of soil {affects its behaviour during leaching and stabilization}~\cite{Hoang,Schaefers,Krisdani}. {Thus, soil texture affects }the degree of contamination {due to different permeability and capacity to retain pollutants. As a result, fine-grained silt and clays }(particle size of 0.002 to 0.6 mm){ typical in the cold environment of Norway}~\cite{Bache,Andersen,Zhangetal2020} retain more contaminants {due to to clogging of the pores }compared to the middle-grained {sand (}0.6 to 2 mm) or coarse-grained {gravel with >} 2 mm of particle size~\cite{Binner,Kulkarni}. 

A key ingredient for effective soil processing {in environmental engineering }is the development of the optimised design fit for experiments ~\cite{Rahinah,Ninic,Lemenkova202226,Zhang,Lemenkova202309,Wei}{, since methods} of soil treatment differ depending on {its }type (fine-, middle-, or coarse-grained) and properties (water content, density, and temperature). Existing methods of soil {remediation} used in environmental engineering include {the }evaluation of leaching potential from soil {and immobilization of contaminants}~\cite{Mohammedetal2016,Simonetal}. Diverse contaminants can be tested, such as {toxic }heavy metals, {(As, Pb, Cd, Cr, Hg, Ni, and  Zn)}~\cite{Naghipour,WangTabbaa,LiLi,Xia}, Polycyclic Aromatic Hydrocarbon (PAH)~\cite{Liang}, Polychlorinated biphenyl (PCB)~\cite{Aydilek}, and organotin compounds Tributyltin (TBT)~\cite{KUCHARSKI2022136133,Lemenkova202214,MARCIC20062322,Lemenkova202131}. Depending {on} the approaches, soil leaching may be implemented with or without the renewed {leachate}~\cite{1997579}{, and using either }dynamic or extraction tests~\cite{Lemenkova202208}. 

{A }popular {approach of the }environmental engineering {evaluation of soil} properties{ includes the} common leaching test~\cite{Chittoori,Beauchesne,Lemenkova202233}. {Other} examples use the water leach test~\cite{Schreck,Becker}, toxicity characteristic leaching~\cite{Fuessle}, column leach test~\cite{Dayioglu,Sauer}, or soil layering over the substance drainage system~\cite{Yu}. Existing methods follow the existing standards to make use of {the }well-known approaches in soil testing. Such methods focus on remediation of {the }ecologically degraded and contaminated soil, assuming that {its} properties conform to the {models} in {the }existing workflow~\cite{Alpaslan,ma16093444,su15064950}. As a consequence, none of the resulting approaches are tailored to {a better }understanding {of how the} variation of soil properties, such as soil moisture{, }grain size{, and permeability, affect its behaviour in leaching}. {At} the same time, {soil quality is significantly affected by the }the selection of binder types {used for stabilization }\cite{Paniagua,TONINIDEARAUJO2023129757,Filho}{, }the proportion of water--binder ratio~\cite{Kukko} and {binder content}~\cite{Lametal2018}. 

Modelling this information, as most dimensionality reduction methods do, introduces additional parameters, and therefore, {improves the workflow }of soil detoxification and removal of chemical pollutants. Workflow strategies of soil decontamination {and remediation }have high costs~\cite{Ghavami,Mishra,Lu}. The process of soil decontamination includes a complex {expensive }chain of remediation techniques {to remove pollutants, heavy metals, organic and inorganic toxins, and organotin compounds}~\cite{Agarwal,DU2014141,SARKER2023138861}. Normally, this includes both {in situ} (excavating the samples) and {ex situ} {(}batch-wise treatment of specimens using percolating water {in a laboratory})~\cite{LU2023130021,CHEN2010638,KAVAMURA201061}. Given the {constraints of high cost and time pressure related to }soil treatment in large-scale construction works~\cite{ChenLi,recycling1030328,Ahmed}, it becomes essential to optimise a workflow {at a small scale }using {parametric simulation and} modelling {of soil performance }prior to {industrial applications}~\cite{7733108}.  

{The} method {proposed in this article }is inspired by the optimization techniques of soil treatment~\cite{Lindh2001,geotechnics2010005,RehmanMoghal,Yin,Kennedy}{. The main idea is to }factorise {the} parameters that affect soil properties into sub-classes having linear and non-linear distribution of data over time for various soil samples. {In this way, }the optimization of the procedure improved the workflow {of soil decontamination }in {a }small-scale using several batches of soil samples {modelled }prior to  construction works. The data obtained from the numerical experiments on {evaluating the leaching behaviour of contaminants} were processed by statistical methods{. Using }modelling algorithms{, we demonstrated the} non-linear {performance} of leaching for various  pollutants. {The} suitability of {the presented} approach to soil treatment {is illustrated }in graphical plots and tables where the behaviour of soil leaching {is compared }for various contaminants.{ The }benefits of {the proposed} approach {are compared with the }existing techniques on soil leaching tests{, immobilisation of contaminants}, and permeability {tests}.

%%%%%%%%%%%%%%%%%%%%%%%%%%%%%%%%%%%%%%%%%%
\section{{Materials and Methods}}

%----------- subsection ----------->
\subsection{{Sampling and Transportation}}

The laboratory experiments on soil quality {and surface leaching }were carried out {in} the Department of Building Materials {of} Lunds Tekniska Högskola LTH (Faculty of Engineering), Lund University. The {geotechnical }tests were {implemented in} the Swedish Geotechnical Institute (SGI). {Sampling of dredged sediment material was performed by} the Norconsult AS. The tests were {evaluated in the framework of} the research project run by SGI and financed by Swedish Transport Administration. {Soil} samples {were} collected from the test sites in Arendal municipality and Kongshavn{, Port of Oslo (southern Norway)}.

The stabilizing agents {and soil specimens }were mixed for 5 min to achieve {the }homogeneity of the material. Afterwards, {soil samples} were moulded into the piston sleeves {which} had an inner diameter of 50 mm and a height of 170 mm. The sampling procedure included four batches. A large Batch 1 of samples is subdivided into the three {additional} produced batches of soil {specimens} (Batches 2, 3, and 4). Batch 2 contains {the }activated carbon (powdered charcoal). Batch 3 was pre-treated with {the }activated carbon one month before the stabilization process{. }Batch 4 was stabilized with {the }Belite Calcium Sulfoaluminate (BCSA) cement. The BCSA cement was selected due to its properties, it is an environmentally-friendly type of cement with low level of %MDPI: Please check if italics should be removed for CO2. %<-- Authors' reply: corrected.
CO\textsubscript{2} emissions.

%----------- subsection ----------->
\subsection{{Soil Stabilization}}

The engineering tests {aimed at evaluating the} environmental properties of soil{. The workflow }consisted of {experiments on soil stabilization}, permeability tests, surface leaching, and shake table testing. Due to vibrations during transportation, there is always a certain separation in the material which {leads} to the damage of some of the containers and leakage of water. Therefore, the {clay} masses were then {homogenised} for each sample and placed in tight containers with screw lids, see Figure \ref{fig01}. 

The screw lids had a gasket to prevent leakage and evaporation. The samples were {then }transported to Lund in {these }containers marked with the respective test labels. After the transfer{, the specimens were} homogenised, and the water ratio and density determined. The water ratio is the ratio of the amount of water {with regard }to the mass of {the }solid phase (soil). After {homogenisation with a mixer, }soil {samples were }stabilized using different combinations of binders. The choice of binder was based on {the }previous experience from several similar projects in Sweden. In the mixing tests, a binder combination of 30\% CEM I (SS-EN 197-1)~\cite{SIS2011} and 70\% slag (SS-EN 15167-1)~\cite{SIS2006} was used. Furthermore, {the tests were }carried out with different ratios of weight of water divided by weight of binder, known as the water--binder number (vbt)%Please check intended meaning has been retained
. The choice of {the }vbt was based on previous experiences, including {the works }from Arendal 2 in Gothenburg~\cite{Lindh202139}.
\vspace{2pt}
\begin{figure}[H]
%\begin{adjustwidth}{-\extralength}{0cm}
%\centering
\includegraphics[width=7.5 cm]{Fig_01.jpg}
%\end{adjustwidth}
\caption{Containers {in which} dredged {soil materials were} stored after transportation. Photo {source}: Per Lindh. \label{fig01}}
\end{figure}


%----------- subsection ----------->
\subsection{Permeability Tests}

{The} permeability tests were performed {using the standard SS-EN ISO 17892-11:2019 for geotechnical investigation and testing of soil using permeability tests in a laboratory}~\cite{SIS2019}. {The aim was }to evaluate the flow of water through a soil sample {and, thus, }to estimate {the suitability and safety of }soil {for} construction works {through the  calculation of the permeability coefficient}. The permeability tests were carried out {after the }homogenisation and stabilisation of specimens, {in} the SGI {using the} specimens {fabricated in} LTH{. The} specimens {were used }from the same series of {soil }samples {that were }tested for compressive strength after {the }stabilization. The permeability test was carried out in a cell pressure permeameter using {a }standard permeameter device{, }according to method SGI No. 15. This method is used {in tests performed} both {in }Arendal {(Norway) }and Östrand {(Sweden)}. Additionally, {the }shake table testing {has} been carried out by the Swedish Geotechnical Institute (SGI) in Linköping; however, this kind of test is less suitable for the monolithic materials such as soil dredged masses.

The objective of the tests on permeability {consists} in sustainable stabilization of the dredged soil {material}. The characteristics of the specimen tested for the permeability tests {are }common for all the samples {and include the following technical parameters}, sample diameter (mm)---50{.}0 and test temperature %MDPI: We changed C$^\circ$ to $^\circ$C. Please confirm. %<-- Authors' reply: yes, we confirm.
 ($^\circ$C)---$20\pm 2$. The experiments are recorded in SGI under {the }diary number 1.1-2107-0587. {The }sample preparation/compaction method included {the }trimming {of }the pieces of specimen from the main stabilized {soil }sample. The individual characteristics of tested soil specimens are summarised in Table \ref{tab01}.

\vspace{-2pt}
\begin{table}[H] 
\caption{Parameters of the six soil specimens tested for permeability using permeameter in SGI.\label{tab01}}
\begin{tabularx}{\textwidth}{Cp{3.5cm}CC}
\toprule
\textbf{Sample ID} & \textbf{Sample Water Ratio (\%)} & \textbf{Bulk Density \boldmath{(t/m\textsuperscript{3})}} & \textbf{Dry Density \boldmath{(t/m\textsuperscript{3}})} \\
\midrule
$1\_5$ & 60{.}8 & 1{.}62 & 1{.}01 \\
$4\_1$ & 75{.}3 & 1{.}49 & 0{.}85 \\
$5\_4$ & 107 & 1{.}40 & 0{.}68 \\
$3\_5$ & 63{.}5 & 1{.}62 & 0{.}99 \\
$6\_5$ & 45{.}2 & 1{.}77 & 1{.}22 \\
$8\_5$ & 45{.}5 & 1{.}72 & 1{.}18 \\
\bottomrule
\end{tabularx}
\end{table}

%----------- subsection ----------->
\subsection{Surface Leaching Tests}

The surface leaching tests have been carried out {in} SGI in Linköping{, Sweden, }and the analysis of the leachate is carried out by ALS Scandinavia AB. The experiments were performed according to standard SS-EN 15863 according to the European Standard used for evaluating the leaching behaviour of monolithic soil under dynamic and changing environmental conditions~\cite{SIS2015}, see Figure \ref{fig02}.
The tests entail {the following procedure. A} stabilized test {sample} with a known surface area is lowered into a water bath. The water is {then }changed after 6 {h}, 1 {day}, 2.25{, }4{, }9{, }16{, }32, and 64 {days}, see Tables in Appendices: \ref{tabA01},  \ref{tabA03} and  \ref{tabA05} for surface leaching, and Tables \ref{tabA02},  \ref{tabA04} and  \ref{tabA06} for cumulative surface leaching. %MDPI: There is no Table 20. Please check. %<-- Authors' reply: corrected. Added following text: "Tables in Appendices: \ref{tabA01},  \ref{tabA03} and  \ref{tabA05} for surface leaching, and Tables \ref{tabA02},  \ref{tabA04} and \ref{tabA06} for cumulative surface leaching"
 At each water change, the leachate is analysed {under fixed experimental conditions and the concentrations are evaluated} for heavy metals (As, Pb, Cd, Cr, Hg, Ni, and Zn), PAHs (16 pcs), PCBs (7 pcs), and organotin compounds (TBT). 

The surface leaching was performed as a triplicate test with specimens having designations within the test series. {According to this standard, the release of pollutants is evaluated as a function of time when the contaminants are immobilised from a monolithic soil mass during a period of active contact with a leachate aqueous solution.} The samples were named Bach1, Sample 6, Sample 7, and Sample 8A. Three tested specimens were fabricated using {the }dredged material collected from various test points. 

\vspace{2pt}
\begin{figure}[H]
%\begin{adjustwidth}{-\extralength}{0cm}
%\centering
\includegraphics[width=13.7 cm]{Fig_02.jpg}
%\end{adjustwidth}
\caption{Process %MDPI: Please cite the figure in the text and ensure the first citation of each figure appears in numerical order. %<-- Authors' reply: added citation above in the phrase "The experiments were performed according to standard SS-EN 15863..."
 of surface leaching test on soil specimens. Photo source: Per Lindh. \label{fig02}}
\end{figure}

The chemical analysis of the dredged masses was carried out during sampling of soil on leaching. {The test method is performed in accordance to the standard SS-EN 15863 to evaluate leaching behaviour in monolithic material using regularly shaped test portions of soil samples with a periodical renewal of leachate from 6 h to 64 days. The main idea of this approach consist in exposing the soil material to a water leachate for a certain period of time, then analysing the obtained solution with concentration of contaminants. The requirement of the minimum size of 40 mm in all directions is satisfied in the current experiments where the sleeves for soil samples had the dimensions of 50 mm. }The aim of the chemical analysis was to determine the background levels of the chemical elements (heavy metals, PAHs, PCB, and TBT) and compare them to level from {the }leached {concentrations} from the tested soil {specimens in the end of the experiment. }In addition, the test aimed to identify, which of the sample points within the sample area in the harbour was the most contaminated and, therefore, would be used in surface leaching tests. 

In previous projects, {the ultrapure }Milli-Q water has been used {obtained from the extra purification water systems; }however, it results in {a} biased {outcome} due to the content of extra pure water. Therefore, the later projects, such as {those run in the }Arendal 2 and Östrand, have been tested with the leachate from the respective local area. In this study, the leachate from the Port of Oslo has been used {to adjust local environmental setting. The} sample point $1\_2$ was selected for surface leaching {based} on the materials of the experiments. {According to} this method, the specimen is suspended {using the} nylon ropes so that surface leaching can take place on the entire surface of the soil specimen. 

%%%%%%%%%%%%%%%%%%%%%%%%%%%%%%%%%%%%%%%%%%
\section{Results and Discussion}

{The results of this study are based on the} optimised {workflow }framework developed for {the updated }environmental engineering tests {with adjusted binder proportions. The study resulted in } evaluated {properties} of soil {obtained during} the following workflow{:} 

\begin{itemize}
	\item The soil samples were collected from the Kongshavn Harbour, Port of Oslo; {the} representative soil specimens were excavated from different locations (test sites) and placed in containers {for} further soil treating in {a }batchwise mode.
	\item The soil was homogenised using a surface mixer and stabilised with various combination of binders (cement/slag) {tested experimentally}. The soil samples were cropped and trimmed in leach piles to a defined height of specimen to minimise the workflow; unnecessary parts were subtracted from the samples using piles. 
	\item {The }permeability of soil was {analysed} using a permeameter device according to {the }method SGI No. 15, and compared with optimal parameters for several specimens; the temperature and humidity conditions were maintained in the laboratory.
	\item The {immobilisation} of {}the {environmentally critical }contaminants was performed through leaching{ with evaluated concentrations of } seven types of heavy metals, PAH, PCB, and TNT. The dynamics of leaching speed and value was {recorded and evaluated for 6 h to 64 days} with labelling of soil leaching over {the }defined time intervals. 
	\item The obtained data were recorded for statistical analysis and computer-based modelling. The {dynamics of concentration of the }contaminants in soil {was modelled} on the graphs{. }The behaviour of soil was evaluated using intervals through repetitive soil sampling in laboratory SGI. 
	\item {Statistical correction of the }outliers {caused by biased sampling was performed in the dataset to obtain statistically sound results}. 
	\item The prognosis of {leaching} behaviour was numerically modelled using the developed equations and statistical data analysis. The comparison of the achieved levels of pollution and leached contaminants were compared on the original ({background}) level to {estimate the degree of contamination for environmental risk assessment}.  
\end{itemize}

{The} presented workflow of soil treatment enabled the processing of samples of soil specimens {in soil samples collected from Kongshavn, Port of Oslo, southern Norway. The experiments resulted in evaluated }dynamics of leaching and permeability over time of {64 days. The data were extracted form prognosis modelling for one year based on the }computer-assisted statistical analysis{. The trends in dynamics} of contamination {were evaluated }over time and {compared against }the original (background) levels of contaminants. 

Using the optimised techniques, we tested different experimental conditions, as follows. In separated batches of soil,   different ratios of binders {were used to estimate their effects on soil quality }({BCSA }cement/slag){;} different ratios of water amount to binder amount and the water binder number (vbt) {were evaluated}. {The }permeability and leaching of soil data {were tested }on six cases to achieve a statistical soundness of the results. {The }numerical experiments and statistical data processing {were performed }to derive {the }average and standard deviation on the measured data. {The} experiments have shown the importance of exploiting {the diverse binder combinations for improving }soil properties {through stabilization/solidification and leaching. The }non-linearity of the permeability {and leaching behaviour of soil were evaluated }with regard to binder content used for {soil} stabilization. 

\subsection{Soil Permeability}

The results of the permeability testing demonstrated low permeability level achieved for all tested samples of {the }dredged soil masses. Figure \ref{fig03} shows the permeability of soil {stabilized with }different combinations of {binders (}cement and slag). The best effect was obtained at  70\% {added} slag admixture {in} the total amount of binder. Table \ref{tab02} shows measured permeability for various tested soil specimens. The maximal achieved permeability for sample $1\_5$ is 1.5%MDPI: We changed the commas appearing between the digits into decimal dots. Please confirm this revision. %<-- Authors' reply: yes, we confirm.
 $\times 10^{-9}$ after 122{.}22 h of treatment; for the specimen $3\_5$---$1.0 \times 10^{-9}$; for the specimen $4\_1$---$1.2 \times 10^{-9}$; for the specimen $5\_4$---$4.8 \times 10^{-9}$; for the specimen $6\_5$---$3.0 \times 10^{-9}$, and for the specimen $8\_5$---$1.5 \times 10^{-9}$.

\vspace{-2pt}
\begin{figure}[H]
	\includegraphics[width=14.0 cm]{Fig_03.jpg}
\caption{\hl{Permeability} %MDPI: Please change the terms into scientific notation in the figure. (e.g., "$8 \times 10^{3}$", not "8E3"). 2. Please change the hyphen (-) into minus sign ($-$, "U+2212"). e.g., "-1" should be "$-$1".
 of soil in relation to the {ratio} between cement (c) and slag (s) for dredged {soil.} The best results with the lowest permeability of soil show the combination of binders for 30\% cement and 70\% slag. The abbreviation $H_WL_B$ stands for ``high water--low binder'' level in a mixture.\label{fig03}}
\end{figure} 
\vspace{-8pt}

\begin{table}[H] 
\caption{Development of permeability (p, m/s) of the soil specimen over time (t, hours). The maximal achieved permeability level in soil specimen after the period of treatment is in bold text. %MDPI: There is no red in table. Please check. %<-- Authors' reply: corrected (not "red" but "bold").
\label{tab02}}
%\renewcommand{\arraystretch}{1}
%\begin{adjustbox}{width=1\textwidth}
%\tablesize{\footnotesize}
\begin{adjustwidth}{-\extralength}{0cm}
%\centering %% If there is a figure in wide page, please release command \centering
\begin{tabularx}{\fulllength}{cCcCcCcCcCcC}
\toprule
\multicolumn{2}{c}{\textbf{Sample 1\_5}} & \multicolumn{2}{c}{\textbf{Sample 3\_5}} & \multicolumn{2}{c}{\textbf{Sample 4\_1}} & \multicolumn{2}{c}{\textbf{Sample 5\_4}} & \multicolumn{2}{c}{\textbf{Sample 6\_5}} & \multicolumn{2}{c}{\textbf{Sample 8\_5}}\\
\midrule
\textbf{t} & \textbf{p} & \textbf{t} & \textbf{p}&\textbf{t} & \textbf{p}& \textbf{t} & \textbf{p}& \textbf{t} & \textbf{p} & \textbf{t} & \textbf{p}\\
\midrule
41{.}90 &1{.}7~$\times 10^{-9}$ & 41{.}94 &1{.}1~$\times 10^{-9}$ & 41{.}91 & 1{.}2~$\times 10^{-9}$ & 65{.}91 & 6{.}2~$\times 10^{-10}$ & 41{.}96 & 2{.}3~$\times 10^{-9}$ & 17{.}98 & 1{.}2~$\times 10^{-9}$ \\
49{.}96 &1{.}8~$\times 10^{-9}$ & 49{.}99 &1{.}2~$\times 10^{-9}$ & 49{.}97 & 1{.}1~$\times 10^{-9}$ & 74{.}05 & 6{.}1~$\times 10^{-10}$ & 46{.}29 & 2{.}3~$\times 10^{-9}$ & 26{.}02 & 1{.}2~$\times 10^{-9}$ \\
65{.}88 &1{.}9~$\times 10^{-9}$ & 65{.}93 &1{.}4~$\times 10^{-9}$ & 65{.}89 & 1{.}5~$\times 10^{-9}$ & 91{.}27 & 5{.}0~$\times 10^{-10}$ & 50{.}00 & 2{.}3~$\times 10^{-9}$ & 41{.}96 & 1{.}7~$\times 10^{-9}$ \\
74{.}02 &1{.}8~$\times 10^{-9}$ & 74{.}07 &1{.}4~$\times 10^{-9}$ & 74{.}03 & 1{.}8~$\times 10^{-9}$ & 122{.}25 & 4{.}9~$\times 10^{-10}$ & 65{.}94 & 2{.}5~$\times 10^{-9}$ & 50{.}11 & 1{.}9~$\times 10^{-9}$ \\
91{.}23 &1{.}5~$\times 10^{-9}$ & 91{.}29 &1{.}2~$\times 10^{-9}$ & 91{.}24 & 1{.}6~$\times 10^{-9}$ & -- &-- & 74{.}09 & 2{.}9~$\times 10^{-9}$ & 67{.}32 & 1{.}7~$\times 10^{-9}$ \\
-- &-- & -- & -- & -- & -- & -- &-- & 91{.}31 & 3{.}0~$\times 10^{-9}$ & -- & --\\ 
\midrule
122{.}22 & \boldmath{\textbf{1{.}5~$\times 10^{-9}$}} %MDPI: Please check if the bold is necessary. %<-- Authors' reply: yes, to highlight the maximal values (please see reply to previous comment)
 & 122{.}26 & \boldmath{\textbf{1{.}0~$\times 10^{-9}$}} &122{.}23 & \boldmath{\textbf{1{.}2~$\times 10^{-9}$}} & 144{.}98 & \boldmath{\textbf{4{.}8~$\times 10^{-10}$}} & 122{.}27 & \boldmath{\textbf{3{.}0~$\times 10^{-9}$}} & 98{.}28 & \boldmath{\textbf{1{.}5~$\times 10^{-9}$}}\\
\bottomrule
\end{tabularx}
\end{adjustwidth}
%\end{adjustbox}
\end{table}
\vspace{-4pt}

Figure \ref{fig04} shows the {dynamics of the }permeability of soil over time{ in hours. The specimens used in} the experiment {are} indicated {by their }IDs. Here, the lines {show} the dynamics and {the }overall trend in permeability {level }with respect to the time of treatment of {soil }samples {during} stabilization/solidification procedure. The original values {of the }permeability plotted on the graph show {their} variations with a general slight decrease in {trend} which proves the effective {selected }technique of soil treatment for the increase in~strength.

\begin{figure}[H]
\captionsetup[subfigure]{justification=centering}
\makebox[1.05\textwidth][r]{%
	\begin{subfigure}[b]{.55\textwidth}
		\centering
			\includegraphics[width=0.8\textwidth]{Fig_04a.jpg}
		\caption{Specimen ID: $1\_5$}
	\end{subfigure}%
	\begin{subfigure}[b]{.55\textwidth}
		\centering
		\includegraphics[width=0.8\textwidth]{Fig_04b.jpg}
		\caption{Specimen ID: $4\_1$}
	\end{subfigure}%
}\\
\vfill \vspace{1mm}
\makebox[1.05\linewidth][r]{%
	\begin{subfigure}[b]{.55\textwidth}
	\centering
		\includegraphics[width=0.8\textwidth]{Fig_04c.jpg}
		\caption{Specimen ID: $5\_4$}
	\end{subfigure}%
	\begin{subfigure}[b]{.55\textwidth}
		\centering
			\includegraphics[width=0.8\textwidth]{Fig_04d.jpg}
		\caption{Specimen ID: $3\_5$}
	\end{subfigure}%
}\\
\vfill \vspace{1mm}
\makebox[1.05\linewidth][r]{%	
	\begin{subfigure}[b]{.55\textwidth}
		\centering
		\includegraphics[width=0.8\textwidth]{Fig_04e.jpg}
		\caption{Specimen ID: $6\_5$}
	\end{subfigure}%
	\begin{subfigure}[b]{.55\textwidth}
		\centering
		\includegraphics[width=0.8\textwidth]{Fig_04f.jpg}
		\caption{Specimen ID: $8\_5$}
	\end{subfigure}%
}\caption{\hl{Permeability} %MDPI: Please change the hyphen (-) into minus sign ($-$, "U+2212"). e.g., "-1" should be "$-$1". 2. Please change the terms into scientific notation in the figure. (e.g., "$8 \times 10^{3}$", not "8E3"). 3. Please use periods as decimal sign instead of commas. e.g., "0,1" should be "0.1".
 of soil samples changing over time of experiment (in hours), Specimen IDs: (\textbf{a})~$1\_5$, (\textbf{b}) $4\_1$, (\textbf{c}) $5\_4$, (\textbf{d}) $3\_5$, (\textbf{e}) $6\_5$, and (\textbf{f}) $8\_5$.}\label{fig04}
\end{figure}

\vspace{-8pt}
\subsection{Concentrations of Contaminants}

{The detailed chemical information was obtained as concentrations of pollutants in the leachate of soil samples evaluating the leaching mechanism of the contaminants. The obtained results with full data are summarised in Tables and presented in the Appendix of this study. The Appendix} \ref{AppA1} {contains the results of leaching of contaminants in sampling passes 2317--2324; Appendix} \ref{AppB1} {presents the results of leaching of contaminants in sampling passes 2325--2332, and Appendix }\ref{AppC1} {summarises the results of leaching of contaminants in sampling passes 2333--2340. }

The background (original) levels of the concentration of contaminants were evaluated and presented in Table \ref{tab03}. {The} analysis of the original level of contaminants enabled to perform the evaluation of the decontamination of soil based on the discriminating {the }actual values of pollutants against those of the original (background) level. In such a way, the degree of concentration of contaminants in soil samples was evaluated using {the }analysis of leaching dynamics for each contaminant, respectively. Reading these data obtained from the leaching experiments include the analysis of information on {the }amount and concentration of the resolved pollutants with respect to the overall amount of soil and, thus, to evaluate the obtained level of contamination to prevent ecological risks. The evaluated {substances} included seven (7) types of heavy metals (mg/kg), total PAH-16 (mg/kg) and Tributyltin-Sn (TBT) ($\upmu$g Sn/kg), as shown in Table \ref{tab03}. 

\begin{table}[H] 
\setcellgapes{1pt}
\makegapedcells
\caption{Original level of contaminants: 7 types of heavy metals (mg/kg), total PAH-16 (mg/kg), and TBT ($\upmu$g Sn/kg) collected in various test sites (TS) of Kongshavn.\label{tab03}}
\newcolumntype{C}{>{\centering\arraybackslash}X}
%\begin{tabularx}{\textwidth}{CCCCCCCC}
\begin{tabularx}{\textwidth}{p{1.2cm} p{0.8cm} p{0.8cm} p{0.8cm}p{0.6cm}p{0.8cm}p{0.8cm}p{0.8cm}p{2.0cm}p{1.1cm}}
\toprule
\multirow{2.7}{*}{\textbf{Test Site}} & \multicolumn{7}{c}{\textbf{Heavy Metals}} & \multirow{2.7}{*}{\textbf{Total PAH-16}} & \multirow{2.7}{*}{\textbf{TBT-Sn}}\\ \cmidrule{2-8}
& \textbf{As} & \textbf{Pb} & \textbf{Cd} & \textbf{Cr} & \textbf{Hg} & \textbf{Ni} & \textbf{Zn} & & \\
\midrule
	1--2 	& 13	& 47 & 0{.}5 & 47 & 0{.}365 & 41 & 200 & 9{.}4 & 37 \\
	3--4 & 10	& 34 & 0{.}2 & 43 & 0{.}17 & 38 & 140 & 0{.}5 & 2{.}8 \\
	5 &    9	& 38 & 0{.}32 & 45 & 0{.}174 & 30 & 140 & 0{.}94 & <2.0 \\
	6 &   7{.}5	& 45 & 0{.}1 & 37 & 0{.}317 & 26 & 180 & 1{.}4 & 38 \\
	8 &   13	& 100 & 0{.}86 & 52 & 0{.}574 & 50 & 350 & 8{.}1 & 5{.}9 \\
\bottomrule
\end{tabularx}
\end{table}

{The} most important parameters in the context of soil contamination by heavy metals have been assessed for the following metal types, arsenic (As), lead (Pb), cadmium (Cd), chromium (Cr), mercury (Hg), nickel (Ni), zinc (Zn), as shown in Table \ref{tab03}. The tests also evaluating the concentrations of elements in soil samples were tested for the organic compounds---16~priority Polycyclic Aromatic Hydrocarbon (PAH) and organotin compounds Tributyltin (TBT), Table \ref{tab03}, and for organic chemicals Polychlorinated biphenyl (PCB) (7 pcs), \mbox{Table \ref{tab04}}. {The hydraulic conductivity evaluated the leachate percolation through the soil monolith, and reported in the Appendix in Table} \ref{tabA01} {showing 251--275 $\upmu$S for samplings 2317--2324, Table} \ref{tabA03} {showing 253--274 $\upmu$S for samplings 2325--2332, and in Table} \ref{tabA05} {showing \linebreak 252--276~$\upmu$S for sampling 2333--2340. Thus, in all cases the conductivity slightly increases in course of the experiment from 6 h to 64 days.}

\vspace{-2pt}
\begin{table}[H] 
\caption{Concentration level of origin for PCBs in Kongshavn test sites (TS) (mg/kg TS).\label{tab04}}
\tablesize{\footnotesize}
\begin{tabularx}{\linewidth}{p{0.5cm}p{1.27cm}p{1.27cm}p{1.27cm}p{1.27cm}p{1.17cm}p{1.17cm}p{1.27cm}p{1.3cm}}
\toprule
\textbf{TS} & \textbf{PCB 28} & \textbf{PCB 52} & \textbf{PCB 101} & \textbf{PCB 118} & \textbf{PCB 153} & \textbf{PCB 138} & \textbf{PCB 180} & \textbf{$\sum$ 7 PCB}\\
\midrule
1--2 & 0{.}0054   & 0{.}019     & 0{.}014 	     & 0{.}0098  & 0{.}20     & 0{.}20       & 0{.}013     & 0{.}1 \\
3--4 & <0{.}00050 & 0{.}00061 & <0{.}00050    & <0{.}00050 & 0{.}00063 & 0{.}00063 & <0{.}00050 & 0{.}0019 \\
5 & <0{.}00053 & <0{.}00053  & <0{.}00053  & <0{.}00053 & 0{.}0011   & 0{.}0011     & 0{.}00070     & 0{.}0029 \\
6 & 0{.}0017   & 0{.}0059 	& 0{.}0026       & 0{.}0031  & 0{.}0033  & 0{.}0033      & 0{.}0022      & 0{.}022 \\
8 & 0{.}0063   & 0{.}017 	& 0{.}015         & 0{.}012     & 0{.}012    & 0{.}012      & 0{.}0071      & 0{.}081\\
\bottomrule
\end{tabularx}
\end{table}

The withdrawal of water for analysis of surface leaching was performed with the following intervals of time periods after 6h of testing, 1d (24 h), 2{.}25days (54 h), and then on days 4, 9, 16, 36, and 64. The repetitions of the tests arranged periodically with a non-linear time gaps in the laboratory across the soil samples {enabled} to indicate {the }marks of leaching levels of {each of the }contaminants separating the amounts of the leached substances from the soil samples. The analysis demonstrated the following leaching values (mg/m\textsuperscript{2}): the highest values of As reached 0.88 with average 0.71; the highest values of Ni reached 2.48 with average 2.20; the maximal leaching of PAH-16 is 0.053 with average 0.039; {and }the leaching of TBT increased to 0.0010 with average of 0.0009{, respectively}. 

In order to make a comparison of the contamination level in current state against the pre-industrial conditions, the deviation classification for heavy metals in sediments has been carried out according to the existing Swedish standards~\cite{ZefferSamuelsson}. The table is recalculated using values from tables 34 and 36 in the Swedish Environmental Protection Agency's report 4914 (mg/kg TS). Thus, the identification of the high level marks in contamination enables to extract information on changes in {concentration of the contaminants }during the {assessed} period, while the amplitude of contamination enables to evaluate the {degree and extent} of {the }increase {in pollution }for each types of contaminant{. Here, the evaluation was performed by classes for target substances, such as }heavy metals, TBT, PCB, and PAH, compared to the pre-industrial period with {a lower} pollution level in soil. 

{The obtained results} show that dredged soil {material} collected from the Kongshavn tested areas are generally above Class 1 (Table \ref{tab05}). The classes in deviation ranking of heavy metals are visually breaking the dataset into five clusters of data according to the value of pollution repetitively for each type of the analysed heavy metals, as shown in Table \ref{tab05}. The description of these values of contamination were stored and used for the comparison of the dynamics of the contaminants for diverse periods, days 1, 2.25, 4, 9, 16, 36, and 64. The full data presented in Appendix of this article included the records on the contamination level, the names of the contaminants, and location of the sample sites. 

\vspace{-2pt}
\begin{table}[H] 
%\begin{threeparttable}
%\centering
\tablesize{\small}
\caption{Deviation classification of heavy metals according to the Swedish standards (mg/kg TS).\label{tab05}}
\begin{tabularx}{\textwidth}{CCCCCC}
\toprule
\textbf{Metal} & \textbf{Class 1} & \textbf{Class 2} & \textbf{Class 3} & \textbf{Class 4} & \textbf{Class 5}\\
\midrule
As & <10 & 10--17 & 17--28 & 28--45 & >45\\
Cd & <0{.}2 & 1{.}2--0{.}5 & 0{.}5--1{.}2 & 1{.}2 & >3\\
Cr & <40 & 40--48 & 48--60 & 60--72 & >72\\
Cu & <15 & 15--30 & 30--49{.}5 & 49{.}5--79{.}5 & >79{.}5\\
Hg & <0{.}04 & 0{.}04--0{.}12 & 0{.}12--0{.}4 & 0{.}4--1 & >1\\
Ni & <30 & 30--45 & 46--66 & 66--99 & >99\\
Pb & <25 & 25--40 & 40--65 & 65--110 & >110\\
Zn & <85 & 85--127{.}5 & 127{.}5--204 & 204--357 & >357\\
\bottomrule
\end{tabularx}
%\begin{tablenotes}
     % \small
     \footnotesize{Notations for Classes in Table \ref{tab05}: %MDPI: Please confirm if the italics should be retained. %<-- Authors' reply: italics is deleted (not necessary).
  Class 1---No insignificant deviation from comparative values; Class 2---Small deviation the comparative value; Class 3---Clear deviation from comparative values; Class 4---Large deviation from comparative values; Class 5---Very large deviation from comparative values.}

   % \end{tablenotes}
%  \end{threeparttable}
\end{table}

\vspace{-14pt}
\subsection{Surface Leaching}

{Tables} summarising the analysis of the leachate are presented in the Appendix. {The analysis of the results of the }surface leaching tests {resulted in a} lower values than the expected limits of contaminants, which is reported in the tests protocols see the \mbox{Appendices \ref{AppA1}--\ref{AppC1}} {for full chemical protocols}. In such cases{, }when the levels of pollutants are lower than the reported limits, it can be concluded that the leaching of the {evaluated }substance {has a} low {level, }and the soil is environmentally safe. Table \ref{tab06} shows the distribution of values above or below the reporting limit for metals summarised in the additional Tables of Appendices %MDPI: Please confirm the change. %<-- Authors' reply: yes, we confirm.
 \ref{AppA1}--\ref{AppC1}.

\vspace{-2pt}
\begin{table}[H] 
%\begin{threeparttable}
\caption{Values of concentrations of leaching contaminants (mg/m\textsuperscript{2}) from soil samples: above or below the reported limits.\label{tab06}}
\begin{tabularx}{\textwidth}{CCCCCCCC}
\toprule
\textbf{Sample} & \textbf{As} & \textbf{Pb} & \textbf{Cd} & \textbf{Cr} & \textbf{Hg} & \textbf{Ni} & \textbf{Zn} \\
\midrule
B 1, S-6 & above & below & below & below & below & above & both \\
B 1, S-7 & above & below & below & below & below & above & below \\
B 1, S-8A & above & below & below & below & below & above & both \\
\bottomrule
\end{tabularx}
%\begin{tablenotes}
     \footnotesize{Notations for abbreviations in Table %MDPI: Please check if this should be Table 6. %<-- Authors' reply: yes, Table 6.
 \ref{tab06}:} %MDPI: Please confirm if the italics should be retained. %<-- Authors' reply: italics is deleted (not necessary).
  B1---Batch 1; S6---Sample 6; S7---Sample 7; S-8A---Sample 8A; As---Arsenic; Pb---Lead; Cd---Cadmium; Cr---Chromium; Hg---Mercury; Ni---Nickel; Zn---Zinc. 
 %   \end{tablenotes}
 % \end{threeparttable}
\end{table}

Among the analysed heavy metals, arsenic (As) and nickel (Ni) {demonstrated the levels of concentration} above the detection limit for {specimens treated during }the entire period from 6 h to 64 days of leaching. Figure \ref{fig05} shows the {cumulative} leaching of arsenic {(As) over} time {where} leaching can be considered as diffusion controlled{, as shown in }Figure \ref{fig05}. The lines on the graphs were interpolated as several segments of splines, separated by the Batches and Samples between each measurement and compared for average values, logarithm curve and standard deviation {between the samples within the same time gap}. The lines on the graph show the contamination level{ as }curvatures for each Batch for evaluated specimens. {A }standard deviation {of the leached As on} day 16 was significantly higher, e.g., in sample {No.} 6. It may indicate a biased measurement in the laboratory on a given specific sample. Nevertheless, for the analysis of leaching{, }the data remain statistical unbiased, and {were} included in the calculation of the amount of {the leached As}. The average value of the three samples is used to visualise a statistical trend line. The equation from the trend line is {calculated to estimated} the expected leaching of As after one year. 
\vspace{-2pt}
\begin{table}[H] 
%\begin{threeparttable}
%\centering
\caption{Surface %MDPI: We movedTable 7 behind it's first citation. Please confirm. %<-- Authors' reply: yes, we confirm.
 leaching of contaminants (heavy metals and TBT, mg/m\textsuperscript{2}) from soil after 9 days.\label{tab07}}

\begin{adjustwidth}{-\extralength}{0cm}
%\centering %% If there is a figure in wide page, please release command \centering
\small
%\begin{tabularx*}{\linewidth}{@{\extracolsep{\fill}}p{0.8cm}p{1.0cm}p{1.0cm}p{1.0cm}p{1.6cm}p{1.6cm}p{1.6cm}p{1.6cm}p{1.6cm}}
\setlength\tabcolsep{0pt}
\begin{tabularx}{\fulllength}{@{\extracolsep{\fill}} ccccccccc }
\toprule
\textbf{Metal} & \boldmath{$O\_P6$} & \boldmath{$O\_P7$} & \boldmath{$O\_P8A$} & \textbf{A \boldmath{$H_WL_B$} 23} & \textbf{A \boldmath{$H_WL_B$} 17} & \textbf{A \boldmath{$H_WL_B$} 18} & \textbf{T \boldmath{$H_WL_B$} 11} & \textbf{T \boldmath{$H_WL_B$} 09}\\
\midrule
As &  0{.}0807 & 0{.}0782 & 0{.}105 & 0{.}0276 & 0{.}0050 & <0{.}0014 & 0{.}0042 & 0{.}0090 \\
Cd & <0{.}004 & <0{.}004 & <0{.}004 & 0{.}0025 & 0{.}00035 & <0{.}00028 & <0{.}0003 & <0{.}0003 \\
Cr & <0{.}04 & <0{.}04 & <0{.}04 & 0{.}006 & 0{.}010 & 0{.}0091 & 0{.}0091 & 0{.}024 \\
Hg & <0{.}002 & <0{.}002 & <0{.}002 & <0{.}0063 & <0{.}0070 & <0{.}0070 & <0{.}0076 & <0{.}0082 \\
Ni & 0{.}272 & 0{.}411 & 0{.}4 & 0{.}52 & 1{.}1 & 0{.}14 & 0{.}46 & 0{.}69 \\
Zn & <0{.}2 & <0{.}2 & <0{.}2 & <0{.}013 & <0{.}014 & <0{.}014 & <0{.}015 & <0{.}02 \\
TBT & 101~$\times 10^{-6}$ & 103~$\times 10^{-6}$ & 157~$\times 10^{-6}$ & 275~$\times 10^{-6}$ & 1180~$\times 10^{-6}$ & 315~$\times 10^{-6}$ & -- & -- \\
\bottomrule
\end{tabularx}
\end{adjustwidth}
%\begin{tablenotes}
      \footnotesize{Notations for Table \ref{tab07} %MDPI: there is no Table 8. Please check. %<-- Authors' reply: corrected (Table 7). Italics is also removed.
:  $H_WL_B$ indicate a ``high water / low binder'' ratio in a mixture. The abbreviations indicate the locations where measurements were performed: Oslo (O), Arendal (A), and Timrå (T). }
  %  \end{tablenotes}
%  \end{threeparttable}
\end{table}
\vspace{-12pt}


\begin{figure}[H]
%\begin{adjustwidth}{-\extralength}{0cm}
%\centering
	\includegraphics[width=13.9 cm]{Fig_05.jpg}
%\end{adjustwidth}
\caption{{Cumulative} leaching of arsenic (As) against time in tested specimen samples.\label{fig05}}
\end{figure}

{The} intensive industrial development in the region of Kongshavn (Port of Oslo) {in recent decades resulted in }significant areas of {degraded }lands and {contaminated }soil. Therefore, the need for remediation of soil requires a comparison of data on {the }contamination level both {in a }retrospective {mode }and as {a prognosis. The} estimation of possible future leaching of the{ selected contaminants} has been {performed }for arsenic (As), nickel (Ni), PAH-16, and TBT. {Using} the regression equation from Figure \ref{fig05} and a time period of 365 days, {the} potential leaching amount {was modelled }per m\textsuperscript{2} {for one year, }as shown in Equation (\ref{eq01}). This gives an estimated leaching of arsenic (As) of 0.9153 mg/m\textsuperscript{2} per year.

\begin{equation}\label{eq01}
	y= 0.128 \ln(time) + 0.161
\end{equation}

Figure \ref{fig06} shows a cumulative leaching of nickel (Ni) over time. The average value of the three samples is used to plot a statistical trend line showing general development {and} the interpolated trend in the increase in leaching of Ni. The equation from the trend line is used to calculate the expected leaching of Ni {in} one year. A trend of Ni leaching {is }based on the evaluated data from Batch 1 {using} Samples 6, 7, and 8A{.} For {Ni}, the {results of estimated} leaching look similar and in the same way as for {As}. Likewise, the regression equation was used from Figure \ref{fig07} {for modelling the} data using {a} time span of 365 days to evaluate possible leaching of Ni per m\textsuperscript{2} during one year{ by} Equation (\ref{eq02}). The computed {data} give an estimated amount of leaching of {Ni} as 2.8178 mg/m\textsuperscript{2} per year.

\begin{equation}\label{eq02}
	y= 0.3522 \ln(time) + 0.7399
\end{equation}

Figure \ref{fig07} shows a graph of the cumulative leaching of PAH-16 in soil samples over time. Here, the results vary a little between sample 8A and samples 6 and 7. However, the standard deviation between those samples is low. The average value of the three samples is used to visualise a statistical trend line. The equation from {this} line is then used to calculate the expected leaching of PAH-16 after 365 days in m\textsuperscript{2}, {according to} Equation (\ref{eq03}). This gives an estimated value of possible leaching of total PAH-16 as 0.0507 mg/m\textsuperscript{2} from the soil in {Kongshavn} per {one }calendar year.

\begin{equation}\label{eq03}
	y= 0.0074 \ln(time) + 0.0071
\end{equation}

\begin{figure}[H]
%\begin{adjustwidth}{-\extralength}{0cm}
%\centering
\includegraphics[width=13.9 cm]{Fig_06.jpg}
%\end{adjustwidth}
\caption{Cumulative leaching of nickel (Ni) {over} time in tested soil specimens.\label{fig06}}
\end{figure}

{The} actual content of the PAH-16 substance {reflects} the results of the statistical analysis. The {location of sampling} marks was identified with each sampling case {representing} a central point of the line in a given transect. 
{The }detection limit{ was }determined {with regard to} the standard deviation (Figure \ref{fig07}) obtained when evaluating soil samples. {Leaching of PAH-16 is }determined as 10 times the standard deviation for specimens and is roughly three~times higher than the detected limits of PAH-16 presented in specimens. The dynamics of the {PAH-16 }leaching was detected based on {the }frequent measurements as regular samplings of the soil, represented on the lines of the graphs as dots encountered once for the period of measurements on each line segment. The lowest concentration level of PAH-16 was determined as {a }limit of quantification, that is, quantitatively assessed with {a }satisfactory certainty. 

The overall trend of {PAH-16 }leaching is increasing with regard to the line path, i.e., the total amount of contaminant is gradually released from the specimens along with soil treatment. This proves a positive effect of solidification and stabilization of soil on the removal of {PAH-16 }contaminants from soil. {The leaching of PCB was not calculated since} the values for leached amount of PCBs are {positive, i.e., }lower than the {threshold} limit{.} 

Figure \ref{fig08} {shows the dynamics} of the cumulative leaching of organotin compound tributyltin (TBT) over time. Using the regression equation from Figure \ref{fig08} extrapolated for a time period 365 days, the {TBT }leaching was modelled for a year in m\textsuperscript{2}, {following} Equation~(\ref{eq04}). The obtained results give an estimated amount of leaching of TBT of 0.00061~mg/m\textsuperscript{2} per year. The average value of the three samples is used to draw a statistical trend line {which was used to calculate the} equation {for} the expected leaching of TBT {within} one year. The {threshold} limit specified for TBT is standard for the concentration parameter{. }The variations in the {TBT} concentration can be affected by{ the }dilution due to matrix disturbances, limited sample quantity, or low dry matter content. The {threshold} limit of chemical {substances was estimated following} the {analytical} method {which evaluates} the lowest {level of }concentration at which {the }contaminants are detected in a soil sample. 

\begin{equation}\label{eq04}
	y= 0.0001 \ln(time) + 0.0002
\end{equation}

\begin{figure}[H]
\begin{adjustwidth}{-\extralength}{0cm}
%\centering
	\hspace{130pt}\includegraphics[width=13.9 cm]{Fig_07.jpg}
\end{adjustwidth}
\caption{Cumulative leaching of PAH-16 in tested soil samples.\label{fig07}}
\end{figure}

For evaluation of the obtained results, a comparison has been completed between the results for surface leaching received in the present study against those from other studies~\cite{,Inui,Holleran,Chenetal2021}. {More specifically, }~\cite{Sandhu} {demonstrates similar results on leaching of As, Pb, and Sb}; {the immobilisation of Cr, Cu, and Zn were assessed by }\cite{Linetal2012}. {Furthermore, }\cite{Kimetal2002} {reports the leaching of diverse heavy metals from remediated soil}. {The removal of PHB that is a waste product from the aluminium industry %Please check intended meaning has been retained
	is reported in }\cite{Bongo} {using stabilization and neutralization methods}.

Table \ref{tab07} summarises the values of leaching recorded in samples from Oslo (O), Arendal (A), and Timrå (T). The results are based on the measurements performed after 9 days of soil testing, and show surface leaching of contaminants (heavy metals and TBT, mg/m\textsuperscript{2}) from soil after 9 days of testing, as shown in Table \ref{tab07}. The {threshold} limits of contaminant concentration of TBT and heavy metals with regard to the content of binder (water--binder ratio) were evaluated in soil samples as a collective designation for the highest levels of the substances in samples and compared for various test locations, as shown in Table \ref{tab07}. The analysis of Table \ref{tab07} indicates that the amount of As is slightly higher in Kongshavn Port of Oslo (O) than for Arendal (A) and Timrå (T); however, the content of Cd, Cr, and Ni is lower. {Similar results were reported earlier in the study on the immobilisation of Mg, Cu, Zn, Al, and K which were evaluated in Swedish soils in Timr\aa }~\cite{Nordmark}.


\vspace{-4pt}
As for TBT, the levels are significantly lower in Oslo than for the corresponding sample {in} Arendal. The TBT was not {analysed} in the Timrå project, and, therefore, {could} only be compared between the {sampling results from }Oslo and Arendal. The {threshold} limit for for contaminants was determined as a quantification limit of seven types of heavy metals, PAH-16, TBT, and PCB. {The applied methodology was optimised to} evaluate and report {the }concentration values in the {contaminants from }soil {stabilised with various binder proportions}. The specimens were treated using geotechnical {and environmental }engineering tests and then evaluated using quantitative methods of the statistical analysis. The comparison of the leaching levels was performed using {estimated} levels over {time }as accumulative leaching {compared to the background level of toxic contaminants}.

\vspace{-2pt}
\begin{figure}[H]
%\begin{adjustwidth}{-\extralength}{0cm}
%\centering
	\includegraphics[width=13.9 cm]{Fig_08.jpg}
%\end{adjustwidth}
\caption{Developed %MDPI: Please check if the blue area need explanation. %<-- Authors' reply: added brief explanation: Blue area represents the significant level of cumulative leaching.
 leaching of TBT from the soil samples over time. Blue area represents the significant level of cumulative leaching.\label{fig08}}
\end{figure}
\vspace{-10pt}

%%%%%%%%%%%%%%%%%%%%%%%%%%%%%%%%%%%%%%%%%%
\section{Conclusions}

{Recent industrialisation has resulted in the increased soil contamination by heavy metals, TBT, PCB, and PAH. The need for decontamination and remediation of soil aimed at removing contaminants in large quantities necessitates the development of effective optimised methods for soil treatment. Using adjusted stabilization workflow with experimentally adjusted proportions of binders, the scheme of soil remediation has been employed in polluted sediments collected from the Port of Oslo (Kongshavn), Norway. This method demonstrated a particular effectiveness for immobilisation of heavy metals, TBT, PAH, and PCB from soil samples. }

We have proposed an {optimised }approach to leaching experiments on soil for physical decontamination from heavy metals and chemical PAH, PCB, and TBT that relies on {processing} the soil samples {using multi-factor} tests. To this end, we have {performed a series of the experiments on} soil testing that allow modelling {the }non-linear behaviour of leaching from soil and {the evaluation of the effects from} the combination of binders. {This method of soil decontamination provided additional }information on the concentration of contaminants {in the sampled soil collected in Kongshavn (Port of Oslo) and enabled the estimation of a prognosis of leaching for one year based on the obtained results simulated for long term release.}

Soil contamination is a significant environmental problem that includes threats to human health and environmental sustainability. At the same time, large construction {and industrial }projects increase {soil} contamination{. This} project contributed to establish the links between the construction engineering and environmental monitoring. {Specifically, }the presented scheme of the workflow enabled to perform accurate decontamination and remediation of soil based on the proposed integrated approaches of {in situ} soil sampling and statistical data analysis. As a result of the implemented framework, soil specimens were processed, stabilized, solidified, and decontaminated. Such workflow {presented} a model for the further processing of large quantities of soil material in harbours and marine ports. {For} environmental purposes, {data} extrapolation was presented for the estimation of possible {leaching }dynamics during one year{. }The visualised examples of {leaching} behaviour {from contaminated soil} are presented {graphically} and {in }tables with provided detailed comments and formulae used for numerical modelling.

Optimization of workflow is a crucial issue in {large-scale} projects, {such }as harbour construction and civil engineering{. }A current limitation of our project is its use of dredged sediments collected from the harbour{. This} has shortcomings, since sampling {of the }marine sediments decreases the availability of suitable disposal specimens and increases {the }transportation and storage costs{. At the same time, }soil processing for industrial projects involves a lengthy, costly, and difficult process and high capital costs. {As} a continuation of this work, we intend to study the soil properties {using samples }collected at different depths. We also plan to investigate how statistical methods could be exploited to model soil collected in large quantities (hundreds of tons) for optimization of large-scale industrial projects such as harbour construction. 

{This study contributes to a better understanding of leaching behaviour in the stabilized soil and can be considered and compared in similar research. The obtained results  can be used as information for the environmental risk assessment and evaluation of soil quality. This article gives an insight on the durability of the treated soil based on the performed environmental engineering workflow. The results can be reused for estimation of the pollution level in the areas of Kongshavn.} As a recommendation for future similar works where large quantity of soil should be processed for industrial construction purposes {with optimised} workflow{, }a pilot project should {always }be carried out to ensure {the reduced cost} of soil {processing} in a large volume before {testing}. Second, when comparing laboratory tests and pilot scale, the temperature of the dredged soil masses should be taken into account {since it} plays a role in the development of soil strength. Moreover, to control soil properties, temperature measurements should be included at different depths. 

%%%%%%%%%%%%%%%%%%%%%%%%%%%%%%%%%%%%%%%%%%
\authorcontributions{Supervision, conceptualization, methodology, visualization, software, resources, data curation, formal analysis, validation, funding acquisition, and project administration, P.L. (Per Lindh); writing---original draft preparation, methodology, software, formal analysis, writing---review and editing, and investigation, P.L. (Polina Lemenkova). All authors have read and agreed to the published version of the manuscript.}

\funding{The publication was funded by the Editorial Office of Algorithms, Multidisciplinary Digital Publishing Institute (MDPI), by providing 100\% discount for the APC of this manuscript. }

%\institutionalreview{Not applicable.}
%
%\informedconsent{Not applicable.}

\dataavailability{Not applicable} 

\acknowledgments{The authors thank the reviewers for reading, suggestions and comments that improved an earlier version of this manuscript.}

\conflictsofinterest{The authors declare no conflicts of interest.} 

\abbreviations{Abbreviations}{
The following abbreviations are used in this manuscript:\\

\noindent 
\begin{tabular}{@{}ll}
BCSA & Belite Calcium Sulfoaluminate \\
BaP & Benz[a]pyrene \\
{BghiP %MDPI: Please check if this is necessary. %<-- Authors' reply: not necessary, corrected (it was a misprint).
} & {Benz[g,h,i]perylene} \\
{DBT} & {Dibutyltin} \\
{DBahA} & {Dibenz[a,h]anthracene} \\
{DOT} & {Dioctyltin} \\

{DPhT} & {Diphenyltin} \\
{IP} & Indeno[1,2,3-cd]pyrene \\
{MBT} & {Monobutyltin} \\
{MOT} & {Monooctyltin} \\ 
{MPhT} & {Monophenyltin} \\
\end{tabular}
\newpage
\setlength{\cellWidtha}{\textwidth/2-2\tabcolsep-2in}

\noindent
\begin{tabular}{@{}ll}

PAH & Polycyclic Aromatic Hydrocarbon \\
PCB & Polychlorinated biphenyl \\
SGI & Swedish Geotechnical Institute \\
TBT & Tributyltin \\
{TPhT} & {Triphenyltin} \\
{TCyT} & {Tricyclohexyltin} \\
{TTBT} & {Tetrabutyltin} \\
\end{tabular}
}

%%%%%%%%%%%%%%%%%%%%%%%%%%%%%%%%%%%%%%%%%%
%% Optional
\appendixtitles{yes} % Leave argument "no" if all appendix headings stay EMPTY (then no dot is printed after "Appendix A"). If the appendix sections contain a heading then change the argument to "yes".
\appendixstart
\appendix
\section[\appendixname~\thesection]{Leaching of Contaminants: Sampling 2317--2324}\label{AppA1}
\subsection[\appendixname~\thesubsection]{Surface Leaching of Contaminants}

Notations for common for all the samples in Table \ref{tabA01}; the applied Methodology is based on standard SS-EN 15863:2015. The start and end days of the tests are from 30 January 2023 until 4 April 2023. The experiments were performed in Swedish Geotechnical Institute. The amount of leachate for all samples is 1.50 L. Conductivity was measured in mS/m by temperature of sampling at 25~$^\circ$C. The period refers to the day of sampling. Surface: 0.0197 m\textsuperscript{2}. Redox (mVolts) refers to oxidation (reduction) potential of solutions or chemical species to acquire/lose electrons and be reduced/oxidised. Leached amount/withdrawal is given in mg/m\textsuperscript{2} for all the samples. 

Abbreviations for organic compounds: Acy---acenaphthylene; BaP---Benz(a)pyrene; B[a]A---Benz[a]anthracene; %MDPI: Please check if [a] should be (a). Please check the whole text and unify the format. %<-- Authors' reply: standard for chemical abbreviation in square brackets (corrected for the whole text): https://pubchem.ncbi.nlm.nih.gov/compound/Benz_a_anthracene
 BjF---Benzo[j]fluoranthene; BkF---Benzo[k]fluoranthene; MBT---Monobutyltin; DBT---Dibutyltin; TBT---Tributyltin; TTBT---Tetrabutyltin; \linebreak MOT---Monooctyltin; DOT---Dioctyltin; TCyT---Tricyclohexyltin; MPhT---Monophenyltin; DPhT---Diphenyltin; IP---Indeno[1,2,3-cd]pyrene; DBahA---Dibenz[a,h]anthracene; BghiP---Benz[g,h,i]perylene; TPhT---Triphenyltin; Total oth. PAH---Total other PAH; Total carc. PAH---Total amount of carcinogenic PAH. 

\vspace{-2pt}
%\begin{footnotesize}
%\begin{scriptsize}
%\begin{center}
%\setlength\LTleft{0pt}
%\setlength\LTright{0pt}
%\begin{longtable}{|p{2.6cm}|cccccccc|}
%\begin{longtable}{@{\extracolsep{\fill}}|p{2.6cm}|cccccccc|@{}}
%\begin{longtable}{@{\extracolsep{\fill}}|p{1.8cm}|cccccccc|@{}}
\begin{table}[H] 
%\begin{threeparttable}
%\centering
\caption{Surface leaching of contaminants (heavy metals, organic compounds, PCB and PAH, mg/m\textsuperscript{2}) from soil specimens during period of 64 days of sampling (2317--2324).} \label{tabA01}

\begin{adjustwidth}{-\extralength}{0cm}
%\centering %% If there is a figure in wide page, please release command \centering
\setlength{\cellWidtha}{\fulllength/9-2\tabcolsep+0.2in}
\setlength{\cellWidthb}{\fulllength/9-2\tabcolsep-.10in}
\setlength{\cellWidthc}{\fulllength/9-2\tabcolsep-0.1in}
\setlength{\cellWidthd}{\fulllength/9-2\tabcolsep-0in}
\setlength{\cellWidthe}{\fulllength/9-2\tabcolsep+0in}
\setlength{\cellWidthf}{\fulllength/9-2\tabcolsep-0in}
\setlength{\cellWidthg}{\fulllength/9-2\tabcolsep-0in}
\setlength{\cellWidthh}{\fulllength/9-2\tabcolsep-0in}
\setlength{\cellWidthi}{\fulllength/9-2\tabcolsep+0in}
\scalebox{1}[1]{\begin{tabularx}{\fulllength}{>{\raggedright\arraybackslash}m{\cellWidtha}>{\centering\arraybackslash}m{\cellWidthb}>{\centering\arraybackslash}m{\cellWidthc}>{\centering\arraybackslash}m{\cellWidthd}>{\centering\arraybackslash}m{\cellWidthe}>{\centering\arraybackslash}m{\cellWidthf}>{\centering\arraybackslash}m{\cellWidthg}>{\centering\arraybackslash}m{\cellWidthh}>{\centering\arraybackslash}m{\cellWidthi}}
\toprule \multicolumn{1}{c}{\textbf{Period}} & {\textbf{6 h}} & {\textbf{1 Days}} & {\textbf{2.25 Days}} & {\textbf{4 Days}} & {\textbf{9 Days}} & {\textbf{16 Days}} & {\textbf{36 Days}} & {\textbf{64 Days}} \\ \midrule 
%\endfirsthead

%\multicolumn{9}{c}%
%{{\bfseries \tablename\ \thetable{} -- continued from previous page}} \\
% \multicolumn{1}{c}{\textbf{Period}} & \multicolumn{1}{c}{\textbf{6 h}} & \multicolumn{1}{c}{\textbf{1 days}} & \multicolumn{1}{c}{\textbf{2.25 days}} & \multicolumn{1}{c}{\textbf{4 days}} & \multicolumn{1}{c}{\textbf{9 days}} & \multicolumn{1}{c}{\textbf{16 days}} & \multicolumn{1}{c}{\textbf{36 days}} & \multicolumn{1}{|c|}{\textbf{64 days}} \\ \hline 
%\endhead

%\hline \multicolumn{9}{|r|}{{Continued on next page}} \\ \hline
%\endfoot

%\hline \hline
%\endlastfoot

Sample No & 2317 & 2318 & 2319 & 2320 & 2321 & 2322 & 2323 & 2324  \\
pH & 8.93 & 9.12 & 8.85 & 9.18 & 9.17 & 8.95 & 8.43 & 8.11  \\
Conductivity & 2510 & 2500 & 2980 & 2760 & 2740 & 2750 & 2770 & 2750  \\
Redox (mV) & 29 & 73 & 33 & 34 & 21 & 19 & 48 & 84  \\
Al & 2.08 & 1.69 & 1.49 & 1.29 & 0.755 & 0.751 & 0.585 & 0.531  \\

As & 0.0553 & 0.0736 & 0.084 & 0.0852 & 0.0807 & 0.35 & 0.0728 & 0.0669  \\
Ba & 6.31 & 7.02 & 8.4 & 7.5 & 13.1 & 13.3 & 21.6 & 22  \\
Ca & 15,500 & 17,000 & 19,600 & 19,300 & 24,300 & 25,300 & 34,600 & 38,200  \\
Cd & <0.004 & <0.004 & <0.004 & <0.004 & <0.004 & <0.004 & <0.004 & <0.004  \\
Co & 0.038 & 0.0306 & 0.0244 & 0.0202 & 0.0333 & 0.0244 & 0.0358 & 0.0508  \\
Cr & <0.04 & <0.04 & <0.04 & <0.04 & <0.04 & <0.04 & <0.04 & <0.04  \\
Cu & 1.58 & 1.03 & 0.744 & 0.532 & 0.653 & 0.5 & 0.86 & 0.634  \\
Fe & 0.898 & 0.739 & 0.664 & 0.486 & <0.3 & <0.3 & <0.3 & <0.3  \\
Hg & <0.002 & <0.002 & <0.002 & <0.002 & <0.002 & <0.002 & <0.002 & <0.002  \\

K & 15,900 & 15,600 & 17,300 & 16,200 & 15,700 & 15,700 & 16,300 & 15,400  \\
Mg & 39,000 & 38,700 & 43,100 & 40,000 & 33,100 & 34,300 & 26,200 & 25,400  \\
Mn & 1.2 & 0.944 & 0.566 & 0.661 & 0.218 & 0.321 & 0.0646 & 0.0921  \\
Mo & 0.807 & 0.709 & 0.807 & 0.727 & 0.91 & 0.776 & 1.13 & 1.18  \\

Na & 374,000 & 374,000 & 414,000 & 390,000 & 368,000 & 370,000 & 356,000 & 360,000  \\
Ni & 0.285 & 0.271 & 0.237 & 0.211 & 0.272 & 0.182 & 0.24 & 0.182  \\
Pb & <0.02 & <0.02 & <0.02 & <0.02 & <0.02 & <0.02 & <0.02 & <0.02  \\

V & 0.153 & 0.212 & 0.24 & 0.282 & 0.391 & 0.264 & 0.155 & 0.139  \\
Zn & 0.18 & <0.2 & <0.2 & 0.179 & <0.2 & <0.2 & <0.2 & <0.2  \\

B & 140 & 142 & 157 & 145 & 126 & 126 & 101 & 101  \\
Sb & 0.0467 & 0.0534 & 0.0496 & 0.0575 & 0.0822 & 0.0688 & 0.11 & 0.116  \\
Se & <0.2 & <0.2 & <0.2 & <0.2 & <0.2 & 7.2 & <0.2 & <0.2  \\
 \bottomrule
\end{tabularx}}
\end{adjustwidth}
\end{table}

\begin{table}[H]\ContinuedFloat
\caption{{\em Cont.}}
\begin{adjustwidth}{-\extralength}{0cm}
%\centering %% If there is a figure in wide page, please release command \centering
\setlength{\cellWidtha}{\fulllength/9-2\tabcolsep+0.2in}
\setlength{\cellWidthb}{\fulllength/9-2\tabcolsep-.10in}
\setlength{\cellWidthc}{\fulllength/9-2\tabcolsep-0.1in}
\setlength{\cellWidthd}{\fulllength/9-2\tabcolsep-0in}
\setlength{\cellWidthe}{\fulllength/9-2\tabcolsep+0in}
\setlength{\cellWidthf}{\fulllength/9-2\tabcolsep-0in}
\setlength{\cellWidthg}{\fulllength/9-2\tabcolsep-0in}
\setlength{\cellWidthh}{\fulllength/9-2\tabcolsep-0in}
\setlength{\cellWidthi}{\fulllength/9-2\tabcolsep+0in}
\scalebox{1}[1]{\begin{tabularx}{\fulllength}{>{\raggedright\arraybackslash}m{\cellWidtha}>{\centering\arraybackslash}m{\cellWidthb}>{\centering\arraybackslash}m{\cellWidthc}>{\centering\arraybackslash}m{\cellWidthd}>{\centering\arraybackslash}m{\cellWidthe}>{\centering\arraybackslash}m{\cellWidthf}>{\centering\arraybackslash}m{\cellWidthg}>{\centering\arraybackslash}m{\cellWidthh}>{\centering\arraybackslash}m{\cellWidthi}}
\toprule \multicolumn{1}{c}{\textbf{Period}} & {\textbf{6 h}} & {\textbf{1 Days}} & {\textbf{2.25 Days}} & {\textbf{4 Days}} & {\textbf{9 Days}} & {\textbf{16 Days}} & {\textbf{36 Days}} & {\textbf{64 Days}} \\ \midrule
Naphthalene & <0.002 & <0.002 & <0.002 & <0.002 & <0.002 & <0.002 & <0.002 & <0.002 \\
Acy & <0.0008 & <0.0008 & <0.0008 & <0.0008 & <0.0008 & <0.0008 & <0.0008 & <0.0008  \\
Acenaphthene & <0.0008 & <0.0008 & <0.0008 & <0.0008 & 0.00099 & 0.00084 & <0.0008 & <0.0008  \\
Fluorine & <0.0008 & <0.0008 & <0.0008 & <0.0008 & <0.0008 & <0.0008 & <0.0008 & <0.0008  \\
Phenanthrene & <0.002 & <0.002 & <0.002 & <0.002 & <0.002 & <0.002 & <0.002 & <0.002  \\
Anthracene & <0.0008 & <0.0008 & <0.0008 & <0.0008 & <0.0008 & <0.0008 & <0.0008 & <0.0008  \\
Fluoranthene & <0.0008 & <0.0008 & <0.0008 & <0.00084 & 0.002 & 0.0014 & 0.0019 & <0.0008  \\
Pyrene & <0.0008 & 0.0013 & 0.002 & 0.0021 & 0.0037 & 0.004 & 0.0052 & 0.0028  \\
B[a]A & <0.0008 & <0.0008 & <0.0008 & <0.0008 & <0.0008 & <0.0008 & <0.0008 & <0.0008  \\
Chrysene & <0.0008 & <0.0008 & <0.0008 & <0.0008 & <0.0008 & <0.0008 & <0.0008 & <0.0008  \\
BjF  & <0.0008 & <0.0008 & <0.0008 & <0.0008 & <0.0008 & <0.0008 & <0.0008 & <0.0008  \\
BkF & <0.0008 & <0.0008 & <0.0008 & <0.0008 & <0.0008 & <0.0008 & <0.0008 & <0.0008  \\
BaP & <0.0008 & <0.0008 & <0.0008 & <0.0008 & <0.0008 & <0.0008 & <0.0008 & <0.0008  \\
DBahA & <0.0008 & <0.0008 & <0.0008 & <0.0008 & <0.0008 & <0.0008 & <0.0008 & <0.0008  \\
BghiP & <0.0008 & <0.0008 & <0.0008 & <0.0008 & <0.0008 & <0.0008 & <0.0008 & <0.0008  \\
IP & <0.0008 & <0.0008 & <0.0008 & <0.0008 & <0.0008 & <0.0008 & <0.0008 & <0.0008  \\
Total PAH 16 & <0.0072 & 0.0013 & 0.002 & 0.0029 & 0.0062 & 0.006 & 0.0071 & 0.0028  \\
Total carc. PAH & <0.0027 & <0.0027 & <0.0027 & <0.0027 & <0.0027 & <0.0027 & <0.0027 & <0.0027  \\
Total oth. PAH & <0.005 & 0.0013 & 0.002 & 0.0029 & 0.0062 & 0.006 & 0.0071 & 0.0028  \\
Total PAH L & <0.0019 & <0.0019 & <0.0019 & <0.0019 & 0.00099 & 0.00084 & <0.0019 & <0.0019  \\
Total PAH M & <0.002 & 0.0013 & 0.002 & 0.0029 & 0.0052 & 0.0052 & 0.0071 & 0.0028  \\
Total PAH H & <0.003 & <0.003 & <0.003 & <0.003 & <0.003 & <0.003 & <0.003 & <0.003  \\
MBT & 0.000404 & 0.00326 & 0.00202 & 0.00223 & 0.00536 & 0.0055 & 0.0044 & 0.00711 \\
DBT & <0.00008 & <0.00008 & <0.00008 & <0.00008 & 0.000142 & 0.000112 & 0.000153 & 0.00024 \\
TBT & 0.0000807 & 0.0000951 & <0.00008 & 0.0000883 & 0.000101 & 0.000104 & 0.0000799 & 0.000133 \\
TTBT & <0.00008 & <0.00008 & <0.00008 & <0.00008 & <0.00008 & <0.00008 & <0.00008 & <0.00008 \\
MOT & <0.00008 & <0.00008 & <0.00008 & <0.00008 & <0.00008 & <0.00008 & <0.00008 & <0.00008 \\
DOT & <0.00008 & <0.00008 & <0.00008 & <0.00008 & <0.00008 & <0.00008 & <0.00008 & <0.00008 \\
TCyT & <0.00008 & <0.00008 & <0.00008 & <0.00008 & <0.00008 & <0.00008 & <0.00008 & <0.00008 \\
MPhT &  <0.00008 &  <0.00008 &  <0.00008 &  <0.00008 &  <0.00008 &  <0.00008 &  <0.00008 &  <0.00008 \\
DPhT & <0.00008 & <0.00008 & <0.00008 & <0.00008 & <0.00008 & <0.00008 & <0.00008 & <0.00008  \\
TPhT &  <0.00008 &  <0.00008 &  <0.00008 &  <0.00008 &  <0.00008 &  <0.00008 &  <0.00008 &  <0.00008 \\
PCB 28 & <0.000084 & <0.000084 & <0.000084 & <0.000084 & <0.000084 & <0.000084 & <0.000084 & <0.000084 \\
PCB 52 & <0.000084 & <0.000084 & <0.000084 & <0.000084 & <0.000084 & <0.000084 & <0.000084 & <0.000084 \\
PCB 101 & <0.000084 & <0.000084 & <0.000084 & <0.000084 & <0.000084 & <0.000084 & <0.000084 & <0.000084  \\
PCB 118 & <0.000084 & <0.000084 & <0.000084 & <0.000084 & <0.000084 & <0.000084 & <0.000084 & <0.000084  \\
PCB 138 & <0.000091 & <0.000091 & <0.000091 & <0.000091 & <0.000091 & <0.000091 & <0.000091 & <0.000091  \\
PCB 153 & <0.000084 & <0.000084 & <0.000084 & <0.000084 & <0.000084 & <0.000084 & <0.000084 & <0.000084  \\
PCB 180  & <0.000084 & <0.000084 & <0.000084 & <0.000084 & <0.000084 & <0.000084 & <0.000084 & <0.000084  \\
Total PCB 7 & <0.0003 & <0.0003 & <0.0003 & <0.0003 & <0.0003 & <0.0003 & <0.0003 & <0.0003  \\
%
\bottomrule
\end{tabularx}}
\end{adjustwidth}
\end{table}
%\end{center}
%\end{footnotesize}
%\end{scriptsize}
\vspace{-14pt}
%------------------>
\subsection[\appendixname~\thesubsection]{Cumulative Surface Leaching of Contaminants}

Notations for Table \ref{tabA02}: The applied methodology is based on the three approaches: 

\begin{itemize}
    \item Standards SS-EN ISO 10523:2012 with measurement uncertainty of $\pm$ 0.05 pH-enh;
    \item SS-EN 27888-1:1994 with measurement uncertainty $\pm$2.9\%;
    \item SGI-method with measurement uncertainty pf $\pm$ 25\%.
\end{itemize}

Abbreviations for organic compounds in Table \ref{tabA02}: Acy---acenaphthylene; \linebreak BaP---Benz[a]pyrene; B[a]A---Benz[a]anthracene; BjF---Benzo[j]fluoranthene; \linebreak BghiP---Benz[g,h,i]perylene; BkF---Benzo[k]fluoranthene; MBT---Monobutyltin; DBT---Dibutyltin; TBT---Tributyltin; TTBT---Tetrabutyltin; MOT---Monooctyltin; DOT---Dioctyltin; TCyT---Tricyclohexyltin; MPhT---Monophenyltin; DPhT---Diphenyltin; IP---Indeno[1,2,3-cd]pyrene; DBahA---Dibenz[a,h]anthracene; TPhT---Triphenyltin; Total oth. PAH---Total other PAH; Total carc. PAH---Total amount of carcinogenic PAH. 

%--------------------------------------------------------->
%\begin{scriptsize}
%\begin{center}
%\setlength\LTleft{0pt}
%\setlength\LTright{0pt}
%%\begin{longtable}{|p{2.6cm}|cccccccc|}
%%\begin{longtable}{@{\extracolsep{\fill}}|p{2.6cm}|cccccccc|@{}}
%\begin{longtable}{@{\extracolsep{\fill}}|p{1.8cm}|cccccccc|@{}}
\begin{table}[H] 
%\begin{threeparttable}
%\centering
\caption{Cumulative surface leaching (heavy metals, organic compounds, PCB and PAH, mg/m\textsuperscript{2}) from soil specimens during period of 64 days of sampling (2317--2324).} \label{tabA02}
\begin{adjustwidth}{-\extralength}{0cm}
%\centering %% If there is a figure in wide page, please release command \centering

\setlength{\cellWidtha}{\fulllength/9-2\tabcolsep+0.2in}
\setlength{\cellWidthb}{\fulllength/9-2\tabcolsep-.10in}
\setlength{\cellWidthc}{\fulllength/9-2\tabcolsep-0.1in}
\setlength{\cellWidthd}{\fulllength/9-2\tabcolsep-0in}
\setlength{\cellWidthe}{\fulllength/9-2\tabcolsep+0in}
\setlength{\cellWidthf}{\fulllength/9-2\tabcolsep-0in}
\setlength{\cellWidthg}{\fulllength/9-2\tabcolsep-0in}
\setlength{\cellWidthh}{\fulllength/9-2\tabcolsep-0in}
\setlength{\cellWidthi}{\fulllength/9-2\tabcolsep+0in}
\scalebox{1}[1]{\begin{tabularx}{\fulllength}{>{\raggedright\arraybackslash}m{\cellWidtha}>{\centering\arraybackslash}m{\cellWidthb}>{\centering\arraybackslash}m{\cellWidthc}>{\centering\arraybackslash}m{\cellWidthd}>{\centering\arraybackslash}m{\cellWidthe}>{\centering\arraybackslash}m{\cellWidthf}>{\centering\arraybackslash}m{\cellWidthg}>{\centering\arraybackslash}m{\cellWidthh}>{\centering\arraybackslash}m{\cellWidthi}}
\toprule
 \multicolumn{1}{c}{\textbf{Period}} & \multicolumn{1}{c}{\textbf{6 h}} & \multicolumn{1}{c}{\textbf{1 Days}} & \multicolumn{1}{c}{\textbf{2.25 Days}} & \multicolumn{1}{c}{\textbf{4 Days}} & \multicolumn{1}{c}{\textbf{9 Days}} & \multicolumn{1}{c}{\textbf{16 Days}} & \multicolumn{1}{c}{\textbf{36 Days}} & \multicolumn{1}{c}{\textbf{64 Days}} \\ \midrule 
%\endfirsthead

%\multicolumn{9}{c}%
%{{\bfseries \tablename\ \thetable{} -- continued from previous page}} \\
%\hline \multicolumn{1}{|c|}{\textbf{Period}} & \multicolumn{1}{c}{\textbf{6 h}} & \multicolumn{1}{c}{\textbf{1 days}} & \multicolumn{1}{c}{\textbf{2.25 days}} & \multicolumn{1}{c}{\textbf{4 days}} & \multicolumn{1}{c}{\textbf{9 days}} & \multicolumn{1}{c}{\textbf{16 days}} & \multicolumn{1}{c}{\textbf{36 days}} & \multicolumn{1}{c}{\textbf{64 days}} \\ \hline 
%\endhead
%
%\hline \multicolumn{9}{|r|}{{Continued on next page}} \\ \hline
%\endfoot
%
%\hline \hline
%\endlastfoot

Sample No & 2317 & 2318 & 2319 & 2320 & 2321 & 2322 & 2323 & 2324  \\
Al & 2.08 & 3.77 & 5.26 & 6.55 & 7.3 & 8.05 & 8.64 & 9.17 \\
As & 0.0553 & 0.129 & 0.21 & 0.298 & 0.379 & 0.73 & 0.802 & 0.868 \\
Ba & 6.31 & 13.3 & 22 & 29.2 & 42.3 & 55.6 & 77.2 & 99 \\
Ca & 15,500 & 32,000 & 51,900 & 71,200 & 95,500 & 121,000 & 155,000 & 194,000 \\
Cd & <0.004 & <0.008 & <0.01 & <0.02 & <0.02 & <0.02 & <0.03 & <0.03 \\
Co & 0.038 & 0.0686 & 0.093 & 0.113 & 0.147 & 0.171 & 0.207 & 0.258 \\
Cr & <0.04 & <0.08 & <0.1 & <0.2 & <0.1 & <0.1 & <0.3 & <0.3 \\
Cu & 1.58 & 2.6 & 3.35 & 3.88 & 4.53 & 5 & 5.89 & 6.53 \\
Fe & 0.898 & 1.64 & 2.3 & 2.79 & <3 & <3 & <4 & <4 \\
Hg & <0.002 &  <0.003 & <0.005 &  <0.006 & <0.008 & <0.009 & <0.01 & <0.01 \\
K & 15,900 & 31,500 & 48,800 & 65,000 & 80,700 & 96,400 & 113,000 & 128,000 \\
Mg & 39,000 & 77,500 & 121,000 & 161,000 & 194,000 & 228,000 & 254,000 & 280,000 \\
Mn & 1.2 & 2.16 & 2.73 & 3.39 & 3.61 & 3.93 & 3.99 & 4.08  \\
Mo & 0.807 & 1.52 & 2.32 & 3.05 & 4 & 4.74 & 5.87 & 7.05  \\
Na & 374,000 & 748,000 & 1,160,000 & 1,550,000 & 1,920,000 & 2,290,000 & 2,650,000 & 3,000,000 \\
Ni & 0.285 & 0.556 & 0.793 & 1 & 1.28 & 1.46 & 1.7 & 1.88 \\
Pb & <0.02 & <0.03 & <0.05 & <0.06 & <0.08 & <0.09 & <0.1 & <0.1 \\
V & 0.153 & 0.365 & 0.605 & 0.887 & 1.28 & 1.54 & 1.7 & 1.84 \\
Zn & 0.18 & <0.3 & <0.5 & 0.666 & <0.8 & <1 & <1 & <1 \\
B & 140 & 282 & 438 & 584 & 710 & 836 & 937 & 1040 \\
Sb & 0.0467 & 0.1 & 0.15 & 0.207 & 0.289 & 0.358 & 0.468 & 0.584 \\
Se & <0.2 & <0.5 & <0.7 & <0.9 & <1 & 8.34 & <9 & <9 \\
Naphthalene & <0.002 & <0.005 & <0.007 & <0.009 & <0.01 & <0.01 & <0.02 & <0.02 \\
Acy & <0.0008 & <0.002 & <0.002 & <0.003 & <0.004 & <0.005 & <0.005 & <0.006 \\
Acenaphthene & <0.0008 & <0.002 & <0.002 & <0.003 & <0.004 & 0.0049 & <0.006 & <0.006 \\
Fluorine & <0.0008 &  <0.002 &  <0.002 &  <0.003 &  <0.004 &  <0.005 &  <0.005 &  <0.006 \\
Phenanthrene & <0.002 & <0.003 & <0.005 & <0.006 & <0.008 & <0.009 & <0.01 & <0.01 \\
Anthracene & <0.0008 & <0.002 & <0.002 & <0.003 & <0.004 & <0.005 & <0.002 & <0.006 \\
Fluoranthene & <0.0008 & <0.002 & <0.002 & 0.0031 & 0.005 & 0.006 & 0.0079 & <0.009 \\
Pyrene & <0.0008 & 0.0021 & 0.004 & 0.0056 & 0.0093 & 0.01 & 0.018 & 0.021 \\

B[a]A & <0.0008 & <0.002 & <0.002 & <0.003 & <0.004 & <0.005 & <0.005 & <0.006 \\
Chrysene & <0.0008 & <0.002 & <0.002 & <0.003 & <0.004 & <0.005 & <0.005 & <0.006 \\
BjF & <0.0008 & <0.002 & <0.002 & <0.003 & <0.004 & <0.005 & <0.005 & <0.006 \\
BkF & <0.0008 & <0.002 & <0.002 & <0.003 & <0.004 & <0.005 & <0.005 & <0.006 \\
BaP & <0.0008 & <0.002 & <0.002 & <0.003 & <0.004 & <0.005 & <0.005 & <0.006 \\
DBahA & <0.0008 & <0.002 & <0.002 & <0.003 & <0.004 & <0.005 & <0.005 & <0.006 \\
BghiP & <0.0008 & <0.002 & <0.002 & <0.003 & <0.004 & <0.005 & <0.005 & <0.006 \\

IP & <0.0008 & <0.002 & <0.002 & <0.003 & <0.004 & <0.005 & <0.005 & <0.006 \\
Total PAH 16 & <0.0072 & 0.0085 & 0.01 & 0.013 & 0.019 & 0.025 & 0.032 & 0.035 \\
Total carc. PAH & <0.0027 & <0.0053 & <0.008 & <0.011 & <0.013 & <0.016 & <0.019 & <0.021 \\
Total oth. PAH & <0.005 & 0.0059 & 0.007 & 0.01 & 0.016 & 0.022 & 0.03 & 0.032 \\

Total PAH L & <0.0019 & <0.0038 & <0.0057 & <0.0076 & 0.0086 & 0.0094 & <0.011 & <0.013 \\
Total PAH M & <0.002 & 0.0036 & 0.005 & 0.008 & 0.013 & 0.018 & 0.025 & 0.028 \\
Total PAH H & <0.003 & <0.006 & <0.009 & <0.01 & <0.02 & <0.02 & <0.02 & <0.02 \\
MBT & 0.000404 & 0.00366 & 0.00569 & 0.00792 & 0.0133 & 0.0188 & 0.0232 & 0.0303 \\
DBT & <0.00008 & <0.0002 & <0.0002 & <0.0003 & 0.000446 & 0.000558 & 0.000711 & 0.00095 \\
TBT & 0.0000807 & 0.000176 & <0.0003 & 0.00034 & 0.000441 & 0.000545 & 0.000625 & 0.000758 \\
TTBT & <0.00008 & <0.0002 & <0.0002 & <0.0003 & <0.0004 & <0.0005 &<0.0005 & <0.0006 \\
MOT & <0.00008 & <0.0002 & <0.0002 & <0.0003 & <0.0004 & <0.0005 &<0.0005 & <0.0006 \\
DOT & <0.00008 & <0.0002 & <0.0002 & <0.0003 & <0.0004 & <0.0005 &<0.0005 & <0.0006 \\
TCyT & <0.00008 & <0.0002 & <0.0002 & <0.0003 & <0.0004 & <0.0005 &<0.0005 & <0.0006 \\
MPhT & <0.00008 & <0.0002 & <0.0002 & <0.0003 & <0.0004 & <0.0005 &<0.0005 & <0.0006 \\
DPhT & <0.00008 & <0.0002 & <0.0002 & <0.0003 & <0.0004 & <0.0005 &<0.0005 & <0.0006 \\
TPhT & <0.00008 & <0.0002 & <0.0002 & <0.0003 & <0.0004 & <0.0005 &<0.0005 & <0.0006 \\
\bottomrule
\end{tabularx}}
\end{adjustwidth}
\end{table}

\begin{table}[H]\ContinuedFloat
%\tablesize{\small}
\caption{{\em Cont.}}
\begin{adjustwidth}{-\extralength}{0cm}
%\centering %% If there is a figure in wide page, please release command \centering

\setlength{\cellWidtha}{\fulllength/9-2\tabcolsep+0.2in}
\setlength{\cellWidthb}{\fulllength/9-2\tabcolsep-.10in}
\setlength{\cellWidthc}{\fulllength/9-2\tabcolsep-0.1in}
\setlength{\cellWidthd}{\fulllength/9-2\tabcolsep-0in}
\setlength{\cellWidthe}{\fulllength/9-2\tabcolsep+0in}
\setlength{\cellWidthf}{\fulllength/9-2\tabcolsep-0in}
\setlength{\cellWidthg}{\fulllength/9-2\tabcolsep-0in}
\setlength{\cellWidthh}{\fulllength/9-2\tabcolsep-0in}
\setlength{\cellWidthi}{\fulllength/9-2\tabcolsep+0in}
\scalebox{1}[1]{\begin{tabularx}{\fulllength}{>{\raggedright\arraybackslash}m{\cellWidtha}>{\centering\arraybackslash}m{\cellWidthb}>{\centering\arraybackslash}m{\cellWidthc}>{\centering\arraybackslash}m{\cellWidthd}>{\centering\arraybackslash}m{\cellWidthe}>{\centering\arraybackslash}m{\cellWidthf}>{\centering\arraybackslash}m{\cellWidthg}>{\centering\arraybackslash}m{\cellWidthh}>{\centering\arraybackslash}m{\cellWidthi}}
\toprule
 \multicolumn{1}{c}{\textbf{Period}} & \multicolumn{1}{c}{\textbf{6 h}} & \multicolumn{1}{c}{\textbf{1 Days}} & \multicolumn{1}{c}{\textbf{2.25 Days}} & \multicolumn{1}{c}{\textbf{4 Days}} & \multicolumn{1}{c}{\textbf{9 Days}} & \multicolumn{1}{c}{\textbf{16 Days}} & \multicolumn{1}{c}{\textbf{36 Days}} & \multicolumn{1}{c}{\textbf{64 Days}} \\ \midrule 
PCB 28 & <0.000084 & <0.00017 & <0.00025 & <0.00033 & <0.00042 & <0.0005 & <0.00059 & <0.00067 \\
PCB 52 & <0.000084 & <0.00017 & <0.00025 & <0.00033 & <0.00042 & <0.0005 & <0.00059 & <0.00067 \\
PCB 101 & <0.000084 & <0.00017 & <0.00025 & <0.00033 & <0.00042 & <0.0005 & <0.00059 & <0.00067 \\
PCB 118 & <0.000084 & <0.00017 & <0.00025 & <0.00033 & <0.00042 & <0.0005 & <0.00059 & <0.00067 \\
PCB 138 & <0.000091 & <0.00018 & <0.00027 & <0.00037 & <0.00046 & <0.00055 & <0.00064 & <0.00073 \\
PCB 153 & <0.000084 & <0.00017 & <0.00025 & <0.00033 & <0.00042 & <0.0005 & <0.00059 & <0.00067 \\
PCB 180 & <0.000084 & <0.00017 & <0.00025 & <0.00033 & <0.00042 & <0.0005 & <0.00059 & <0.00067 \\
Total PCB 7 & <0.0003 & <0.00059 & <0.00089 & <0.0012 & <0.0015 & <0.0018 & <0.0021 & <0.0024 \\
\bottomrule
\end{tabularx}}
\end{adjustwidth}
\end{table}
%\end{center}
%%\end{footnotesize}
%\end{scriptsize}
%---------------------------------------------------------<
\vspace{-14pt}
\section[\appendixname~\thesection]{Leaching of Contaminants: Sampling 2325--2332}\label{AppB1}
\subsection[\appendixname~\thesubsection]{Surface Leaching of Contaminants}

Notations for common for all the samples in Table \ref{tabA03}. The applied Methodology contains Batch 1, Prov 7. The start and end days of the tests are from 30 January 2023 until \mbox{4 April 2023}. The experiments were performed by Robert Selegård in Swedish Geotechnical Institute, registered under Diary Nr. 1.1-2107-0587. The amount of leachate for all samples is 1.50 L. Conductivity was measured in mS/m by temperature of sampling at 25~$^\circ$C. The period refers to the day of sampling. Surface: 0.0197 m\textsuperscript{2}. Redox (mVolts) refers to oxidation (reduction) potential of solutions or chemical species to acquire/lose electrons and be reduced/oxidised. Leached amount/withdrawal is given in mg/m\textsuperscript{2} for all the samples. 

Abbreviations for organic compounds: Acy---acenaphthylene; BaP---Benz[a]pyrene; B[a]A---Benz[a]anthracene; BjF---Benzo[j]fluoranthene; BkF---Benzo[k]fluoranthene; MBT---Monobutyltin; DBT---Dibutyltin; TBT---Tributyltin; TTBT---Tetrabutyltin; \linebreak MOT---Monooctyltin; DOT---Dioctyltin; TCyT---Tricyclohexyltin; MPhT---Monophenyltin; DPhT---Diphenyltin; IP---Indeno[1,2,3-cd]pyrene; DBahA---Dibenz[a,h]anthracene; BghiP---Benz[g,h,i]perylene; TPhT---Triphenyltin; Total oth. PAH---Total other PAH; Total carc. PAH---Total amount of carcinogenic PAH. 
%--------------------------------------------------------->
%\begin{scriptsize}
%\begin{center}
%\setlength\LTleft{0pt}
%\setlength\LTright{0pt}
%\begin{longtable}{|p{2.6cm}|cccccccc|}
%\begin{longtable}{@{\extracolsep{\fill}}|p{2.6cm}|cccccccc|@{}}
%\begin{longtable}{@{\extracolsep{\fill}}|p{1.8cm}|cccccccc|@{}}

\vspace{-2pt}
\begin{table}[H] 
%\begin{threeparttable}
%\centering
\caption{Surface leaching of contaminants (heavy metals, organic compounds, PCB and PAH, mg/m\textsuperscript{2}) from soil specimens during period of 64 days of sampling (2325--2332).} \label{tabA03}
\begin{adjustwidth}{-\extralength}{0cm}
%\centering %% If there is a figure in wide page, please release command \centering

\setlength{\cellWidtha}{\fulllength/9-2\tabcolsep+0.2in}
\setlength{\cellWidthb}{\fulllength/9-2\tabcolsep-.10in}
\setlength{\cellWidthc}{\fulllength/9-2\tabcolsep-0.1in}
\setlength{\cellWidthd}{\fulllength/9-2\tabcolsep-0in}
\setlength{\cellWidthe}{\fulllength/9-2\tabcolsep+0in}
\setlength{\cellWidthf}{\fulllength/9-2\tabcolsep-0in}
\setlength{\cellWidthg}{\fulllength/9-2\tabcolsep-0in}
\setlength{\cellWidthh}{\fulllength/9-2\tabcolsep-0in}
\setlength{\cellWidthi}{\fulllength/9-2\tabcolsep+0in}
\scalebox{1}[1]{\begin{tabularx}{\fulllength}{>{\raggedright\arraybackslash}m{\cellWidtha}>{\centering\arraybackslash}m{\cellWidthb}>{\centering\arraybackslash}m{\cellWidthc}>{\centering\arraybackslash}m{\cellWidthd}>{\centering\arraybackslash}m{\cellWidthe}>{\centering\arraybackslash}m{\cellWidthf}>{\centering\arraybackslash}m{\cellWidthg}>{\centering\arraybackslash}m{\cellWidthh}>{\centering\arraybackslash}m{\cellWidthi}}
\toprule \multicolumn{1}{c}{\textbf{Period}} & \multicolumn{1}{c}{\textbf{6 h}} & \multicolumn{1}{c}{\textbf{1 Days}} & \multicolumn{1}{c}{\textbf{2.25 Days}} & \multicolumn{1}{c}{\textbf{4 Days}} & \multicolumn{1}{c}{\textbf{9 Days}} & \multicolumn{1}{c}{\textbf{16 Days}} & \multicolumn{1}{c}{\textbf{36 Days}} & \multicolumn{1}{c}{\textbf{64 Days}} \\ \midrule 
%\endfirsthead
%
%\multicolumn{9}{c}%
%{{\bfseries \tablename\ \thetable{} -- continued from previous page}} \\
%\hline \multicolumn{1}{|c}{\textbf{Period}} & \multicolumn{1}{c}{\textbf{6 h}} & \multicolumn{1}{c}{\textbf{1 days}} & \multicolumn{1}{c}{\textbf{2.25 days}} & \multicolumn{1}{c}{\textbf{4 days}} & \multicolumn{1}{c}{\textbf{9 days}} & \multicolumn{1}{c}{\textbf{16 days}} & \multicolumn{1}{c}{\textbf{36 days}} & \multicolumn{1}{c}{\textbf{64 days}} \\ \hline 
%\endhead
%
%\hline \multicolumn{9}{|r|}{{Continued on next page}} \\ \hline
%\endfoot

%\hline \hline
%\endlastfoot

Sample No & 2325 & 2326 & 2327 & 2328 & 2329 & 2330 & 2331 & 2332 \\
pH & 8.94 & 9.21 & 9.1 & 9.2 & 9.22 & 9.06 & 8.42 & 8.21 \\
Conductivity & 2530 & 2510 & 2840 & 2860 & 2600 & 2760 & 2730 & 2740 \\
Redox (mV) & 36 & 70 & 36 & 39 & 30 & 27 & 58 & 87 \\
Al & 1.86 & 1.73 & 1.4 & 1.34 & 0.68 & 0.61 & 0.417 & 0.387 \\
As & 0.0581 & 0.0753 & 0.0859 & 0.0642 & 0.0782 & 0.0821 & 0.0569 & 0.0828 \\
Ba & 6.04 & 6.89 & 7.66 & 6.65 & 13.1 & 13.4 & 21 & 20.1 \\
Ca & 15,900 & 17,400 & 19,600 & 19,000 & 24,200 & 24,300 & 32,300 & 35,000 \\
Cd & 0.0042 & 0.00505 & <0.004 & <0.004 & <0.004 & <0.004 & <0.004 & 0.00431 \\
Co & 0.037 & 0.0295 & 0.0239 & 0.0266 & 0.0419 & 0.0271 & 0.0292 & 0.0413 \\
Cr & <0.04 & <0.04 & <0.04 & <0.04 & <0.04 & <0.04 & <0.04 & <0.04  \\
Cu & 1.3 & 0.936 & 0.674 & 0.539 & 0.599 & 0.594 & 0.89 & 0.967 \\

Fe & 0.859 & 0.66 & 0.59 & 0.463 & <0.3 & <0.3 & <0.3 & <0.3 \\
Hg & <0.002 & <0.002 & <0.002 & <0002 & <0.002 & <0.002 & <0.002 & <0.002 \\
K & 16,000 & 15,600 & 17,500 & 16,300 & 16,000 & 15,400 & 16,200 & 15,500 \\
Mg & 39,400 & 39,400 & 43,600 & 40,700 & 34,000 & 34,300 & 26,700 & 27,200 \\
Mn & 1.33 & 1.03 & 0.542 & 0.561 & 0.196 & 0.271 & 0.0798 & 0.0427 \\
Mo & 0.782 & 0.72 & 0.782 & 0.68 & 0.898 & 0.798 & 1.1 & 1.17 \\
Na & 380,000 & 371,000 & 418,000 & 394,000 & 371,000 & 364,000 & 350,000 & 359,000 \\
Ni & 0.411 & 0.315 & 0.224 & 0.301 & 0.411 & 0.233 & 0.173 & 0.158 \\
\bottomrule
\end{tabularx}}
\end{adjustwidth}
\end{table}

\begin{table}[H]\ContinuedFloat
\tablesize{\small}
\caption{{\em Cont.}}\label{tabA03}
\begin{adjustwidth}{-\extralength}{0cm}
%\centering %% If there is a figure in wide page, please release command \centering

\setlength{\cellWidtha}{\fulllength/9-2\tabcolsep+0.2in}
\setlength{\cellWidthb}{\fulllength/9-2\tabcolsep-.10in}
\setlength{\cellWidthc}{\fulllength/9-2\tabcolsep-0.1in}
\setlength{\cellWidthd}{\fulllength/9-2\tabcolsep-0in}
\setlength{\cellWidthe}{\fulllength/9-2\tabcolsep+0in}
\setlength{\cellWidthf}{\fulllength/9-2\tabcolsep-0in}
\setlength{\cellWidthg}{\fulllength/9-2\tabcolsep-0in}
\setlength{\cellWidthh}{\fulllength/9-2\tabcolsep-0in}
\setlength{\cellWidthi}{\fulllength/9-2\tabcolsep+0in}
\scalebox{1}[1]{\begin{tabularx}{\fulllength}{>{\raggedright\arraybackslash}m{\cellWidtha}>{\centering\arraybackslash}m{\cellWidthb}>{\centering\arraybackslash}m{\cellWidthc}>{\centering\arraybackslash}m{\cellWidthd}>{\centering\arraybackslash}m{\cellWidthe}>{\centering\arraybackslash}m{\cellWidthf}>{\centering\arraybackslash}m{\cellWidthg}>{\centering\arraybackslash}m{\cellWidthh}>{\centering\arraybackslash}m{\cellWidthi}}
\toprule \multicolumn{1}{c}{\textbf{Period}} & \multicolumn{1}{c}{\textbf{6 h}} & \multicolumn{1}{c}{\textbf{1 Days}} & \multicolumn{1}{c}{\textbf{2.25 Days}} & \multicolumn{1}{c}{\textbf{4 Days}} & \multicolumn{1}{c}{\textbf{9 Days}} & \multicolumn{1}{c}{\textbf{16 Days}} & \multicolumn{1}{c}{\textbf{36 Days}} & \multicolumn{1}{c}{\textbf{64 Days}} \\ \midrule 
Pb & <0.02 & <0.02 & <0.02 & <0.02 & <0.02 & <0.02 & <0.02 & <0.02 \\
V & 0.141 & 0.21 & 0.243 & 0.242 & 0.351 & 0.24 & 0.15 & 0.116 \\
Zn & 0.278 & <0.2 & <0.2 & <0.2 & <0.2 & <0.2 & <0.2 & <0.2 \\
B & 144 & 144 & 157 & 149 & 128 & 124 & 101 & 104 \\

Sb & 0.044 & 0.0485 & 0.0501 & 0.0511 & 0.0882 & 0.0725 & 0.114 & 0.11 \\
Se & <0.2 & <0.2 & <0.2 & <0.2 & <0.2 & 4.2 & <0.2 & <0.2 \\
Naphthalene & <0.002 & <0.002 & <0.002 & <0.002 & <0.002 & <0.002 & <0.002 & <0.002 \\
Acy & <0.0008 & <0.0008 & <0.0008 & <0.0008 & <0.0008 & <0.0008 & <0.0008 & <0.0008  \\
Acenaphthene & <0.0008 & <0.0008 & <0.0008 & <0.0008 & 0.00092 & 0.0008 & <0.0008 & <0.0008 \\
Fluorine & <0.0008 & <0.0008 & <0.0008 & <0.0008 & <0.0008 & <0.0008 & <0.0008 & <0.0008 \\
Phenanthrene & <0.002 & <0.002 & <0.002 & <0.002 & <0.002 & <0.002 & <0.002 & <0.002 \\
Anthracene & <0.0008 & <0.0008 & <0.0008 & <0.0008 & <0.0008 & <0.0008 & <0.0008 & <0.0008 \\
Fluoranthene & <0.0008 & <0.0008 & <0.0008 & <0.0008 & 0.0012 & 0.0012 & 0.0018 & 0.001 \\
Pyrene & <0.0008 & 0.0012 & 0.0017 & 0.002 & 0.0035 & 0.0037 & 0.0053 & 0.0044 \\

B[a]A & <0.0008 & <0.0008 & <0.0008 & <0.0008 & <0.0008 & <0.0008 & <0.0008 & <0.0008  \\
Chrysene & <0.0008 & <0.0008 & <0.0008 & <0.0008 & <0.0008 & <0.0008 & <0.0008 & <0.0008  \\
BjF & <0.0008 & <0.0008 & <0.0008 & <0.0008 & <0.0008 & <0.0008 & <0.0008 & <0.0008  \\
BkF & <0.0008 & <0.0008 & <0.0008 & <0.0008 & <0.0008 & <0.0008 & <0.0008 & <0.0008  \\
BaP & <0.0008 & <0.0008 & <0.0008 & <0.0008 & <0.0008 & <0.0008 & <0.0008 & <0.0008  \\
DBahA & <0.0008 & <0.0008 & <0.0008 & <0.0008 & <0.0008 & <0.0008 & <0.0008 & <0.0008  \\
BghiP & <0.0008 & <0.0008 & <0.0008 & <0.0008 & <0.0008 & <0.0008 & <0.0008 & <0.0008  \\
IP & <0.0008 & <0.0008 & <0.0008 & <0.0008 & <0.0008 & <0.0008 & <0.0008 & <0.0008  \\
Total PAH 16 & <0.0073 & 0.0012 & 0.0017 & 0.002 & 0.0057 & 0.0057 & 0.0071 & 0.0054  \\
Total carc. PAH & <0.0027 &  <0.0027 &  <0.0027 &  <0.0027 &  <0.0027 &  <0.0027 &  <0.0027 &  <0.0027 \\
Total oth. PAH & <0.005 & 0.0012 & 0.0017 & 0.002 & 0.0057 & 0.0057 & 0.0071 & 0.0054 \\
Total PAH L & <0.0019 & <0.0019 & <0.0019 & <0.0019 & 0.00092 & 0.0008 & <0.0019 & <0.0019 \\
Total PAH M & <0.002 & 0.0012 & 0.0017 & 0.002 & 0.0048 & 0.0049 & 0.0071 & 0.0054 \\
Total PAH H & <0.003 & <0.003 & <0.003 & <0.003 & <0.003 & <0.003 & <0.003 & <0.003 \\
MBT & 0.000333 & 0.0027 & 0.00229 & 0.00335 & 0.00398 & 0.0153 & 0.00515 & 0.0174 \\
DBT & <0.00008 & <0.00008 & <0.00008 & <0.00008 & 0.00013 & 0.000108 & 0.000132 & 0.000272 \\
TBT & 0.000108 & 0.0000936 & <0.00008 & 0.0000944 & 0.000103 & 0.000109 & 0.000089 & 0.000144 \\
TTBT & <0.00008 & <0.00008 & <0.00008 & <0.00008 & <0.00008 & <0.00008 & <0.00008 & <0.00008 \\
MOT & <0.00008 & <0.00008 & <0.00008 & <0.00008 & <0.00008 & <0.00008 & <0.00008 & <0.00008 \\
DOT & <0.00008 & <0.00008 & <0.00008 & <0.00008 & <0.00008 & <0.00008 & <0.00008 & <0.00008 \\
TCyT & <0.00008 & <0.00008 & <0.00008 & <0.00008 & <0.00008 & <0.00008 & <0.00008 & <0.00008 \\
MPhT & <0.00008 & <0.00008 & <0.00008 & <0.00008 & <0.00008 & <0.00008 & <0.00008 & <0.00008 \\

DPhT & <0.00008 & <0.00008 & <0.00008 & <0.00008 & <0.00008 & <0.00008 & <0.00008 & <0.00008 \\
TPhT & <0.00008 & <0.00008 & <0.00008 & <0.00008 & <0.00008 & <0.00008 & <0.00008 & <0.00008 \\
PCB 28 & <0.000084 & <0.000084 & <0.000084 & <0.000084 & <0.000084 & <0.000084 & <0.000084 & <0.000084  \\

PCB 52 & <0.000084 & <0.000084 & <0.000084 & <0.000084 & <0.000084 & <0.000084 & <0.000084 & <0.000084  \\
PCB 101 & <0.000084 & <0.000084 & <0.000084 & <0.000084 & <0.000084 & <0.000084 & <0.000084 & <0.000084  \\
PCB 118 & <0.000084 & <0.000084 & <0.000084 & <0.000084 & <0.000084 & <0.000084 & <0.000084 & <0.000084  \\
PCB 138 & <0.000092 & <0.000092 & <0.000092 & <0.000092 & <0.000092 & <0.000092 & <0.000092 & <0.000092 \\
PCB 153 & <0.000084 & <0.000084 & <0.000084 & <0.000084 & <0.000084 & <0.000084 & <0.000084 & <0.000084  \\
PCB 180 & <0.000084 & <0.000084 & <0.000084 & <0.000084 & <0.000084 & <0.000084 & <0.000084 & <0.000084  \\
Total PCB 7 & <0.0003 & <0.0003 & <0.0003 & <0.0003 & <0.0003 & <0.0003 & <0.0003 & <0.0003  \\
\bottomrule
\end{tabularx}}
\end{adjustwidth}
\end{table}
%\end{center}
%%\end{footnotesize}
%\end{scriptsize}
%---------------------------------------------------------<
\vspace{-16pt}
\subsection[\appendixname~\thesubsection]{Cumulative Surface Leaching of Contaminants}

Notations for Table \ref{tabA04}.  The applied methodology is based on the three approaches: 

\begin{itemize}
    \item Standards SS-EN ISO 10523:2012 with measurement uncertainty of $\pm$0.05 pH-enh;
    \item SS-EN 27888-1:1994 with measurement uncertainty $\pm$2.9\%; 
    \item SGI-method with measurement uncertainty pf $\pm$ 25\%.
\end{itemize}

\vspace{-6pt}
%--------------------------------------------------------->
%\begin{scriptsize}
%\begin{center}
%\setlength\LTleft{0pt}
%\setlength\LTright{0pt}
%\begin{longtable}{|p{2.6cm}|cccccccc|}
%\begin{longtable}{@{\extracolsep{\fill}}|p{2.6cm}|cccccccc|@{}}
%\begin{longtable}{@{\extracolsep{\fill}}|p{1.8cm}|cccccccc|@{}}
\begin{table}[H] 
%\begin{threeparttable}
%\centering
\caption{Cumulative surface leaching (heavy metals, organic compounds, PCB and PAH, mg/m\textsuperscript{2}) from soil specimens during period of 64 days of sampling (2325--2332).} \label{tabA04}

\begin{adjustwidth}{-\extralength}{0cm}
%\centering %% If there is a figure in wide page, please release command \centering
\setlength{\cellWidtha}{\fulllength/9-2\tabcolsep+0.2in}
\setlength{\cellWidthb}{\fulllength/9-2\tabcolsep-.10in}
\setlength{\cellWidthc}{\fulllength/9-2\tabcolsep-0.1in}
\setlength{\cellWidthd}{\fulllength/9-2\tabcolsep-0in}
\setlength{\cellWidthe}{\fulllength/9-2\tabcolsep+0in}
\setlength{\cellWidthf}{\fulllength/9-2\tabcolsep-0in}
\setlength{\cellWidthg}{\fulllength/9-2\tabcolsep-0in}
\setlength{\cellWidthh}{\fulllength/9-2\tabcolsep-0in}
\setlength{\cellWidthi}{\fulllength/9-2\tabcolsep+0in}
\scalebox{1}[1]{\begin{tabularx}{\fulllength}{>{\raggedright\arraybackslash}m{\cellWidtha}>{\centering\arraybackslash}m{\cellWidthb}>{\centering\arraybackslash}m{\cellWidthc}>{\centering\arraybackslash}m{\cellWidthd}>{\centering\arraybackslash}m{\cellWidthe}>{\centering\arraybackslash}m{\cellWidthf}>{\centering\arraybackslash}m{\cellWidthg}>{\centering\arraybackslash}m{\cellWidthh}>{\centering\arraybackslash}m{\cellWidthi}}
\toprule \multicolumn{1}{c}{\textbf{Period}} & \multicolumn{1}{c}{\textbf{6 h}} & \multicolumn{1}{c}{\textbf{1 Days}} & \multicolumn{1}{c}{\textbf{2.25 Days}} & \multicolumn{1}{c}{\textbf{4 Days}} & \multicolumn{1}{c}{\textbf{9 Days}} & \multicolumn{1}{c}{\textbf{16 Days}} & \multicolumn{1}{c}{\textbf{36 Days}} & \multicolumn{1}{c}{\textbf{64 Days}} \\ \midrule 
%\endfirsthead
%
%\multicolumn{9}{c}%
%{{\bfseries \tablename\ \thetable{} -- continued from previous page}} \\
%\hline \multicolumn{1}{c}{\textbf{Period}} & \multicolumn{1}{c}{\textbf{6 h}} & \multicolumn{1}{c}{\textbf{1 days}} & \multicolumn{1}{c}{\textbf{2.25 days}} & \multicolumn{1}{c}{\textbf{4 days}} & \multicolumn{1}{c}{\textbf{9 days}} & \multicolumn{1}{c}{\textbf{16 days}} & \multicolumn{1}{c}{\textbf{36 days}} & \multicolumn{1}{c}{\textbf{64 days}} \\ \hline 
%\endhead
%
%\hline \multicolumn{9}{|r|}{{Continued on next page}} \\ \hline
%\endfoot
%
%\hline \hline
%\endlastfoot

Sample No & 2325 & 2326 & 2327 & 2328 & 2329 & 2330 & 2331 & 2332 \\
Al & 1.86 & 3.6 & 5 & 6.32 & 7 & 7.6 & 8.02 & 8.41 \\
As & 0.0581 & 0.133 & 0.219 & 0.284 & 0.362 & 0.444 & 0.501 & 0.584  \\
Ba & 6.04 & 12.9 & 20.6 & 27.2 & 40.3 & 53.8 & 74 & 94.6  \\
Ca & 15,900 & 33,300 & 52,900 & 72,000 & 96,300 & 121,000 & 15,300 & 190,000 \\
Cd & 0.0042 & 0.00927 & <0.01 & <0.02 & <0.02 & <0.02 & <0.03 & 0.0328 \\
Co & 0.037 & 0.0663 & 0.0902 & 0.117 & 0.159 & 0.186 & 0.215 & 0.256 \\
Cr & <0.04 & <0.08 & <0.1 & <0.2 & <0.2 & <0.2 & <0.3 & <0.3  \\
Cu & 1.3 & 2.23 & 2.91 & 3.45 & 4.04 & 4.64 & 5.53 & 6.5 \\
Fe & 0.859 & 1.52 & 2.1 & 2.57 & <3 & <3 & <3 & <4 \\
Hg & <0.002 & <0.003 & <0.005 & <0.006 & <0.008 & <0.009 & <0.01 & <0.01 \\
K & 16,000 & 31,600 & 49,100 & 65,400 & 81,300 & 96,700 & 113,000 & 128,000 \\
Mg & 39,400 & 78,800 & 122,000 & 163,000 & 197,000 & 231,000 & 258,000 & 285,000 \\
Mn & 1.33 & 2.35 & 2.9 & 3.46 & 3.65 & 3.92 & 4 & 4.05  \\
Mo & 0.782 & 1.5 & 2.29 & 3 & 3.87 & 4.66 & 5.7 & 6.9 \\
Na & 380,000 & 747,000 & 1,170,000 & 1,560,000 & 1,930,000 & 2,290,000 & 2,600,000 & 3,010,000 \\
Ni & 0.411 & 0.726 & 0.95 & 1.25 & 1.66 & 1.9 & 2.07 & 2.23 \\
Pb & <0.02 & <0.03 & <0.05 & <0.06 & <0.08 & <0.09 & <0.1 & <0.1 \\
V & 0.141 & 0.35 & 0.591 & 0.833 & 1.18 & 1.4 & 1.57 & 1.69 \\
Zn & 0.278 & <0.4 & <0.6 & <0.7 & <0.9 & <1 & <1 & <1  \\
B & 144 & 288 & 445 & 594 & 722 & 846 & 947 & 1050 \\
Sb & 0.044 & 0.0924 & 0.143 & 0.194 & 0.282 & 0.354 & 0.469 & 0.579  \\
Se & <0.2 & <0.5 & <0.7 & <0.9 & <1 & 5.35 & <6 & <6  \\
Naphthalene & <0.002 & <0.005 & <0.007 & <0.009 & <0.01 & <0.01 & <0.02 & <0.02 \\
Acy & <0.0008 & <0.002 & <0.002 & <0.003 & <0.004 & <0.005 & <0.005 & <0.006 \\
Acenaphthene & <0.0008 & <0.002 & <0.002 & <0.003 & 0.004 & 0.005 & <0.006 & <0.006 \\
Fluorine & <0.0008 & <0.002 & <0.002 & <0.003 & 0.004 & 0.005 & <0.005 & <0.006 \\
Phenanthrene & <0.0008 & <0.002 & <0.002 & <0.003 & 0.004 & 0.005 & <0.005 & <0.006 \\

Anthracene & <0.0008 & <0.002 & <0.002 & <0.003 & 0.004 & 0.005 & <0.005 & <0.006 \\
Fluoranthene & <0.0008 & <0.002 & <0.002 & <0.003 & 0.0043 & 0.0055 & <0.0074 & <0.0084 \\
Pyrene & <0.0008 & 0.0019 & 0.0036 & 0.0056 & 0.0091 & 0.013 & 0.018 & 0.023 \\
B[a]A & <0.0008 & <0.002 & <0.002 & <0.003 & <0.004 & <0.005 & <0.005 & <0.006 \\

Chrysene & <0.0008 & <0.002 & <0.002 & <0.003 & <0.004 & <0.005 & <0.005 & <0.006 \\
BjF & <0.0008 & <0.002 & <0.002 & <0.003 & <0.004 & <0.005 & <0.005 & <0.006 \\

BkF & <0.0008 & <0.002 & <0.002 & <0.003 & <0.004 & <0.005 & <0.005 & <0.006 \\
BaP & <0.0008 & <0.002 & <0.002 & <0.003 & <0.004 & <0.005 & <0.005 & <0.006 \\

DBahA & <0.0008 & <0.002 & <0.002 & <0.003 & <0.004 & <0.005 & <0.005 & <0.006 \\
BghiP & <0.0008 & <0.002 & <0.002 & <0.003 & <0.004 & <0.005 & <0.005 & <0.006 \\
IP & <0.0008 & <0.002 & <0.002 & <0.003 & <0.004 & <0.005 & <0.005 & <0.006 \\
Total PAH 16 & <0.0073 & 0.0084 & 0.01 & 0.012 & 0.018 & 0.023 & 0.031 & 0.036 \\
Total carc. PAH & <0.0027 & <0.0054 & <0.0081 & <0.011 & <0.013 & <0.016 & <0.019 & <0.021 \\
Total oth. PAH & <0.005 & 0.0058 & 0.0074 & 0.0094 & 0.015 & 0.021 & 0.028 & 0.033 \\
Total PAH L & <0.0019 & <0.0038 & <0.0058 & <0.0077 & 0.0086 & 0.009 & <0.011 & <0.013 \\
Total PAH M & <0.002 & 0.0035 & 0.0051 & 0.0071 & 0.012 & 0.017 & 0.024 & 0.029 \\
Total PAH H & <0.003 & <0.006 & <0.009 & <0.01 & <0.02 & <0.02 & <0.02 & <0.02 \\

MBT & 0.000333 & 0.003 & 0.00531 & 0.00866 & 0.0126 & 0.0279 & 0.0331 & 0.0505 \\
DBT & <0.00008 & <0.0002 & <0.0002 & <0.0003 & 0.000436 & 0.000545 & 0.000677 & 0.000948 \\
TBT & 0.000108 & 0.000202 & <0.0003 & 0.000373 & 0.000476 & 0.000585 & 0.000674 & 0.000818 \\
TTBT & <0.00008 & <0.0002 & <0.0002 & <0.0003 & <0.0004 & <0.0005 & <0.0005 & <0.0006 \\
MOT & <0.00008 & <0.0002 & <0.0002 & <0.0003 & <0.0004 & <0.0005 & <0.0005 & <0.0006 \\
DOT & <0.00008 & <0.0002 & <0.0002 & <0.0003 & <0.0004 & <0.0005 & <0.0005 & <0.0006 \\
TCyT & <0.00008 & <0.0002 & <0.0002 & <0.0003 & <0.0004 & <0.0005 & <0.0005 & <0.0006 \\
MPhT & <0.00008 & <0.0002 & <0.0002 & <0.0003 & <0.0004 & <0.0005 & <0.0005 & <0.0006 \\
DPhT & <0.00008 & <0.0002 & <0.0002 & <0.0003 & <0.0004 & <0.0005 & <0.0005 & <0.0006 \\
TPhT & <0.00008 & <0.0002 & <0.0002 & <0.0003 & <0.0004 & <0.0005 & <0.0005 & <0.0006 \\
\bottomrule
\end{tabularx}}
\end{adjustwidth}
\end{table}

\begin{table}[H]\ContinuedFloat
\tablesize{\small}
\caption{{\em Cont.}}
 \label{tabA04}


\begin{adjustwidth}{-\extralength}{0cm}
%\centering %% If there is a figure in wide page, please release command \centering
\setlength{\cellWidtha}{\fulllength/9-2\tabcolsep+0.2in}
\setlength{\cellWidthb}{\fulllength/9-2\tabcolsep-.10in}
\setlength{\cellWidthc}{\fulllength/9-2\tabcolsep-0.1in}
\setlength{\cellWidthd}{\fulllength/9-2\tabcolsep-0in}
\setlength{\cellWidthe}{\fulllength/9-2\tabcolsep+0in}
\setlength{\cellWidthf}{\fulllength/9-2\tabcolsep-0in}
\setlength{\cellWidthg}{\fulllength/9-2\tabcolsep-0in}
\setlength{\cellWidthh}{\fulllength/9-2\tabcolsep-0in}
\setlength{\cellWidthi}{\fulllength/9-2\tabcolsep+0in}
\scalebox{1}[1]{\begin{tabularx}{\fulllength}{>{\raggedright\arraybackslash}m{\cellWidtha}>{\centering\arraybackslash}m{\cellWidthb}>{\centering\arraybackslash}m{\cellWidthc}>{\centering\arraybackslash}m{\cellWidthd}>{\centering\arraybackslash}m{\cellWidthe}>{\centering\arraybackslash}m{\cellWidthf}>{\centering\arraybackslash}m{\cellWidthg}>{\centering\arraybackslash}m{\cellWidthh}>{\centering\arraybackslash}m{\cellWidthi}}
\toprule \multicolumn{1}{c}{\textbf{Period}} & \multicolumn{1}{c}{\textbf{6 h}} & \multicolumn{1}{c}{\textbf{1 Days}} & \multicolumn{1}{c}{\textbf{2.25 Days}} & \multicolumn{1}{c}{\textbf{4 Days}} & \multicolumn{1}{c}{\textbf{9 Days}} & \multicolumn{1}{c}{\textbf{16 Days}} & \multicolumn{1}{c}{\textbf{36 Days}} & \multicolumn{1}{c}{\textbf{64 Days}} \\ \midrule 
PCB 28 & <0.000084 & <0.00017 & <0.00025 & <0.00034 & <0.00042 & <0.00051 & <0.00059 & <0.00068 \\
PCB 52 & <0.000084 & <0.00017 & <0.00025 & <0.00034 & <0.00042 & <0.00051 & <0.00059 & <0.00068 \\
PCB 101 & <0.000084 & <0.00017 & <0.00025 & <0.00034 & <0.00042 & <0.00051 & <0.00059 & <0.00068 \\
PCB 118 & <0.000084 & <0.00017 & <0.00025 & <0.00034 & <0.00042 & <0.00051 & <0.00059 & <0.00068 \\
PCB 138 & <0.000092 & <0.00018 & <0.00028 & <0.00037 & <0.00046 & <0.00055 & <0.00064 & <0.00074 \\
PCB 153 & <0.000084 & <0.00017 & <0.00025 & <0.00034 & <0.00042 & <0.00051 & <0.00059 & <0.00068 \\
PCB 180  & <0.000084 & <0.00017 & <0.00025 & <0.00034 & <0.00042 & <0.00051 & <0.00059 & <0.00068 \\
Total PCB 7 & <0.0003 & <0.0006 & <0.0009 & <0.0012 & <0.0015 & <0.0018 & <0.0021 & <0.0024 \\
\bottomrule
\end{tabularx}}
\end{adjustwidth}
\end{table}
%\end{center}
%%\end{footnotesize}
%\end{scriptsize}
%---------------------------------------------------------<
\vspace{-16pt}
\section[\appendixname~\thesection]{Leaching of Contaminants: Sampling 2333--2340}\label{AppC1}
\subsection[\appendixname~\thesubsection]{Surface Leaching of Contaminants}

Notations for all the samples in Table \ref{tabA05}. The applied Methodology contains Batch 1, Prov 8A. The start and end days of the tests are from 30 January 2023 until 4 April 2023. The experiments were performed by Robert Selegård in Swedish Geotechnical Institute, registered under Diary Nr. 1.1-2107-0587. The amount of leachate for all samples is 1.5 L. Conductivity was measured in mS/m by temperature of sampling at 25~$^\circ$C. The period refers to the day of sampling. Surface: 0.0197 m\textsuperscript{2}. Redox (mVolts) refers to oxidation (reduction) potential of solutions or chemical species to acquire / lose electrons and be reduced/oxidised. Leached amount/withdrawal is given in mg/m\textsuperscript{2} for all the samples.

%--------------------------------------------------------->
%\begin{scriptsize}
%\begin{center}
%\setlength\LTleft{0pt}
%\setlength\LTright{0pt}
%\begin{longtable}{|p{2.6cm}|cccccccc|}
%\begin{longtable}{@{\extracolsep{\fill}}|p{2.6cm}|cccccccc|@{}}
%\begin{longtable}{@{\extracolsep{\fill}}|p{1.8cm}|cccccccc|@{}}
\vspace{-2pt}
\begin{table}[H] 
%\begin{threeparttable}
%\centering
\caption{Surface leaching of contaminants (heavy metals, organic compounds, PCB and PAH, mg/m\textsuperscript{2}) from soil specimens during period of 64 days of sampling (2333--2340).} \label{tabA05}
\begin{adjustwidth}{-\extralength}{0cm}
%\centering %% If there is a figure in wide page, please release command \centering
\setlength{\cellWidtha}{\fulllength/9-2\tabcolsep+0.2in}
\setlength{\cellWidthb}{\fulllength/9-2\tabcolsep-.10in}
\setlength{\cellWidthc}{\fulllength/9-2\tabcolsep-0.1in}
\setlength{\cellWidthd}{\fulllength/9-2\tabcolsep-0in}
\setlength{\cellWidthe}{\fulllength/9-2\tabcolsep+0in}
\setlength{\cellWidthf}{\fulllength/9-2\tabcolsep-0in}
\setlength{\cellWidthg}{\fulllength/9-2\tabcolsep-0in}
\setlength{\cellWidthh}{\fulllength/9-2\tabcolsep-0in}
\setlength{\cellWidthi}{\fulllength/9-2\tabcolsep+0in}
\scalebox{1}[1]{\begin{tabularx}{\fulllength}{>{\raggedright\arraybackslash}m{\cellWidtha}>{\centering\arraybackslash}m{\cellWidthb}>{\centering\arraybackslash}m{\cellWidthc}>{\centering\arraybackslash}m{\cellWidthd}>{\centering\arraybackslash}m{\cellWidthe}>{\centering\arraybackslash}m{\cellWidthf}>{\centering\arraybackslash}m{\cellWidthg}>{\centering\arraybackslash}m{\cellWidthh}>{\centering\arraybackslash}m{\cellWidthi}}
\toprule \multicolumn{1}{c}{\textbf{Period}} & \multicolumn{1}{c}{\textbf{6 h}} & \multicolumn{1}{c}{\textbf{1 Days}} & \multicolumn{1}{c}{\textbf{2.25 Days}} & \multicolumn{1}{c}{\textbf{4 Days}} & \multicolumn{1}{c}{\textbf{9 Days}} & \multicolumn{1}{c}{\textbf{16 Days}} & \multicolumn{1}{c}{\textbf{36 Days}} & \multicolumn{1}{c}{\textbf{64 Days}} \\ \midrule 
%\endfirsthead
%
%\multicolumn{9}{c}%
%{{\bfseries \tablename\ \thetable{} -- continued from previous page}} \\
%\hline \multicolumn{1}{c}{\textbf{Period}} & \multicolumn{1}{c}{\textbf{6 h}} & \multicolumn{1}{c}{\textbf{1 days}} & \multicolumn{1}{c}{\textbf{2.25 days}} & \multicolumn{1}{c}{\textbf{4 days}} & \multicolumn{1}{c}{\textbf{9 days}} & \multicolumn{1}{c}{\textbf{16 days}} & \multicolumn{1}{c}{\textbf{36 days}} & \multicolumn{1}{c}{\textbf{64 days}} \\ \hline 
%\endhead
%
%\hline \multicolumn{9}{|r|}{{Continued on next page}} \\ \hline
%\endfoot
%
%\hline \hline
%\endlastfoot

Sample No & 2333 & 2334 & 2335 & 2336 & 2337 & 2338 & 2339 & 2340 \\
pH & 8.75 & 8.68 & 8.57 & 8.72 & 8.8 & 8.49 & 8.15 & 7.87  \\
Conductivity & 2520 & 2510 & 2660 & 2860 & 2750 & 2760 & 2740 & 2760 \\
Redox (mV) & 46 & 84 & 74 & 57 & 45 & 42 & 74 & 75 \\
Al & 2.25 & 2.03 & 1.75 & 1.6 & 0.749 & 0.78 & 0.624 & 0.646 \\
As & 0.0637 & 0.074 & 0.084 & 0.0687 & 0.105 & 0.135 & 0.071 & 0.0895 \\
Ba & 8.11 & 9.25 & 11.3 & 10.4 & 19.4 & 20.5 & 30.3 & 26.8 \\

Ca & 16,700 & 18,100 & 22,000 & 22,000 & 30,100 & 31,500 & 50,000 & 47,000 \\

Cd & 0.0056 & 0.00637 & <0.004 & <0.004 & <0.004 & <0.004 & <0.004 & 0.00499  \\

Co & 0.0414 & 0.036 & 0.033 & 0.0226 & 0.042 & 0.028 & 0.0486 & 0.0499 \\
Cr & 0.0414 & 0.036 & 0.033 & 0.0226 & 0.042 & 0.028 & 0.0486 & 0.0499 \\
Cu & 1.46 & 1.13 & 0.803 & 0.722 & 0.739 & 0.741 & 0.941 &  0.941 \\
Fe & 0.902 & 0.84 & 0.556 & 0.593 & <0.3 & <0.3 & <0.3 & <0.3 \\
Hg & <0.002 & <0.002 & <0.002 & <0.002 & <0.002 & <0.002 & <0.002 & <0.002 \\
K & 16,000 & 15,800 & 17,500 & 16,200 & 15,900 & 15,600 & 16,400 &  20,000 \\
Mg & 38,800 & 37,900 & 41,800 & 38,400 & 29,500 & 29,900 & 20,000 & 20,500 \\
Mn & 1.28 & 1.1 & 0.612 & 0.772 & 0.135 & 0.0856 & 0.0263 & 0.0247 \\
Mo & 0.811 & 0.772 & 0.895 & 0.772 & 1.04 & 0.948 & 1.19 & 1.29 \\
Na & 375,000 & 370,000 & 415,000 & 389,000 & 370,000 & 366,000 & 360,000 & 364,000 \\
Ni & 0.372 & 0.362 & 0.292 & 0.268 & 0.4 & 0.259 & 0.338 & 0.129 \\
Pb & <0.02 & <0.02 & <0.02 & <0.02 & <0.02 & <0.02 & <0.02 & <0.02 \\
V & 0.172 & 0.284 & 0.353 & 0.364 & 0.45 & 0.295 & 0.154 &  0.122 \\
Zn & 0.294 & 0.183 & 0.216 & 0.2 & <0.2 & <0.2 & <0.2 & <0.2 \\
B & 139 & 139 & 151 & 141 & 116 & 110 & 88.7 & 91.8  \\
Sb & 0.0561 & 0.0664 & 0.0699 & 0.0724 & 0.108 & 0.0971 & 0.135 & 0.135 \\
Se & <0.2 & <0.2 & <0.2 & <0.2 & 0.55 & 4.47 & <0.2 & <0.2 \\

Naphthalene & <0.002 & <0.002 & <0.002 & <0.002 & 0.004 & <0.002 & <0.002 & <0.002 \\
Acy & <0.0008 & <0.0008 & <0.0008 & <0.0008 & <0.0008 & <0.0008 & <0.0008 & <0.0008 \\
Acenaphthene & <0.0008 & <0.0008 & <0.0008 & <0.0008 & 0.0017 & <0.0008 & <0.0008 & <0.0008  \\
Fluorine & <0.0008 & <0.0008 & <0.0008 & <0.0008 & 0.00099 & <0.0008 & <0.0008 & <0.0008 \\
Phenanthrene & <0.002 & <0.002 & <0.002 & <0.002 & 0.0041 & <0.002 & <0.002 & <0.002 \\
\bottomrule
\end{tabularx}}
\end{adjustwidth}
\end{table}

\begin{table}[H]\ContinuedFloat
\tablesize{\small}
\caption{{\em Cont.}}
 \label{tabA05}
\begin{adjustwidth}{-\extralength}{0cm}
%\centering %% If there is a figure in wide page, please release command \centering
\setlength{\cellWidtha}{\fulllength/9-2\tabcolsep+0.2in}
\setlength{\cellWidthb}{\fulllength/9-2\tabcolsep-.10in}
\setlength{\cellWidthc}{\fulllength/9-2\tabcolsep-0.1in}
\setlength{\cellWidthd}{\fulllength/9-2\tabcolsep-0in}
\setlength{\cellWidthe}{\fulllength/9-2\tabcolsep+0in}
\setlength{\cellWidthf}{\fulllength/9-2\tabcolsep-0in}
\setlength{\cellWidthg}{\fulllength/9-2\tabcolsep-0in}
\setlength{\cellWidthh}{\fulllength/9-2\tabcolsep-0in}
\setlength{\cellWidthi}{\fulllength/9-2\tabcolsep+0in}
\scalebox{1}[1]{\begin{tabularx}{\fulllength}{>{\raggedright\arraybackslash}m{\cellWidtha}>{\centering\arraybackslash}m{\cellWidthb}>{\centering\arraybackslash}m{\cellWidthc}>{\centering\arraybackslash}m{\cellWidthd}>{\centering\arraybackslash}m{\cellWidthe}>{\centering\arraybackslash}m{\cellWidthf}>{\centering\arraybackslash}m{\cellWidthg}>{\centering\arraybackslash}m{\cellWidthh}>{\centering\arraybackslash}m{\cellWidthi}}
\toprule \multicolumn{1}{c}{\textbf{Period}} & \multicolumn{1}{c}{\textbf{6 h}} & \multicolumn{1}{c}{\textbf{1 Days}} & \multicolumn{1}{c}{\textbf{2.25 Days}} & \multicolumn{1}{c}{\textbf{4 Days}} & \multicolumn{1}{c}{\textbf{9 Days}} & \multicolumn{1}{c}{\textbf{16 Days}} & \multicolumn{1}{c}{\textbf{36 Days}} & \multicolumn{1}{c}{\textbf{64 Days}} \\ \midrule 
Anthracene & <0.0008 & <0.0008 & <0.0008 & <0.0008 & <0.0008 & <0.0008 & <0.0008 & <0.0008 \\
Fluoranthene & 0.00099 & 0.002 & 0.0022 & 0.0028 & 0.0043 & 0.0036 & <0.0008 & <0.0008 \\
Pyrene & 0.0011 & 0.0022 & 0.0029 & 0.0034 & 0.0057 & 0.0057 & 0.0037 & 0.0013 \\
B[a]A & <0.0008 & <0.0008 & <0.0008 & <0.0008 & <0.0008 & <0.0008 & <0.0008 & <0.0008 \\
Chrysene & <0.0008 & <0.0008 & <0.0008 & <0.0008 & <0.0008 & <0.0008 & <0.0008 & <0.0008 \\
BjF & <0.0008 & <0.0008 & <0.0008 & <0.0008 & <0.0008 & <0.0008 & <0.0008 & <0.0008 \\
BkF & <0.0008 & <0.0008 & <0.0008 & <0.0008 & <0.0008 & <0.0008 & <0.0008 & <0.0008 \\

BaP & <0.0008 & <0.0008 & <0.0008 & <0.0008 & <0.0008 & <0.0008 & <0.0008 & <0.0008 \\
DBahA & <0.0008 & <0.0008 & <0.0008 & <0.0008 & <0.0008 & <0.0008 & <0.0008 & <0.0008 \\
BghiP & <0.0008 & <0.0008 & <0.0008 & <0.0008 & <0.0008 & <0.0008 & <0.0008 & <0.0008 \\
IP & <0.0008 & <0.0008 & <0.0008 & <0.0008 & <0.0008 & <0.0008 & <0.0008 & <0.0008 \\
Total PAH 16 & 0.0021 & 0.0042 & 0.0051 & 0.0062 & 0.021 & 0.00933 & 0.0037 & 0.0013 \\
Total carc. PAH & <0.0027 & <0.0027 & <0.0027 & <0.0027 & <0.0027 & <0.0027 & <0.0027 & <0.0027 \\
Total oth. PAH & 0.0021 & 0.0042 & 0.0051 & 0.0062 & 0.021 & 0.00933 &0.0037 & 0.0013  \\
Total PAH L & <0.0019 & <0.0019 & <0.0019 & <0.0019 & 0.0057 & <0.0019 & <0.0019 & <0.0019 \\
Total PAH M & 0.0021 & 0.0042 & 0.0051 & 0.0062 & 0.015 & 0.00933 & 0.0037 & 0.0013 \\
Total PAH H & <0.003 & <0.003 & <0.003 & <0.003 & <0.003 & <0.003 & <0.003 & <0.003  \\
MBT & 0.00142 & 0.00255 & 0.00189 & 0.00262 & 0.00453 & 0.00902 & 0.00554 & 0.00661 \\
DBT& 0.000091 & 0.00011 & 0.0000849 & 0.00011 & 0.000211 & 0.000248 & 0.000236 & 0.000593 \\
TBT & 0.000142 & 0.000118 & 0.000109 & 0.000151 & 0.000157 & 0.000121 & 0.000125 & 0.000113 \\
TTBT & <0.00008 & <0.00008 & <0.00008 & <0.00008 & <0.00008 & <0.00008 & <0.00008 & <0.00008  \\
MOT & <0.00008 & <0.00008 & <0.00008 & <0.00008 & <0.00008 & <0.00008 & <0.00008 & <0.00008  \\
DOT & <0.00008 & <0.00008 & <0.00008 & <0.00008 & <0.00008 & <0.00008 & <0.00008 & <0.00008  \\
TCyT & <0.00008 & <0.00008 & <0.00008 & <0.00008 & <0.00008 & <0.00008 & <0.00008 & <0.00008  \\
MPhT & <0.00008 & <0.00008 & <0.00008 & <0.00008 & <0.00008 & <0.00008 & <0.00008 & <0.00008  \\
DPhT & <0.00008 & <0.00008 & <0.00008 & <0.00008 & <0.00008 & <0.00008 & <0.00008 & <0.00008  \\
TPhT & <0.00008 & <0.00008 & <0.00008 & <0.00008 & <0.00008 & <0.00008 & <0.00008 & <0.00008  \\

PCB 28 & <0.000084 & <0.000084 & <0.000084 & <0.000084 & <0.000084 & <0.000084 & <0.000084 & <0.000084 \\
PCB 52 & <0.000084 & <0.000084 & <0.000084 & <0.000084 & <0.000084 & <0.000084 & <0.000084 & <0.000084 \\
PCB 101 & <0.000084 & <0.000084 & <0.000084 & <0.000084 & <0.000084 & <0.000084 & <0.000084 & <0.000084 \\
PCB 118 & <0.000084 & <0.000084 & <0.000084 & <0.000084 & <0.000084 & <0.000084 & <0.000084 & <0.000084 \\
PCB 138 & <0.000092 & <0.000092 & <0.000092 & <0.000092 & <0.000092 & <0.000092 & <0.000092 & <0.000092 \\
PCB 153 & <0.000084 & <0.000084 & <0.000084 & <0.000084 & <0.000084 & <0.000084 & <0.000084 & <0.000084 \\
PCB 180 & <0.000084 & <0.000084 & <0.000084 & <0.000084 & <0.000084 & <0.000084 & <0.000084 & <0.000084 \\
Total PCB 7 & <0.0003 & <0.0003 & <0.0003 & <0.0003 & <0.0003 & <0.0003 & <0.0003 & <0.0003 \\
\bottomrule
\end{tabularx}}
\end{adjustwidth}
\end{table}
%\end{center}
%%\end{footnotesize}
%\end{scriptsize}
%---------------------------------------------------------<
\vspace{-14pt}
\subsection[\appendixname~\thesubsection]{Cumulative Surface Leaching of Contaminants}
\vspace{-6pt}
%--------------------------------------------------------->
%\begin{scriptsize}
%\begin{center}
%\setlength\LTleft{0pt}
%\setlength\LTright{0pt}
%\begin{longtable}{@{\extracolsep{\fill}}|p{1.8cm}|cccccccc|@{}}
\begin{table}[H] 
%\begin{threeparttable}
%\centering
\caption{Cumulative surface leaching (heavy metals, organic compounds, PCB and PAH, mg/m\textsuperscript{2}) from soil specimens during period of 64 days of sampling (2333--2340).} \label{tabA06}
\begin{adjustwidth}{-\extralength}{0cm}
%\centering %% If there is a figure in wide page, please release command \centering
\setlength{\cellWidtha}{\fulllength/9-2\tabcolsep+0.2in}
\setlength{\cellWidthb}{\fulllength/9-2\tabcolsep-.10in}
\setlength{\cellWidthc}{\fulllength/9-2\tabcolsep-0.1in}
\setlength{\cellWidthd}{\fulllength/9-2\tabcolsep-0in}
\setlength{\cellWidthe}{\fulllength/9-2\tabcolsep+0in}
\setlength{\cellWidthf}{\fulllength/9-2\tabcolsep-0in}
\setlength{\cellWidthg}{\fulllength/9-2\tabcolsep-0in}
\setlength{\cellWidthh}{\fulllength/9-2\tabcolsep-0in}
\setlength{\cellWidthi}{\fulllength/9-2\tabcolsep+0in}
\scalebox{1}[1]{\begin{tabularx}{\fulllength}{>{\raggedright\arraybackslash}m{\cellWidtha}>{\centering\arraybackslash}m{\cellWidthb}>{\centering\arraybackslash}m{\cellWidthc}>{\centering\arraybackslash}m{\cellWidthd}>{\centering\arraybackslash}m{\cellWidthe}>{\centering\arraybackslash}m{\cellWidthf}>{\centering\arraybackslash}m{\cellWidthg}>{\centering\arraybackslash}m{\cellWidthh}>{\centering\arraybackslash}m{\cellWidthi}}
\toprule \multicolumn{1}{c}{\textbf{Period}} & \multicolumn{1}{c}{\textbf{6 h}} & \multicolumn{1}{c}{\textbf{1 Days}} & \multicolumn{1}{c}{\textbf{2.25 Days}} & \multicolumn{1}{c}{\textbf{4 Days}} & \multicolumn{1}{c}{\textbf{9 Days}} & \multicolumn{1}{c}{\textbf{16 Days}} & \multicolumn{1}{c}{\textbf{36 Days}} & \multicolumn{1}{c}{\textbf{64 Days}} \\ \midrule 
%\endfirsthead
%
%\multicolumn{9}{c}%
%{{\bfseries \tablename\ \thetable{} -- continued from previous page}} \\
%\hline \multicolumn{1}{c}{\textbf{Period}} & \multicolumn{1}{c}{\textbf{6 h}} & \multicolumn{1}{c}{\textbf{1 days}} & \multicolumn{1}{c}{\textbf{2.25 days}} & \multicolumn{1}{c}{\textbf{4 days}} & \multicolumn{1}{c}{\textbf{9 days}} & \multicolumn{1}{c}{\textbf{16 days}} & \multicolumn{1}{c}{\textbf{36 days}} & \multicolumn{1}{c}{\textbf{64 days}} \\ \hline 
%\endhead
%
%\hline \multicolumn{9}{|r|}{{Continued on next page}} \\ \hline
%\endfoot
%
%\hline \hline
%\endlastfoot

Sample No & 2333 & 2334 & 2335 & 2336 & 2337 & 2338 & 2339 & 2340 \\
Al & 2.25 & 4.28 & 6.03 & 7.63 & 8.38 & 9.16 & 9.78 & 10.4 \\
As & 0.0637 & 0.138 & 0.22 & 0.29 & 0.395 & 0.53 & 0.6 & 0.69 \\
Ba & 8.11 & 17.4 & 28.7 & 39.1 & 58.5 & 79 & 109 & 136 \\
Ca & 16,700 & 34,800 & 57,000 & 79,000 & 109,000 & 141,000 & 200,000 & 230,000 \\
Cd & 0.0056 & 0.012 & <0.02 & <0.02 & <0.02 & <0.03 & <0.03 & 0.0361 \\
Co & 0.0414 & 0.0775 & 0.11 & 0.133 & 0.18 & 0.2 & 0.251 & 0.301 \\
Cr & <0.04 & <0.08 & <0.1 & <0.2 & <0.2 & <0.2 & <0.3 & <0.3 \\

Cu & 1.46 & 2.59 & 3.4 & 4.12 & 4.86 & 5.6 & 6.54 & 7.48 \\
Fe & 0.902 & 1.7 & 2.3 & 2.89 & <3 & <4 & <4 & <4 \\
Hg & <0.002 & <0.003 & <0.005 & <0.006 & <0.008 & <0.009 & <0.01 & <0.01  \\
K & 16,000 & 31,800 & 49,300 & 65,500 & 81,400 & 97,000 & 113,000 & 100,000 \\
Mg & 38,800 & 76,800 & 119,000 & 157,000 & 186,000 & 216,000 & 236,000 & 257,000 \\
Mn & 1.28 & 2.39 & 3 & 3.77 & 3.91 & 3.99 & 4.02 & 4.04 \\
\bottomrule
\end{tabularx}}
\end{adjustwidth}
\end{table}

\begin{table}[H]\ContinuedFloat
\tablesize{\small}
\caption{{\em Cont.}}
 \label{tabA06}
\begin{adjustwidth}{-\extralength}{0cm}
%\centering %% If there is a figure in wide page, please release command \centering
\setlength{\cellWidtha}{\fulllength/9-2\tabcolsep+0.2in}
\setlength{\cellWidthb}{\fulllength/9-2\tabcolsep-.10in}
\setlength{\cellWidthc}{\fulllength/9-2\tabcolsep-0.1in}
\setlength{\cellWidthd}{\fulllength/9-2\tabcolsep-0in}
\setlength{\cellWidthe}{\fulllength/9-2\tabcolsep+0in}
\setlength{\cellWidthf}{\fulllength/9-2\tabcolsep-0in}
\setlength{\cellWidthg}{\fulllength/9-2\tabcolsep-0in}
\setlength{\cellWidthh}{\fulllength/9-2\tabcolsep-0in}
\setlength{\cellWidthi}{\fulllength/9-2\tabcolsep+0in}
\scalebox{1}[1]{\begin{tabularx}{\fulllength}{>{\raggedright\arraybackslash}m{\cellWidtha}>{\centering\arraybackslash}m{\cellWidthb}>{\centering\arraybackslash}m{\cellWidthc}>{\centering\arraybackslash}m{\cellWidthd}>{\centering\arraybackslash}m{\cellWidthe}>{\centering\arraybackslash}m{\cellWidthf}>{\centering\arraybackslash}m{\cellWidthg}>{\centering\arraybackslash}m{\cellWidthh}>{\centering\arraybackslash}m{\cellWidthi}}
\toprule \multicolumn{1}{c}{\textbf{Period}} & \multicolumn{1}{c}{\textbf{6 h}} & \multicolumn{1}{c}{\textbf{1 Days}} & \multicolumn{1}{c}{\textbf{2.25 Days}} & \multicolumn{1}{c}{\textbf{4 Days}} & \multicolumn{1}{c}{\textbf{9 Days}} & \multicolumn{1}{c}{\textbf{16 Days}} & \multicolumn{1}{c}{\textbf{36 Days}} & \multicolumn{1}{c}{\textbf{64 Days}} \\ \midrule 
Mo & 0.811 & 1.58 & 2.48 & 3.25 & 4.29 & 5.24 & 6.43 & 7.72  \\
Na & 375,000 & 750,000 & 1,170,000 & 1,550,000 & 1,900,000 & 2,290,000 & 2,600,000 & 3,010,000 \\
Ni & 0.372 & 0.733 & 1.03 & 1.29 & 1.7 & 1.95 & 2.29 & 2.42  \\
Pb & <0.02 & <0.03 & <0.05 & <0.06 & <0.08 & <0.09 & <0.1 &  <0.1 \\
V & 0.172 & 0.456 & 0.81 & 1.17 & 1.6 & 1.92 & 2.07 & 2.2 \\
Zn & 0.294 & 0.477 & 0.693 & 0.893 & <1 & <1 & <1 & <2 \\
B & 139 & 278 & 430 & 570 & 686 & 796 & 885 & 976 \\
Sb & 0.0561 & 0.123 & 0.192 & 0.265 & 0.373 & 0.47 & 0.604 & 0.74 \\
Se & <0.2 & <0.5 & <0.7 & <0.9 & 1.5 & 5.94 & <6 & <6 \\
Naphthalene & <0.002 & <0.005 & <0.007 & <0.009 & 0.013 & <0.02 & <0.02 & <0.02 \\
Acy & <0.0008 & <0.002 & <0.002 & <0.003 & <0.004 & <0.005 & <0.005 & <0.006 \\
Acenaphthene & <0.0008 & <0.002 & <0.002 & <0.003 & 0.0047 & <0.006 & <0.006 & <0.007 \\
Fluorine & <0.0008 & <0.002 & <0.002 & <0.003 & <0.0041 & <0.005 & <0.006 & <0.006 \\
Phenanthrene & <0.002 & <0.003 & <0.005 & <0.006 & 0.01 & <0.01 & <0.01 & <0.01  \\
Anthracene & <0.0008 & <0.002 & <0.002 & <0.003 & <0.004 & <0.005 & <0.005 & <0.006  \\
Fluoranthene & 0.00099 & 0.003 & 0.0052 & 0.008 & 0.012 & 0.016 & <0.02 & <0.02 \\
Pyrene & 0.0011 & 0.0033 & 0.0062 & 0.0096 & 0.015 & 0.021 & 0.025 & 0.026 \\
B[a]A & <0.0008 & <0.002 & <0.002 & <0.003 & <0.004 & <0.005 & <0.005 & <0.006  \\
Chrysene & <0.0008 & <0.002 & <0.002 & <0.003 & <0.004 & <0.005 & <0.005 & <0.006  \\
BjF  & <0.0008 & <0.002 & <0.002 & <0.003 & <0.004 & <0.005 & <0.005 & <0.006  \\
BkF & <0.0008 & <0.002 & <0.002 & <0.003 & <0.004 & <0.005 & <0.005 & <0.006  \\
BaP & <0.0008 & <0.002 & <0.002 & <0.003 & <0.004 & <0.005 & <0.005 & <0.006  \\
DBahA & <0.0008 & <0.002 & <0.002 & <0.003 & <0.004 & <0.005 & <0.005 & <0.006  \\
BghiP & <0.0008 & <0.002 & <0.002 & <0.003 & <0.004 & <0.005 & <0.005 & <0.006  \\
IP & <0.0008 & <0.002 & <0.002 & <0.003 & <0.004 & <0.005 & <0.005 & <0.006  \\
Total PAH 16 & 0.0021 & 0.0063 & 0.011 & 0.018 & 0.038 & 0.0476 & 0.051 & 0.053 \\
Total carc. PAH & <0.0027 & <0.0054 & <0.008 & <0.011 & <0.013 & <0.016 & <0.019 & <0.021 \\
Total oth. PAH & 0.0021 & 0.0063 & 0.011 & 0.018 & 0.038 & 0.0476 & 0.051 & 0.053 \\
Total PAH L & <0.0019 & <0.0038 & <0.0057 & <0.0076 & 0.013 & <0.015 & <0.017 &  <0.019 \\
Total PAH M & 0.0021 & 0.0063 & 0.011 & 0.018 & 0.0326 & 0.0419 & 0.046 &  0.047 \\
Total PAH H & <0.003 & <0.006 & <0.009 & <0.01 & <0.02 & <0.02 & <0.02 & <0.02 \\
MBT & 0.00142 & 0.00398 & 0.00586 & 0.00849 & 0.013 & 0.022 & 0.0276 & 0.0342 \\
DBT & 0.000091 & 0.0002 & 0.000283 & 0.00039 & 0.000601 & 0.000849 & 0.00108 & 0.00168 \\
TBT & 0.000142 & 0.00026 & 0.000369 & 0.000519 & 0.000676 & 0.000797 & 0.000922 & 0.00104 \\
TTBT & <0.00008 & <0.0002 & <0.0002 & <0.0003 & <0.0004 & <0.0005 & <0.0005 & <0.0006 \\
MOT & <0.00008 & <0.0002 & <0.0002 & <0.0003 & <0.0004 & <0.0005 & <0.0005 & <0.0006 \\
DOT & <0.00008 & <0.0002 & <0.0002 & <0.0003 & <0.0004 & <0.0005 & <0.0005 & <0.0006 \\
TCyT & <0.00008 & <0.0002 & <0.0002 & <0.0003 & <0.0004 & <0.0005 & <0.0005 & <0.0006 \\

MPhT & <0.00008 & <0.0002 & <0.0002 & <0.0003 & <0.0004 & <0.0005 & <0.0005 & <0.0006 \\
DPhT & <0.00008 & <0.0002 & <0.0002 & <0.0003 & <0.0004 & <0.0005 & <0.0005 & <0.0006 \\
TPhT & <0.00008 & <0.0002 & <0.0002 & <0.0003 & <0.0004 & <0.0005 & <0.0005 & <0.0006 \\
PCB 28 & <0.000084 & <0.00017 & <0.00025 & <0.00034 & <0.00042 & <0.0005 & <0.00059 & <0.00067 \\
PCB 52 & <0.000084 & <0.00017 & <0.00025 & <0.00034 & <0.00042 & <0.0005 & <0.00059 & <0.00067 \\
PCB 101 & <0.000084 & <0.00017 & <0.00025 & <0.00034 & <0.00042 & <0.0005 & <0.00059 & <0.00067 \\
PCB 118 & <0.000084 & <0.00017 & <0.00025 & <0.00034 & <0.00042 & <0.0005 & <0.00059 & <0.00067 \\
PCB 138 & <0.000092 & <0.00018 & <0.00028 & <0.00037 & <0.00046 & <0.00055 & <0.00064 & <0.00073 \\
PCB 153 & <0.000084 & <0.00017 & <0.00025 & <0.00034 & <0.00042 & <0.0005 & <0.00059 & <0.00067 \\
PCB 180  & <0.000084 & <0.00017 & <0.00025 & <0.00034 & <0.00042 & <0.0005 & <0.00059 & <0.00067 \\
Total PCB 7 & <0.0003 & <0.0006 & <0.00089 & <0.0012 & <0.0015 & <0.0018 & <0.0021 & <0.0024 \\
\bottomrule
\end{tabularx}}
\end{adjustwidth}
\end{table}
%\end{center}
%%\end{footnotesize}
%\end{scriptsize}
%---------------------------------------------------------<

%%%%%%%%%%%%%%%%%%%%%%%%%%%%%%%%%%%%%%%%%%
\begin{adjustwidth}{-\extralength}{0cm}
%\printendnotes[custom] % Un-comment to print a list of endnotes

\reftitle{References}

%=====================================
% References, variant A: external bibliography
%=====================================
\begin{thebibliography}{999}

\bibitem[Chesworth(2008)]{ref1}
Chesworth, W. (Ed.) Soil stabilization.
\newblock In {\em Encyclopedia of Soil Science}; Chesworth, W., Ed.; Springer: Dordrecht,The Netherlands,  2008; pp. 705--705.
\newblock {\url{https://doi.org/10.1007/978-1-4020-3995-9_556}}.

\bibitem[Arias-Jaramillo \em{et~al.}(2023)Arias-Jaramillo, Gómez-Cano,
  Carvajal, Hidalgo, and Muñoz]{ma16113996}
Arias-Jaramillo, Y.P.; Gómez-Cano, D.; Carvajal, G.I.; Hidalgo, C.A.; Muñoz,
  F.
\newblock Evaluation of the Effect of Binary Fly Ash-Lime Mixture on the
  Bearing Capacity of Natural Soils: A Comparison with Two Conventional
  Stabilizers Lime and Portland Cement.
\newblock {\em Materials} {\bf 2023}, {\em 16}, 3996.
\newblock {\url{https://doi.org/10.3390/ma16113996}}.

\clearpage 
\bibitem[Yi \em{et~al.}(2014)Yi, Liska, and Al-Tabbaa]{Yietal2014}
Yi, Y.; Liska, M.; Al-Tabbaa, A.
\newblock Properties of Two Model Soils Stabilized with Different Blends and
  Contents of GGBS, MgO, Lime, and PC.
\newblock {\em J. Mater. Civ. Eng.} {\bf 2014}, {\em
  26},~267--274.
\newblock {\url{https://doi.org/10.1061/(ASCE)MT.1943-5533.0000806}}.

\bibitem[Correia \em{et~al.}(2023)Correia, Figueiredo, and
  Rasteiro]{app13084880}
Correia, A.A.S.; Figueiredo, D.; Rasteiro, M.G.
\newblock An Experimental Design Methodology to Evaluate the Key Parameters on
  Dispersion of Carbon Nanotubes Applied in Soil Stabilization.
\newblock {\em Appl. Sci.} {\bf 2023}, {\em 13}, 4880.
\newblock {\url{https://doi.org/10.3390/app13084880}}.

\bibitem[Benschoten \em{et~al.}(1997)Benschoten, Matsumoto, and
  Young]{Benschoten}
Benschoten, J.E.V.; Matsumoto, M.R.; Young, W.H.
\newblock Evaluation and Analysis of Soil Washing for Seven Lead-Contaminated
  Soils.
\newblock {\em J. Environ. Eng.} {\bf 1997}, {\em
  123},~217--224.
\newblock {\url{https://doi.org/10.1061/(ASCE)0733-9372(1997)123:3(217)}}.

\bibitem[Chu(2003)]{Chu}
Chu, W.
\newblock Remediation of Contaminated Soils by Surfactant-Aided Soil Washing.
\newblock In {\em Practice Periodical of Hazardous, Toxic, and
Radioactive Waste Management}; ASCE: Reston, VA, USA, %MDPI: We added publisher and it's location. Please confirm. %<-- Authors' reply: yes, we confirm.
2003; {Volume 7},~pp. 19--24.
\newblock {\url{https://doi.org/10.1061/(ASCE)1090-025X(2003)7:1(19)}}.

\bibitem[Bhandari \em{et~al.}(1994)Bhandari, Dove, and Novak]{Bhandari}
Bhandari, A.; Dove, D.C.; Novak, J.T.
\newblock Soil Washing and Biotreatment of Petroleum\&\#x2010;Contaminated
  Soils.
\newblock {\em J. Environ. Eng.} {\bf 1994}, {\em
  120},~1151--1169.
\newblock {\url{https://doi.org/10.1061/(ASCE)0733-9372(1994)120:5(1151)}}.

\bibitem[Rangeard \em{et~al.}(2003)Rangeard, Y.~Hicher, and Zentar]{Rangeard}
Rangeard, D.; ~Hicher, Y.P.; Zentar, R.
\newblock Determining soil permeability from pressuremeter tests.
\newblock {\em Int. J. Numer. Anal. Methods Geomech.} {\bf 2003}, {\em 27},~1--24.
\newblock {\url{https://doi.org/10.1002/nag.258}}.

\bibitem[Lindh and Lemenkova(2022)]{Lemenkova202222}
Lindh, P.; Lemenkova, P.
\newblock {Permeability, compressive strength and Proctor parameters of silts
  stabilised by Portland cement and ground granulated blast furnace slag
  (GGBFS)}.
\newblock {\em Arch. Mech. Eng.} {\bf 2022}, {\em
  69},~667--692.
\newblock {\url{https://doi.org/10.24425/ame.2022.141522}}.

\bibitem[Schultz(2012)]{Schultz}
Schultz, C.
\newblock Model estimates soil permeability using induced polarization.
\newblock {\em Eos Trans. Am. Geophys. Union} {\bf 2012}, {\em
  93},~150--151.
\newblock {\url{https://doi.org/https://doi.org/10.1029/2012EO140016}}.

\bibitem[Mohammed \em{et~al.}(2002)Mohammed, Kennedy, and Smith]{Mohammedetal}
Mohammed, H.; Kennedy, J.B.; Smith, P.
\newblock Improving the Response of Soil-Metal Structures during Construction.
\newblock {\em J. Bridge Eng.} {\bf 2002}, {\em 7},~6--13.
\newblock {\url{https://doi.org/10.1061/(ASCE)1084-0702(2002)7:1(6)}}.

\bibitem[Zayyat \em{et~al.}(2005)Zayyat, Jackson, Tanaka, Hino, and
  Matsuba]{Zayyat}
Zayyat, M.M.; Jackson, A.W.; Tanaka, T.; Hino, T.; Matsuba, Y., Foundation Soil
  Improvement Using Vibro Compaction Combined with Geo Grids.
\newblock In {\em Innovations in Grouting and Soil Improvement}; ASCE: Reston, VA, USA, %<-- Authors' reply: yes, we confirm.
  2005;
  pp. 1--10.
\newblock {\url{https://doi.org/10.1061/40783(162)25}}.

\bibitem[Sheob \em{et~al.}(2023)Sheob, Sajid, {Asim Ansari}, Rais, Sadique, and
  Ahmad]{SHEOB2023}
Sheob, M.; Sajid, M.; {Asim Ansari}, M.; Rais, I.; Sadique, M.; Ahmad, S.
\newblock Using a blend of cement and waste glass powder to improve the
  properties of clayey soil.
\newblock {\em Mater. Today Proc.} {\bf 2023}.
\newblock {\url{https://doi.org/10.1016/j.matpr.2023.05.440}}.

\bibitem[Tiwari \em{et~al.}(2014)Tiwari, Ajmera, Moubayed, Lemmon, Styler, and
  Martinez]{Tiwari}
Tiwari, B.; Ajmera, B.; Moubayed, S.; Lemmon, A.; Styler, K.; Martinez, J.G.
  Improving Geotechnical Behavior of Clayey Soils with Shredded Rubber
  Tires\&\#x2014;Preliminary Study.
\newblock In {\em Geo-Congress 2014 Technical Papers}; ASCE: Reston, VA, USA, 2014; pp. %<-- Authors' reply: yes, we confirm.
  3734--3743.
\newblock {\url{https://doi.org/10.1061/9780784413272.362}}.

\bibitem[Dharini \em{et~al.}(2023)Dharini, Balamaheswari, and {Nevis
  Presentia}]{DHARINI2023}
Dharini, V.; Balamaheswari, M.; {Nevis Presentia}, A.
\newblock Enhancing the strength of expansive clayey soil using lime as soil
  stabilizing agent along with sodium silicate as grouting chemical.
\newblock {\em Mater. Today Proc.} {\bf 2023}.
\newblock {\url{https://doi.org/10.1016/j.matpr.2023.05.156}}.

\bibitem[{El Hariri} \em{et~al.}(2023){El Hariri}, {Elawad Eltayeb Ahmed}, and
  Kiss]{ELHARIRI202347}
{El Hariri}, A.; {Elawad Eltayeb Ahmed}, A.; Kiss, P.
\newblock Review on soil shear strength with loam sand soil results using
  direct shear test.
\newblock {\em J. Terramechanics} {\bf 2023}, {\em 107},~47--59.
\newblock {\url{https://doi.org/10.1016/j.jterra.2023.03.003}}.

\bibitem[Eissa \em{et~al.}(2022)Eissa, Bassuoni, Ghazy, and Alfaro]{Eissa}
Eissa, A.; Bassuoni, M.T.; Ghazy, A.; Alfaro, M.
\newblock Improving the Properties of Soft Clay Using Cement, Slag, and
  Nanosilica: Experimental and Statistical Modeling.
\newblock {\em J. Mater. Civ. Eng.} {\bf 2022}, {\em
  34},~04022031.
\newblock {\url{https://doi.org/10.1061/(ASCE)MT.1943-5533.0004172}}.

\bibitem[Zhang \em{et~al.}(2023)Zhang, Zhu, Liu, Yang, Wan, Huo, and
  Yang]{ZHANG2023131887}
Zhang, C.; Zhu, Z.; Liu, F.; Yang, Y.; Wan, Y.; Huo, W.; Yang, L.
\newblock Efficient machine learning method for evaluating compressive strength
  of cement stabilized soft soil.
\newblock {\em Constr. Build. Mater.} {\bf 2023}, {\em
  392},~131887.
\newblock {\url{https://doi.org/10.1016/j.conbuildmat.2023.131887}}.

\bibitem[Hogentogler and Willis(1938)]{HogentoglerWillis}
Hogentogler, C.A.; Willis, E.A.
\newblock Essential Consideration in the Stabilization of Soil.
\newblock {\em Trans. Am. Soc. Civ. Eng.} {\bf
  1938}, {\em 103},~1163--1183.
\newblock {\url{https://doi.org/10.1061/TACEAT.0004955}}.

\bibitem[Fristad(1995)]{Fristad}
Fristad, W.E.
\newblock Case study: Using soil washing/leaching for the removal of heavy
  metal at the twin cities army ammunition plant.
\newblock {\em Remediat. J.} {\bf 1995}, {\em 5},~61--72.
\newblock {\url{https://doi.org/10.1002/rem.3440050407}}.

\bibitem[Finney \em{et~al.}(2015)Finney, Himmer, Morris, and
  Coladonato]{Finney}
Finney, D.S.; Himmer, T.; Morris, J.J.; Coladonato, S.
\newblock {Demonstrating Leaching Reductions of NAPL-Impacted Soils Treated
  with Stabilization/Solidification Using Modified EPA Method 1315}.
\newblock {\em J. Hazardous Toxic, Radioact. Waste} {\bf 2015},
  {\em 19},~C4014002.
\newblock {\url{https://doi.org/10.1061/(ASCE)HZ.2153-5515.0000250}}.

\bibitem[Hanson \em{et~al.}(1993)Hanson, Dwyer, Samani, and York]{Hanson}
Hanson, A.T.; Dwyer, B.; Samani, Z.A.; York, D.
\newblock Remediation of Chromium\&\#x2010;Containing Soils by Heap Leaching:
  Column Study.
\newblock {\em J. Environ. Eng.} {\bf 1993}, {\em
  119},~825--841.
\newblock {\url{https://doi.org/10.1061/(ASCE)0733-9372(1993)119:5(825)}}.

\bibitem[He \em{et~al.}(2022)He, Zhou, Cao, and Shen]{He}
He, L.; Zhou, X.; Cao, J.; Shen, L.
\newblock Ultrasound-Assisted Soil Washing for Metals-Contaminated Soil Using
  Various Washing Solutions.
\newblock {\em Clean---Soil Air Water} {\bf 2022}, {\em 50},~2100419.
\newblock {\url{https://doi.org/10.1002/clen.202100419}}.

\bibitem[Mizutani \em{et~al.}(2016)Mizutani, Ikegami, Sakanakura, and
  Kanjo]{Mizutani2016}
Mizutani, S.; Ikegami, M.; Sakanakura, H.; Kanjo, Y. Test Methods for the
  Evaluation of Heavy Metals in Contaminated Soil.
\newblock In {\em Environmental Remediation Technologies for Metal-Contaminated
  Soils}; Springer: Tokyo, Japan, 2016; pp. 67--97.
\newblock {\url{https://doi.org/10.1007/978-4-431-55759-3_4}}.

\bibitem[Hoang \em{et~al.}(2020)Hoang, Alleman, Cetin, and Choi]{Hoang}
Hoang, T.; Alleman, J.; Cetin, B.; Choi, S.G.
\newblock Engineering Properties of Biocementation Coarse- and Fine-Grained
  Sand Catalyzed By Bacterial Cells and Bacterial Enzyme.
\newblock {\em J. Mater. Civ. Eng.} {\bf 2020}, {\em
  32},~04020030.
\newblock {\url{https://doi.org/10.1061/(ASCE)MT.1943-5533.0003083}}.

\bibitem[Schaefers and Nam(2015)]{Schaefers}
Schaefers, K.; Nam, S., Improving the Shear Strength of Soils by Adding
  Agricultural By-Products.
\newblock In {\em IFCEE 2015}; ASCE: Reston, VA, USA, 2015; pp. 2777--2786. %<-- Authors' reply: yes, we confirm.
\newblock {\url{https://doi.org/10.1061/9780784479087.259}}.

\bibitem[Krisdani \em{et~al.}(2005)Krisdani, Rahardjo, and Leong]{Krisdani}
Krisdani, H.; Rahardjo, H.; Leong, E.C., Behaviour of Capillary Barrier System
  Constructed using Residual Soil.
\newblock In {\em Waste Containment and Remediation}; ASCE: Reston, VA, USA, 2005; pp. 1--15. %<-- Authors' reply: yes, we confirm.
\newblock {\url{https://doi.org/10.1061/40789(168)26}}.

\bibitem[Bache \em{et~al.}(2022)Bache, Wiersholm, Paniagua, and Emdal]{Bache}
Bache, B.K.F.; Wiersholm, P.; Paniagua, P.; Emdal, A.
\newblock Effect of Temperature on the Strength of Lime-Cement Stabilized
  Norwegian Clays.
\newblock {\em J. Geotech. Geoenvironmental Eng.} {\bf
  2022}, {\em 148},~04021198.
\newblock {\url{https://doi.org/10.1061/(ASCE)GT.1943-5606.0002699}}.

\bibitem[Andersen \em{et~al.}(1980)Andersen, Rosenbrand, Brown, and
  Pool]{Andersen}
Andersen, K.H.; Rosenbrand, W.F.; Brown, S.F.; Pool, J.H.
\newblock Cyclic and Static Laboratory Tests on Drammen Clay.
\newblock {\em J. Geotech. Eng. Div.} {\bf 1980},
  {\em 106},~499--529.
\newblock {\url{https://doi.org/10.1061/AJGEB6.0000957}}.

\bibitem[Zhang \em{et~al.}(2020)Zhang, Andersen, Jeanjean, Karlsrud, and
  Haugen]{Zhangetal2020}
Zhang, Y.; Andersen, K.H.; Jeanjean, P.; Karlsrud, K.; Haugen, T.
\newblock Validation of Monotonic and Cyclic p-y Framework by Lateral Pile Load
  Tests in Stiff, Overconsolidated Clay at the Haga Site.
\newblock {\em J. Geotech. Geoenvironmental Eng.} {\bf
  2020}, {\em 146},~04020080.
\newblock {\url{https://doi.org/10.1061/(ASCE)GT.1943-5606.0002318}}.

\bibitem[Binner \em{et~al.}(2010)Binner, Homberg, Prohaska, Kalbe, and
  Witt]{Binner}
Binner, R.; Homberg, U.; Prohaska, S.; Kalbe, U.; Witt, K.J. Identification of
  Descriptive Parameters of the Soil Pore Structure Using Experiments and CT
  Data.
\newblock In {\em Scour and Erosion, Proceedings of the Chapter International
  Conference on Scour and Erosion (ICSE-5), San Francisco, CA, USA, 7--10 November 2010 %MDPI: We added the location and date of the conference. Please confirm. %<-- Authors' reply: yes, we confirm.
}; ASCE: Reston, VA, USA, 2010; pp. 397--407. %<-- Authors' reply: yes, we confirm.
\newblock {\url{https://doi.org/10.1061/41147(392)38}}.

\bibitem[Kulkarni \em{et~al.}(2023)Kulkarni, Basavaraj, and Babu]{Kulkarni}
Kulkarni, K.; Basavaraj, P.; Babu, S.
\newblock Effect of Grain Size Distribution of Soil on Immobilization of
  Cadmium and Nickel in Contaminated Soil Using Nano Zerovalent Iron: A
  Factorial Design and Response Surface Methodology Approach.
\newblock {\em J. Hazard. Toxic Radioact. Waste} {\bf 2023},
  {\em 27},~04023008.
\newblock {\url{https://doi.org/10.1061/JHTRBP.HZENG-1195}}.

\bibitem[Ibrahim and Nissen(2012)]{Rahinah}
Ibrahim, R.; Nissen, M., Emerging Technology to Model Dynamic Knowledge
  Creation and Flow among Construction Industry Stakeholders during the
  Critical Feasibility-Entitlements Phase.
\newblock In {\em Towards a Vision for Information Technology in Civil
  Engineering}; ASCE: Reston, VA, USA, 2012; pp. 1--14. %<-- Authors' reply: yes, we confirm.
\newblock {\url{https://doi.org/10.1061/40704(2003)40}}.

\bibitem[Ninić \em{et~al.}(2019)Ninić, Bui, Koch, and Meschke]{Ninic}
Ninić, J.; Bui, H.G.; Koch, C.; Meschke, G.
\newblock Computationally Efficient Simulation in Urban Mechanized Tunneling
  Based on Multilevel BIM Models.
\newblock {\em J. Comput. Civ. Eng.} {\bf 2019}, {\em
  33},~04019007.
\newblock {\url{https://doi.org/10.1061/(ASCE)CP.1943-5487.0000822}}.

\bibitem[Lindh and Lemenkova(2022)]{Lemenkova202226}
Lindh, P.; Lemenkova, P.
\newblock {Simplex Lattice Design and X-ray Diffraction for Analysis of Soil
  Structure: A Case of Cement-Stabilised Compacted Tills Reinforced with Steel
  Slag and Slaked Lime}.
\newblock {\em Electronics} {\bf 2022}, {\em 11},~3726.
\newblock {\url{https://doi.org/10.3390/electronics11223726}}.

\bibitem[Zhang \em{et~al.}(2019)Zhang, Wu, Liu, and Skibniewski]{Zhang}
Zhang, L.; Wu, X.; Liu, W.; Skibniewski, M.J.
\newblock Optimal Strategy to Mitigate Tunnel-Induced Settlement in Soft Soils:
  Simulation Approach.
\newblock {\em J. Perform. Constr. Facil.} {\bf 2019},
  {\em 33},~04019058.
\newblock {\url{https://doi.org/10.1061/(ASCE)CF.1943-5509.0001322}}.

\bibitem[Lindh and Lemenkova(2023)]{Lemenkova202309}
Lindh, P.; Lemenkova, P.
\newblock {Utilising Pareto efficiency and RSM to adjust binder content in clay
  stabilisation for Yttre Ringvägen, Malmö}.
\newblock {\em Acta Polytech.} {\bf 2023}, {\em 63},~140--157.
\newblock {\url{https://doi.org/10.14311/AP.2023.63.0140}}.

\bibitem[Wei(2023)]{Wei}
Wei, W. Optimization of the Mixing in a Produced Water Storage Tank with CFD.
\newblock In {\em World Environmental and Water Resources Congress 2023}; ASCE:
 Reston, VA, USA, 2023; pp. 1175--1184. %<-- Authors' reply: yes, we confirm.
\newblock {\url{https://doi.org/10.1061/9780784484852.107}}.

\bibitem[Mohammed \em{et~al.}(2016)Mohammed, Sanaulla, Alnuaim, and
  Moghal]{Mohammedetal2016}
Mohammed, S.A.S.; Sanaulla, P.F.; Alnuaim, A.M.; Moghal, A.A.B., Role of
  Different Leaching Methods to Arrest the Transport of Ni\textsuperscript{2+} in Soil
  and Soil Amended with Nano Calcium Silicate.
\newblock In Proceedings of the  Geo-China 2016, Shandong, China, 25--27 July 2016; ASCE: Reston, VA, USA, 2016; pp. 49--56. %<-- Authors' reply: yes, we confirm.
\newblock {\url{https://doi.org/10.1061/9780784480045.007}}.

\bibitem[Simon \em{et~al.}(2008)Simon, Kalbe, and Berger]{Simonetal}
Simon, F.G.; Kalbe, U.; Berger, W. Waste Characterization by Leaching and
  Extraction Procedures.
\newblock In Proceedings of the GeoCongress 2008, New Orleans, LA, USA, 9--12 March 2008; ASCE: Reston, VA, USA, 2008; pp. 676--683. %<-- Authors' reply: yes, we confirm.
\newblock {\url{https://doi.org/10.1061/40970(309)85}}.

\bibitem[Naghipour \em{et~al.}(2017)Naghipour, Jaafari, Ashrafi, and
  Mahvi]{Naghipour}
Naghipour, D.; Jaafari, J.; Ashrafi, S.D.; Mahvi, A.H.
\newblock Remediation of Heavy Metals Contaminated Silty Clay Loam Soil by
  Column Extraction with Ethylenediaminetetraacetic Acid and Nitrilo Triacetic
  Acid.
\newblock {\em J. Environ. Eng.} {\bf 2017}, {\em
  143},~04017026.
\newblock {\url{https://doi.org/10.1061/(ASCE)EE.1943-7870.0001219}}.

\bibitem[Wang and Al-Tabbaa(2014)]{WangTabbaa}
Wang, F.; Al-Tabbaa, A., Leachability of 17-Year-Old Stabilized/Solidified
  Contaminated Site Soils.
\newblock In {\em Geo-Congress 2014 Technical Papers}; ASCE: Reston, VA, USA, 2014; pp. %<-- Authors' reply: yes, we confirm.
  1612--1624.
\newblock {\url{https://doi.org/10.1061/9780784413272.158}}.

\bibitem[Li and Li(2001)]{LiLi}
Li, L.Y.; Li, F.
\newblock Heavy Metal Sorption and Hydraulic Conductivity Studies Using Three
  Types of Bentonite Admixes.
\newblock {\em J. Environ. Eng.} {\bf 2001}, {\em
  127},~420--429.
\newblock {\url{https://doi.org/10.1061/(ASCE)0733-9372(2001)127:5(420)}}.

\bibitem[Xia \em{et~al.}(2018)Xia, Feng, Du, Reddy, and Wei]{Xia}
Xia, W.Y.; Feng, Y.S.; Du, Y.J.; Reddy, K.R.; Wei, M.L.
\newblock Solidification and Stabilization of Heavy Metal\&\#x2013;Contaminated
  Industrial Site Soil Using KMP Binder.
\newblock {\em J. Mater. Civ. Eng.} {\bf 2018}, {\em
  30},~04018080.
\newblock {\url{https://doi.org/10.1061/(ASCE)MT.1943-5533.0002264}}.

\bibitem[Liang(2005)]{Liang}
Liang, Y. Biodegradation of Pyrene in Soil Microcosms: Identification of a
  Toxic Intermediate.
\newblock In {\em Impacts of Global Climate Change}; ASCE: Reston, VA, USA, 2005; pp. 1--9. %<-- Authors' reply: yes, we confirm.
\newblock {\url{https://doi.org/10.1061/40792(173)4}}.

\bibitem[Aydilek and Edil(2008)]{Aydilek}
Aydilek, A.H.; Edil, T.B. Solidification/Stabilization of PCB-Contaminated
  Wastewater Treatment Sludges.
\newblock In Proceedings of the  GeoCongress 2008, New Orleans, LA, USA, 9--12 March 2008; ASCE: Reston, VA, USA, 2008; pp. 724--731. %<-- Authors' reply: yes, we confirm.
\newblock {\url{https://doi.org/10.1061/40970(309)91}}.

\bibitem[Kucharski \em{et~al.}(2022)Kucharski, Giebułtowicz, Drobniewska,
  Nałęcz-Jawecki, Skowronek, Strzelecka, Mianowicz, and
  Drzewicz]{KUCHARSKI2022136133}
Kucharski, D.; Giebułtowicz, J.; Drobniewska, A.; Nałęcz-Jawecki, G.;
  Skowronek, A.; Strzelecka, A.; Mianowicz, K.; Drzewicz, P.
\newblock The study on contamination of bottom sediments from the Odra River
  estuary (SW Baltic Sea) by tributyltin using environmetric methods.
\newblock {\em Chemosphere} {\bf 2022}, {\em 308},~136133.
\newblock {\url{https://doi.org/10.1016/j.chemosphere.2022.136133}}.

\bibitem[Lindh and Lemenkova(2022)]{Lemenkova202214}
Lindh, P.; Lemenkova, P.
\newblock {Soil contamination from heavy metals and persistent organic
  pollutants (PAH, PCB and HCB) in the coastal area of Västernorrland,
  Sweden}.
\newblock {\em Gospod. Surowcami Miner.---Miner. Resour. Manag.} {\bf 2022}, {\em 38},~147--168.
\newblock {\url{https://doi.org/10.24425/gsm.2022.141662}}.

\bibitem[Marcic \em{et~al.}(2006)Marcic, {Le Hecho}, Denaix, and
  Lespes]{MARCIC20062322}
Marcic, C.; {Le Hecho}, I.; Denaix, L.; Lespes, G.
\newblock TBT and TPhT persistence in a sludged soil.
\newblock {\em Chemosphere} {\bf 2006}, {\em 65},~2322--2332.
\newblock   {\url{https://doi.org/10.1016/j.chemosphere.2006.05.007}}.

\bibitem[Lindh and Lemenkova(2021)]{Lemenkova202131}
Lindh, P.; Lemenkova, P.
\newblock {Evaluation of Different Binder Combinations of Cement, Slag and CKD
  for S/S Treatment of TBT Contaminated Sediments}.
\newblock {\em Acta Mech. Autom.} {\bf 2021}, {\em 15},~236--248.
\newblock {\url{https://doi.org/10.2478/ama-2021-0030}}.

\bibitem[Chandler \em{et~al.}(1997)Chandler, Eighmy, Hartlén, Hjelmar, Kosson,
  Sawell, {van der Sloot}, and Vehlow]{1997579}
Chandler, A.J.; Eighmy, T.T.; Hartlén, J.; Hjelmar, O.; Kosson, D.S.; Sawell,
  S.E.; {van der Sloot}, H.A.; Vehlow, J.
\newblock Chapter 14---Leaching tests. In {\em Municipal Solid Waste
  Incinerator Residues}; Studies in
  Environmental Science; Elsevier: Amsterdam, The Netherlands, %We added the location of publisher. Please confirm. %<-- Authors' reply: yes, we confirm.
 1997; Volume~67, pp. 579--606.
\newblock {\url{https://doi.org/10.1016/S0166-1116(97)80020-9}}.

\clearpage 
\bibitem[Lindh and Lemenkova(2022)]{Lemenkova202208}
Lindh, P.; Lemenkova, P.
\newblock {Geochemical tests to study the effects of cement ratio on potassium
  and TBT leaching and the pH of the marine sediments from the Kattegat Strait,
  Port of Gothenburg, Sweden}.
\newblock {\em Baltica} {\bf 2022}, {\em 35},~47--59.
\newblock {\url{https://doi.org/10.5200/baltica.2022.1.4}}.

\bibitem[Chittoori \em{et~al.}(2013)Chittoori, Puppala, Wejrungsikul, and
  Hoyos]{Chittoori}
Chittoori, B.C.S.; Puppala, A.J.; Wejrungsikul, T.; Hoyos, L.R.
\newblock Experimental Studies on Stabilized Clays at Various Leaching Cycles.
\newblock {\em J. Geotech. Geoenvironmental Eng.} {\bf
  2013}, {\em 139},~1665--1675.
\newblock {\url{https://doi.org/10.1061/(ASCE)GT.1943-5606.0000920}}.

\bibitem[Beauchesne \em{et~al.}(2008)Beauchesne, Blais, Mercier, and
  Ouarda]{Beauchesne}
Beauchesne, I.; Blais, J.F.; Mercier, G.; Ouarda, T.B.
\newblock Multicriteria Optimization of a Chemical Leaching Process for Sewage
  Sludge Decontamination.
\newblock In {\em Practice Periodical of Hazardous, Toxic, and Radioactive Waste
  Management}; ASCE: Reston, VA, USA, 2008; {Volume 12}, pp.~150--158. %<-- Authors' reply: yes, we confirm.
\newblock {\url{https://doi.org/10.1061/(ASCE)1090-025X(2008)12:3(150)}}.

\bibitem[Lindh and Lemenkova(2022)]{Lemenkova202233}
Lindh, P.; Lemenkova, P.
\newblock {Leaching of Heavy Metals from Contaminated Soil Stabilised by
  Portland Cement and Slag Bremen}.
\newblock {\em Ecol. Chem. Eng. S} {\bf 2022}, {\em
  29},~537--552.
\newblock {\url{https://doi.org/10.2478/eces-2022-0039}}.

\bibitem[Schreck \em{et~al.}(2020)Schreck, Mahedi, and Cetin]{Schreck}
Schreck, S.; Mahedi, M.; Cetin, B. Leaching Behavior of Metals and Sulfate
  from Taconite Tailings Used in Pavement Construction.
\newblock In Proceedings of the  Geo-Congress 2020, Minneapolis, MN, USA, 25--28 February 2020; ASCE: Reston, VA, USA,  2020; pp. 159--168. %<-- Authors' reply: yes, we confirm.
\newblock {\url{https://doi.org/10.1061/9780784482827.018}}.

\bibitem[Becker \em{et~al.}(2013)Becker, Aydilek, Davis, and Seagren]{Becker}
Becker, J.; Aydilek, A.H.; Davis, A.P.; Seagren, E.A.
\newblock Evaluation of Leaching Protocols for Testing of High-Carbon Coal Fly
  Ash-Soil Mixtures.
\newblock {\em J. Environ. Eng.} {\bf 2013}, {\em
  139},~642--653.
\newblock {\url{https://doi.org/10.1061/(ASCE)EE.1943-7870.0000668}}.

\bibitem[Fuessle and Taylor(2004)]{Fuessle}
Fuessle, R.W.; Taylor, M.A.
\newblock Long-Term Solidification/Stabilization and Toxicity Characteristic
  Leaching Procedure for an Electric Arc Furnace Dust.
\newblock {\em J. Environ. Eng.} {\bf 2004}, {\em
  130},~492--498.
\newblock {\url{https://doi.org/10.1061/(ASCE)0733-9372(2004)130:5(492)}}.

\bibitem[Dayioglu \em{et~al.}(2018)Dayioglu, Aydilek, Cimen, and
  Cimen]{Dayioglu}
Dayioglu, A.Y.; Aydilek, A.H.; Cimen, O.; Cimen, M.
\newblock Trace Metal Leaching from Steel Slag Used in Structural Fills.
\newblock {\em J. Geotech. Geoenvironmental Eng.} {\bf
  2018}, {\em 144},~04018089.
\newblock {\url{https://doi.org/10.1061/(ASCE)GT.1943-5606.0001980}}.

\bibitem[Sauer \em{et~al.}(2012)Sauer, Benson, Aydilek, and Edil]{Sauer}
Sauer, J.J.; Benson, C.H.; Aydilek, A.H.; Edil, T.B.
\newblock {Trace Elements Leaching from Organic Soils Stabilized with High
  Carbon Fly Ash}.
\newblock {\em J. Geotech. Geoenvironmental Eng.} {\bf
  2012}, {\em 138},~968--980.
\newblock {\url{https://doi.org/10.1061/(ASCE)GT.1943-5606.0000653}}.

\bibitem[Yu \em{et~al.}(2018)Yu, Dan, and Xin]{Yu}
Yu, X.; Dan, H.C.; Xin, P.
\newblock Method for Improving Leaching Efficiency of Coastal Subsurface
  Drainage Systems.
\newblock {\em J. Irrig. Drain. Eng.} {\bf 2018}, {\em
  144},~04018019.
\newblock {\url{https://doi.org/10.1061/(ASCE)IR.1943-4774.0001330}}.

\bibitem[Alpaslan and Yukselen(2002)]{Alpaslan}
Alpaslan, B.; Yukselen, M.A.
\newblock {Remediation of Lead Contaminated Soils by
  Stabilization/Solidification}.
\newblock {\em Water Air Soil Pollut.} {\bf 2002}, {\em 133},~253--263.
\newblock {\url{https://doi.org/10.1023/A:1012977829536}}.

\bibitem[Liu \em{et~al.}(2023)Liu, Wu, Tan, Yu, and Xu]{ma16093444}
Liu, J.; Wu, D.; Tan, X.; Yu, P.; Xu, L.
\newblock Review of the Interactions between Conventional Cementitious
  Materials and Heavy Metal Ions in Stabilization/Solidification Processing.
\newblock {\em Materials} {\bf 2023}, {\em 16}, 3444.
\newblock {\url{https://doi.org/10.3390/ma16093444}}.

\bibitem[Huang \em{et~al.}(2023)Huang, Shi, Chen, and Yuan]{su15064950}
Huang, H.; Shi, L.; Chen, R.; Yuan, J.
\newblock Effect of Modified Illite on Cd Immobilization and Fertility
  Enhancement of Acidic Soils.
\newblock {\em Sustainability} {\bf 2023}, {\em 15}, 4950.
\newblock {\url{https://doi.org/10.3390/su15064950}}.

\bibitem[Paniagua \em{et~al.}(2023)Paniagua, Ritter, Moseid, and
  Okkenhaug]{Paniagua}
Paniagua, P.; Ritter, S.; Moseid, M.; Okkenhaug, G. Bioashes and Steel Slag as
  Alternative Binders in Ground Improvement of Quick Clays.
\newblock In {\em Geo-Congress 2023 : Soil Improvement, Geoenvironmental, and
  Sustainability}; ASCE: Reston, VA, USA, 2023; pp. 25--34. %<-- Authors' reply: yes, we confirm.
\newblock {\url{https://doi.org/10.1061/9780784484661.003}}.

\bibitem[{Tonini de Araújo} \em{et~al.}(2023){Tonini de Araújo}, {Tonatto
  Ferrazzo}, {Mansur Chaves}, {Gravina da Rocha}, and {Cesar
  Consoli}]{TONINIDEARAUJO2023129757}
{Tonini de Araújo}, M.; {Tonatto Ferrazzo}, S.; {Mansur Chaves}, H.; {Gravina
  da Rocha}, C.; {Cesar Consoli}, N.
\newblock Mechanical behavior, mineralogy, and microstructure of
  alkali-activated wastes-based binder for a clayey soil stabilization.
\newblock {\em Constr. Build. Mater.} {\bf 2023}, {\em
  362},~129757.
\newblock {\url{https://doi.org/10.1016/j.conbuildmat.2022.129757}}.

\bibitem[Filho \em{et~al.}(2021)Filho, Saldanha, da~Rocha, and Consoli]{Filho}
Filho, H.C.S.; Saldanha, R.B.; da~Rocha, C.G.; Consoli, N.C.
\newblock Sustainable Binders Stabilizing Dispersive Clay.
\newblock {\em J. Mater. Civ. Eng.} {\bf 2021}, {\em
  33},~06020026.
\newblock {\url{https://doi.org/10.1061/(ASCE)MT.1943-5533.0003595}}.

\bibitem[Kukko(2000)]{Kukko}
Kukko, H.
\newblock Stabilization of Clay with Inorganic By-Products.
\newblock {\em J. Mater. Civ. Eng.} {\bf 2000}, {\em
  12},~307--309.
\newblock {\url{https://doi.org/10.1061/(ASCE)0899-1561(2000)12:4(307)}}.

\bibitem[Lam \em{et~al.}(2018)Lam, lei Kou, Xie, Chu, and He]{Lametal2018}
Lam, K.P.; lei Kou, H.; Xie, B.; Chu, J.; He, J.
\newblock Use of a Waste-Based Binder for High Water Content Soil Treatment.
\newblock {\em J. Mater. Civ. Eng.} {\bf 2018}, {\em
  30},~06018009.
\newblock {\url{https://doi.org/10.1061/(ASCE)MT.1943-5533.0002385}}.

\bibitem[Ghavami \em{et~al.}(2021)Ghavami, Jahanbakhsh, Saeedi~Azizkandi, and
  Nejad]{Ghavami}
Ghavami, S.; Jahanbakhsh, H.; Saeedi~Azizkandi, A.; Nejad, F.M.
\newblock {Influence of sodium chloride on cement kiln dust-treated clayey
  soil: Strength properties, cost analysis, and environmental impact}.
\newblock {\em Environ. Dev. Sustain.} {\bf 2021}, {\em
  23},~683--702.
\newblock {\url{https://doi.org/10.1007/s10668-020-00603-6}}.

\bibitem[Mishra \em{et~al.}(2019)Mishra, Sachdeva, and Manocha]{Mishra}
Mishra, S.; Sachdeva, S.N.; Manocha, R.
\newblock {Subgrade Soil Stabilization Using Stone Dust and Coarse Aggregate: A
  Cost Effective Approach}.
\newblock {\em Int. J. Geosynth. Ground Eng.}
  {\bf 2019}, {\em 5},~20.
\newblock {\url{https://doi.org/10.1007/s40891-019-0171-0}}.

\bibitem[Lu \em{et~al.}(2021)Lu, Wu, Chen, Li, Shen, Xuan, and Xu]{Lu}
Lu, S.F.; Wu, Y.L.; Chen, Z.; Li, T.; Shen, C.; Xuan, L.K.; Xu, L.
\newblock {Remediation of contaminated soil and groundwater using chemical
  reduction and solidification/stabilization method: A case study}.
\newblock {\em Environ. Sci. Pollut. Res.} {\bf 2021}, {\em
  28},~12766--12779.
\newblock {\url{https://doi.org/10.1007/s11356-020-11337-3}}.

\bibitem[Agarwal \em{et~al.}(2007)Agarwal, Al-Abed, and Dionysiou]{Agarwal}
Agarwal, S.; Al-Abed, S.R.; Dionysiou, D.D.
\newblock In Situ Technologies for Reclamation of PCB-Contaminated Sediments:
  Current Challenges and Research Thrust Areas.
\newblock {\em J. Environ. Eng.} {\bf 2007}, {\em
  133},~1075--1078.
\newblock {\url{https://doi.org/10.1061/(ASCE)0733-9372(2007)133:12(1075)}}.

\bibitem[Du \em{et~al.}(2014)Du, Chadalavada, Chen, and Naidu]{DU2014141}
Du, J.; Chadalavada, S.; Chen, Z.; Naidu, R.
\newblock Environmental remediation techniques of tributyltin contamination in
  soil and water: A review.
\newblock {\em Chem. Eng. J.} {\bf 2014}, {\em 235},~141--150.
\newblock {\url{https://doi.org/10.1016/j.cej.2013.09.044}}.

\bibitem[Sarker \em{et~al.}(2023)Sarker, Masud, Deepo, Das, Nandi, Ansary,
  Islam, and Islam]{SARKER2023138861}
Sarker, A.; Masud, M.A.A.; Deepo, D.M.; Das, K.; Nandi, R.; Ansary, M.W.R.;
  Islam, A.R.M.T.; Islam, T.
\newblock Biological and green remediation of heavy metal contaminated water
  and soils: A state-of-the-art review.
\newblock {\em Chemosphere} {\bf 2023}, {\em 332},~138861.
\newblock {\url{https://doi.org/10.1016/j.chemosphere.2023.138861}}.

\clearpage 
\bibitem[Lu \em{et~al.}(2023)Lu, Xie, Yang, and Chen]{LU2023130021}
Lu, L.; Xie, Y.; Yang, Z.; Chen, B.
\newblock Sustainable decontamination of heavy metal in wastewater and soil
  with novel rectangular wave asymmetrical alternative current
  electrochemistry.
\newblock {\em J. Hazard. Mater.} {\bf 2023}, {\em 442},~130021.
\newblock {\url{https://doi.org/10.1016/j.jhazmat.2022.130021}}.

\bibitem[Chen \em{et~al.}(2010)Chen, Lin, Lu, Ma, Bai, Chen, and
  Li]{CHEN2010638}
Chen, A.; Lin, C.; Lu, W.; Ma, Y.; Bai, Y.; Chen, H.; Li, J.
\newblock Chemical dynamics of acidity and heavy metals in a mine
  water-polluted soil during decontamination using clean water.
\newblock {\em J. Hazard. Mater.} {\bf 2010}, {\em 175},~638--645.
\newblock {\url{https://doi.org/10.1016/j.jhazmat.2009.10.055}}.

\bibitem[Kavamura and Esposito(2010)]{KAVAMURA201061}
Kavamura, V.N.; Esposito, E.
\newblock Biotechnological strategies applied to the decontamination of soils
  polluted with heavy metals.
\newblock {\em Biotechnol. Adv.} {\bf 2010}, {\em 28},~61--69.
\newblock {\url{https://doi.org/10.1016/j.biotechadv.2009.09.002}}.

\bibitem[Chen and Li(2018)]{ChenLi}
Chen, W.; Li, H.
\newblock {Cost-Effectiveness Analysis for Soil Heavy Metal Contamination
  Treatments}.
\newblock {\em Water Air Soil Pollut. Vol.} {\bf 2018}, {\em
  229},~126.
\newblock {\url{https://doi.org/10.1007/s11270-018-3784-3}}.

\bibitem[Karachaliou \em{et~al.}(2016)Karachaliou, Protonotarios, Kaliampakos,
  and Menegaki]{recycling1030328}
Karachaliou, T.; Protonotarios, V.; Kaliampakos, D.; Menegaki, M.
\newblock Using Risk Assessment and Management Approaches to Develop
  Cost-Effective and Sustainable Mine Waste Management Strategies.
\newblock {\em Recycling} {\bf 2016}, {\em 1},~328--342.
\newblock {\url{https://doi.org/10.3390/recycling1030328}}.

\bibitem[Ahmed \em{et~al.}(2022)Ahmed, Nwaichi, Ugwoha, Ugbebor, and
  Arokoyu]{Ahmed}
Ahmed, I.B.; Nwaichi, E.O.; Ugwoha, E.; Ugbebor, J.N.; Arokoyu, S.B.
\newblock {Cost reduction strategies in the remediation of petroleum
  hydrocarbon contaminated soil}.
\newblock {\em Open Res. Afr.} {\bf 2022}, {\em 5},~21.
\newblock {\url{https://doi.org/10.12688/openresafrica.13383.1}}.

\bibitem[Zhang \em{et~al.}(2016)Zhang, Wu, Wang, Mao, and Wu]{7733108}
Zhang, J.; Wu, C.; Wang, L.; Mao, X.; Wu, Y.
\newblock The Work Flow and Operational Model for Geotechnical Investigation
  Based on BIM.
\newblock {\em IEEE Access} {\bf 2016}, {\em 4},~7500--7508.
\newblock {\url{https://doi.org/10.1109/ACCESS.2016.2606158}}.

\bibitem[Lindh(2001)]{Lindh2001}
Lindh, P.
\newblock Optimising binder blends for shallow stabilisation of fine-grained
  soils.
\newblock  {\em Proc. Inst. Civ. Eng.---Ground Improv.} {\bf 2001}, {\em 5},~23--34.
\newblock {\url{https://doi.org/10.1680/grim.2001.5.1.23}}.

\bibitem[Boumezerane(2022)]{geotechnics2010005}
Boumezerane, D.
\newblock Recent Tendencies in the Use of Optimization Techniques in
  Geotechnics: A Review.
\newblock {\em Geotechnics} {\bf 2022}, {\em 2},~114--132.
\newblock {\url{https://doi.org/10.3390/geotechnics2010005}}.

\bibitem[Rehman and Moghal(2018)]{RehmanMoghal}
Rehman, A.U.; Moghal, A.A.B.
\newblock {The Influence and Optimization of Treatment Strategy in Enhancing
  Semiarid Soil Geotechnical Properties}.
\newblock {\em Arab. J. Sci. Eng.} {\bf 2018}, {\em
  43},~5129--5141.
\newblock {\url{https://doi.org/10.1007/s13369-017-2942-z}}.

\bibitem[Yin \em{et~al.}(2018)Yin, Jin, Shen, and Hicher]{Yin}
Yin, Z.Y.; Jin, Y.F.; Shen, J.S.; Hicher, P.Y.
\newblock Optimization techniques for identifying soil parameters in
  geotechnical engineering: Comparative study and enhancement.
\newblock {\em Int. J. Numer. Anal. Methods Geomech.} {\bf 2018}, {\em 42},~70--94.
\newblock {\url{https://doi.org/10.1002/nag.2714}}.

\bibitem[Onyelowe \em{et~al.}(2023)Onyelowe, Mojtahedi, Ebid, Rezaei, Osinubi,
  Eberemu, Salahudeen, Gadzama, Rezazadeh, Jahangir, Yohanna, Onyia, Jalal,
  Iqbal, Ikpa, Obianyo, and Rehman]{Kennedy}
Onyelowe, K.C.; Mojtahedi, F.F.; Ebid, A.M.; Rezaei, A.; Osinubi, K.J.;
  Eberemu, A.O.; Salahudeen, B.; Gadzama, E.W.; Rezazadeh, D.; Jahangir, H.;
  et~al.
\newblock Selected AI optimization techniques and applications in geotechnical
  engineering.
\newblock {\em Cogent Eng.} {\bf 2023}, {\em 10},~2153419.
\newblock {\url{https://doi.org/10.1080/23311916.2022.2153419}}.

\bibitem[{Swedish Institute for Standards}(2011)]{SIS2011}
\emph{SS-EN 197-1:2011}; 
\newblock {Cement---Part 1: Composition, specifications and conformity criteria
  for common cements}.
\newblock {Swedish Institute for Standards}: Stockholm, Sweden, 2011; 48p. %<-- Authors' reply: yes, we confirm.

\bibitem[{Swedish Institute for Standards}(2006)]{SIS2006}
\emph{SS-EN 15167-1:2006};
\newblock {Ground granulated blast furnace slag for use in concrete, mortar and
  grout---Part 1: Definitions, specifications and conformity criteria}.
\newblock {Swedish Institute for Standards}: Stockholm, Sweden, 2006. %<-- Authors' reply: yes, we confirm.

\bibitem[Lindh and Lemenkova(2021)]{Lindh202139}
Lindh, P.; Lemenkova, P.
\newblock {Resonant Frequency Ultrasonic P-Waves for Evaluating Uniaxial
  Compressive Strength of the Stabilized Slag--Cement Sediments}.
\newblock {\em Nord. Concr. Res.} {\bf 2021}, {\em 65},~39--62.
\newblock {\url{https://doi.org/10.2478/ncr-2021-0012}}.

\bibitem[for Standards~(SIS)(2019)]{SIS2019}
\emph{SS-EN ISO 17892-11:2019};
\newblock {Geotechnical investigation and testing---Laboratory testing of soil---Part 11: Permeability tests (ISO 17892-11:2019)}.
\newblock Swedish Institute for Standards: Stockholm, Sweden, 2019. %<-- Authors' reply: yes, we confirm.

\bibitem[{Swedish Institute for Standards}(2015)]{SIS2015}
 \emph{STD-8013628};
\newblock {Karaktärisering av Avfall---Test av GrundläGgande Lakegenskaper---Dynamisk Laktest för Monoliter Med Periodiskt Utbyte av LakväTska vid
  BestäMda TestföRhåLlanden [Characterization of waste---Test of basic leach
  properties---Dynamic Leach Test for Monoliths with Periodic Exchange of Leach
  Liquid at Specified Test Conditions]}.
\newblock Swedish Institute for Standards: Stockholm, Sweden, 2015. (In Swedish) %<-- Authors' reply: yes, we confirm.

\bibitem[Zeffer and Samuelsson(2011)]{ZefferSamuelsson}
Zeffer, A.; Samuelsson, P.O.
\newblock {\em Sedimentprovtagning i SmåBåTshamnar i Stenungsund [Sediment
  Sampling in Marinas in Stenungsund]};
\newblock Technical Report; Stenungsund Municipality: Stenungsund, Sweden,
  2011;
\newblock  62p. (In Swedish)

\bibitem[Inui \em{et~al.}(2008)Inui, Kamon, Katsumi, and Kida]{Inui}
Inui, T.; Kamon, M.; Katsumi, T.; Kida, A. Evaluating Cr(VI) Leaching from
  Recycled Waste Concrete Aggregate Using Acceleration Tests.
\newblock In GeoCongress 2008, New Orleans, LA, USA, 9--12 March 2008; ASCE: Reston, VA, USA, 2008; pp. 280--287. %<-- Authors' reply: yes, we confirm.
\newblock {\url{https://doi.org/10.1061/40970(309)35}}.

\bibitem[Holleran \em{et~al.}(2016)Holleran, Wilson, Holleran, and
  Walubita]{Holleran}
Holleran, I.; Wilson, D.J.; Holleran, G.; Walubita, L.F. Reclaimed Asphalt
  Pavements and Contaminant Leaching-A Literature Review Study.
\newblock In Proceedings of the  Geo-China 2016, Shandong, China, 25--27 July 2016; ASCE: Reston, VA, USA, 2016; pp. 19--28. %<-- Authors' reply: yes, we confirm.
\newblock {\url{https://doi.org/10.1061/9780784480052.003}}.

\bibitem[Chen \em{et~al.}(2021)Chen, Eun, Feng, and Tinjum]{Chenetal2021}
Chen, J.; Eun, J.; Feng, Y.; Tinjum, J.M.
\newblock Long-Term Leaching Behavior of Chromite Ore Processing Residue as
  Backfill Material and the Propagation of Chromium in the Surrounding Soil.
\newblock {\em J. Hazardous Toxic Radioact. Waste} {\bf 2021},
  {\em 25},~04021017.
\newblock {\url{https://doi.org/10.1061/(ASCE)HZ.2153-5515.0000622}}.

\bibitem[Sandhu \em{et~al.}(2013)Sandhu, Axe, Jahan, Ramanujachary, and
  Magdaleno]{Sandhu}
Sandhu, N.K.; Axe, L.B.; Jahan, K.; Ramanujachary, K.V.; Magdaleno, T.F.
\newblock Leaching of Arsenic, Lead, and Antimony from Highway-Marking Glass
  Beads.
\newblock {\em J. Environ. Eng.} {\bf 2013}, {\em
  139},~1168--1177.
\newblock {\url{https://doi.org/10.1061/(ASCE)EE.1943-7870.0000721}}.

\bibitem[Lin \em{et~al.}(2012)Lin, Weng, and Lee]{Linetal2012}
Lin, Y.T.; Weng, C.H.; Lee, S.Y.
\newblock Spatial Distribution of Heavy Metals in Contaminated Agricultural
  Soils Exemplified by Cr, Cu, and Zn.
\newblock {\em J. Environ. Eng.} {\bf 2012}, {\em
  138},~299--306.
\newblock {\url{https://doi.org/10.1061/(ASCE)EE.1943-7870.0000417}}.

\bibitem[Kim \em{et~al.}(2002)Kim, Kim, and Stüben]{Kimetal2002}
Kim, S.O.; Kim, K.W.; Stüben, D.
\newblock Evaluation of Electrokinetic Removal of Heavy Metals from Tailing
  Soils.
\newblock {\em J. Environ. Eng.} {\bf 2002}, {\em
  128},~705--715.
\newblock {\url{https://doi.org/10.1061/(ASCE)0733-9372(2002)128:8(705)}}.

\clearpage 
\bibitem[Bongo \em{et~al.}(2009)Bongo, Mercier, Chartier, Dhenain, and
  Blais]{Bongo}
Bongo, G.; Mercier, G.; Chartier, M.; Dhenain, A.; Blais, J.F.
\newblock Treatment of Aluminum Plant Hazardous Wastes Containing Fluorides and
  PAH.
\newblock {\em J. Environ. Eng.} {\bf 2009}, {\em
  135},~159--166.
\newblock {\url{https://doi.org/10.1061/(ASCE)0733-9372(2009)135:3(159)}}.

 
\bibitem[Nordmark \em{et~al.}(2014)Nordmark, Vestin, Lagerkvist, Lind, Arm, and
  Hallgren]{Nordmark}
Nordmark, D.; Vestin, J.; Lagerkvist, A.; Lind, B.B.; Arm, M.; Hallgren, P.
\newblock Geochemical Behavior of a Gravel Road Upgraded with Wood Fly Ash.
\newblock {\em J. Environ. Eng.} {\bf 2014}, {\em
  140},~05014002.
\newblock {\url{https://doi.org/10.1061/(ASCE)EE.1943-7870.0000821}}.

\end{thebibliography}
\PublishersNote{}

%\nocite{*}

%%%%%%%%%%%%%%%%%%%%%%%%%%%%%%%%%%%%%%%%%%
\end{adjustwidth}
\end{document}
